\chapter{消息渠道}

% === 源文件: channels/bluebubbles.md ===
\section{BlueBubbles（macOS REST）}


状态：内置插件，通过 HTTP 与 BlueBubbles macOS 服务器通信。由于其更丰富的 API 和更简便的设置，\textbf{推荐用于 iMessage 集成}，优于旧版 imsg 渠道。

\subsection{概述}


\begin{itemize}
  \item 通过 BlueBubbles 辅助应用在 macOS 上运行（\href{https://bluebubbles.app}{bluebubbles.app}）。
  \item 推荐/已测试版本：macOS Sequoia (15)。macOS Tahoe (26) 可用；但在 Tahoe 上编辑功能目前不可用，群组图标更新可能显示成功但实际未同步。
  \item OpenClaw 通过其 REST API 与之通信（\texttt{GET /api/v1/ping}、\texttt{POST /message/text}、\texttt{POST /chat/:id/*}）。
  \item 传入消息通过 webhook 到达；发出的回复、输入指示器、已读回执和 tapback 均为 REST 调用。
  \item 附件和贴纸作为入站媒体被接收（并在可能时呈现给智能体）。
  \item 配对/白名单的工作方式与其他渠道相同（\texttt{/channels/pairing} 等），使用 \texttt{channels.bluebubbles.allowFrom} + 配对码。
  \item 回应作为系统事件呈现，与 Slack/Telegram 类似，智能体可以在回复前"提及"它们。
  \item 高级功能：编辑、撤回、回复线程、消息效果、群组管理。
\end{itemize}


\subsection{快速开始}


\begin{enumerate}
  \item 在你的 Mac 上安装 BlueBubbles 服务器（按照 \href{https://bluebubbles.app/install}{bluebubbles.app/install} 的说明操作）。
  \item 在 BlueBubbles 配置中，启用 web API 并设置密码。
  \item 运行 \texttt{openclaw onboard} 并选择 BlueBubbles，或手动配置：
\begin{lstlisting}[language=json]
   {
     channels: {
       bluebubbles: {
         enabled: true,
         serverUrl: "http://192.168.1.100:1234",
         password: "example-password",
         webhookPath: "/bluebubbles-webhook",
       },
     },
   }
\end{lstlisting}

  \item 将 BlueBubbles webhook 指向你的 Gateway 网关（示例：\texttt{https://your-gateway-host:3000/bluebubbles-webhook?password=}）。
  \item 启动 Gateway 网关；它将注册 webhook 处理程序并开始配对。
\end{enumerate}


\subsection{新手引导}


BlueBubbles 可在交互式设置向导中使用：

\begin{lstlisting}
openclaw onboard
\end{lstlisting}


向导会提示输入：

\begin{itemize}
  \item \textbf{服务器 URL}（必填）：BlueBubbles 服务器地址（例如 \texttt{http://192.168.1.100:1234}）
  \item \textbf{密码}（必填）：来自 BlueBubbles 服务器设置的 API 密码
  \item \textbf{Webhook 路径}（可选）：默认为 \texttt{/bluebubbles-webhook}
  \item \textbf{私信策略}：配对、白名单、开放或禁用
  \item \textbf{白名单}：电话号码、电子邮件或聊天目标
\end{itemize}


你也可以通过 CLI 添加 BlueBubbles：

\begin{lstlisting}
openclaw channels add bluebubbles --http-url http://192.168.1.100:1234 --password <password>
\end{lstlisting}


\subsection{访问控制（私信 + 群组）}


私信：

\begin{itemize}
  \item 默认：\texttt{channels.bluebubbles.dmPolicy = "pairing"}。
  \item 未知发送者会收到配对码；在批准之前消息会被忽略（配对码 1 小时后过期）。
  \item 批准方式：
  \item \texttt{openclaw pairing list bluebubbles}
  \item \texttt{openclaw pairing approve bluebubbles }
  \item 配对是默认的令牌交换方式。详情：\href{/channels/pairing}{配对}
\end{itemize}


群组：

\begin{itemize}
  \item \texttt{channels.bluebubbles.groupPolicy = open | allowlist | disabled}（默认：\texttt{allowlist}）。
  \item 当设置为 \texttt{allowlist} 时，\texttt{channels.bluebubbles.groupAllowFrom} 控制谁可以在群组中触发。
\end{itemize}


\subsubsection{提及门控（群组）}


BlueBubbles 支持群聊的提及门控，与 iMessage/WhatsApp 行为一致：

\begin{itemize}
  \item 使用 \texttt{agents.list[].groupChat.mentionPatterns}（或 \texttt{messages.groupChat.mentionPatterns}）检测提及。
  \item 当群组启用 \texttt{requireMention} 时，智能体仅在被提及时响应。
  \item 来自授权发送者的控制命令会绕过提及门控。
\end{itemize}


单群组配置：

\begin{lstlisting}[language=json]
{
  channels: {
    bluebubbles: {
      groupPolicy: "allowlist",
      groupAllowFrom: ["+15555550123"],
      groups: {
        "*": { requireMention: true }, // 所有群组的默认设置
        "iMessage;-;chat123": { requireMention: false }, // 特定群组的覆盖设置
      },
    },
  },
}
\end{lstlisting}


\subsubsection{命令门控}


\begin{itemize}
  \item 控制命令（例如 \texttt{/config}、\texttt{/model}）需要授权。
  \item 使用 \texttt{allowFrom} 和 \texttt{groupAllowFrom} 确定命令授权。
  \item 授权发送者即使在群组中未被提及也可以运行控制命令。
\end{itemize}


\subsection{输入状态 + 已读回执}


\begin{itemize}
  \item \textbf{输入指示器}：在响应生成前和生成期间自动发送。
  \item \textbf{已读回执}：由 \texttt{channels.bluebubbles.sendReadReceipts} 控制（默认：\texttt{true}）。
  \item \textbf{输入指示器}：OpenClaw 发送输入开始事件；BlueBubbles 在发送或超时时自动清除输入状态（通过 DELETE 手动停止不可靠）。
\end{itemize}


\begin{lstlisting}[language=json]
{
  channels: {
    bluebubbles: {
      sendReadReceipts: false, // 禁用已读回执
    },
  },
}
\end{lstlisting}


\subsection{高级操作}


BlueBubbles 在配置中启用时支持高级消息操作：

\begin{lstlisting}[language=json]
{
  channels: {
    bluebubbles: {
      actions: {
        reactions: true, // tapback（默认：true）
        edit: true, // 编辑已发送消息（macOS 13+，在 macOS 26 Tahoe 上不可用）
        unsend: true, // 撤回消息（macOS 13+）
        reply: true, // 通过消息 GUID 进行回复线程
        sendWithEffect: true, // 消息效果（slam、loud 等）
        renameGroup: true, // 重命名群聊
        setGroupIcon: true, // 设置群聊图标/照片（在 macOS 26 Tahoe 上不稳定）
        addParticipant: true, // 将参与者添加到群组
        removeParticipant: true, // 从群组移除参与者
        leaveGroup: true, // 离开群聊
        sendAttachment: true, // 发送附件/媒体
      },
    },
  },
}
\end{lstlisting}


可用操作：

\begin{itemize}
  \item \textbf{react}：添加/移除 tapback 回应（\texttt{messageId}、\texttt{emoji}、\texttt{remove}）
  \item \textbf{edit}：编辑已发送的消息（\texttt{messageId}、\texttt{text}）
  \item \textbf{unsend}：撤回消息（\texttt{messageId}）
  \item \textbf{reply}：回复特定消息（\texttt{messageId}、\texttt{text}、\texttt{to}）
  \item \textbf{sendWithEffect}：带 iMessage 效果发送（\texttt{text}、\texttt{to}、\texttt{effectId}）
  \item \textbf{renameGroup}：重命名群聊（\texttt{chatGuid}、\texttt{displayName}）
  \item \textbf{setGroupIcon}：设置群聊图标/照片（\texttt{chatGuid}、\texttt{media}）— 在 macOS 26 Tahoe 上不稳定（API 可能返回成功但图标未同步）。
  \item \textbf{addParticipant}：将某人添加到群组（\texttt{chatGuid}、\texttt{address}）
  \item \textbf{removeParticipant}：将某人从群组移除（\texttt{chatGuid}、\texttt{address}）
  \item \textbf{leaveGroup}：离开群聊（\texttt{chatGuid}）
  \item \textbf{sendAttachment}：发送媒体/文件（\texttt{to}、\texttt{buffer}、\texttt{filename}、\texttt{asVoice}）
  \item 语音备忘录：将 \texttt{asVoice: true} 与 \textbf{MP3} 或 \textbf{CAF} 音频一起设置，以 iMessage 语音消息形式发送。BlueBubbles 在发送语音备忘录时会将 MP3 转换为 CAF。
\end{itemize}


\subsubsection{消息 ID（短格式 vs 完整格式）}


OpenClaw 可能会显示\textit{短}消息 ID（例如 \texttt{1}、\texttt{2}）以节省 token。

\begin{itemize}
  \item \texttt{MessageSid} / \texttt{ReplyToId} 可以是短 ID。
  \item \texttt{MessageSidFull} / \texttt{ReplyToIdFull} 包含提供商的完整 ID。
  \item 短 ID 存储在内存中；它们可能在重启或缓存清除后过期。
  \item 操作接受短或完整的 \texttt{messageId}，但如果短 ID 不再可用将会报错。
\end{itemize}


对于持久化自动化和存储，请使用完整 ID：

\begin{itemize}
  \item 模板：\texttt{\{\{MessageSidFull\}\}}、\texttt{\{\{ReplyToIdFull\}\}}
  \item 上下文：入站负载中的 \texttt{MessageSidFull} / \texttt{ReplyToIdFull}
\end{itemize}


参见\href{/gateway/configuration}{配置}了解模板变量。

\subsection{分块流式传输}


控制响应是作为单条消息发送还是分块流式传输：

\begin{lstlisting}[language=json]
{
  channels: {
    bluebubbles: {
      blockStreaming: true, // 启用分块流式传输（默认关闭）
    },
  },
}
\end{lstlisting}


\subsection{媒体 + 限制}


\begin{itemize}
  \item 入站附件会被下载并存储在媒体缓存中。
  \item 媒体上限通过 \texttt{channels.bluebubbles.mediaMaxMb} 设置（默认：8 MB）。
  \item 出站文本按 \texttt{channels.bluebubbles.textChunkLimit} 分块（默认：4000 字符）。
\end{itemize}


\subsection{配置参考}


完整配置：\href{/gateway/configuration}{配置}

提供商选项：

\begin{itemize}
  \item \texttt{channels.bluebubbles.enabled}：启用/禁用渠道。
  \item \texttt{channels.bluebubbles.serverUrl}：BlueBubbles REST API 基础 URL。
  \item \texttt{channels.bluebubbles.password}：API 密码。
  \item \texttt{channels.bluebubbles.webhookPath}：Webhook 端点路径（默认：\texttt{/bluebubbles-webhook}）。
  \item \texttt{channels.bluebubbles.dmPolicy}：\texttt{pairing | allowlist | open | disabled}（默认：\texttt{pairing}）。
  \item \texttt{channels.bluebubbles.allowFrom}：私信白名单（句柄、电子邮件、E.164 号码、\texttt{chat\_id:\textit{}、\texttt{chat\_guid:}}）。
  \item \texttt{channels.bluebubbles.groupPolicy}：\texttt{open | allowlist | disabled}（默认：\texttt{allowlist}）。
  \item \texttt{channels.bluebubbles.groupAllowFrom}：群组发送者白名单。
  \item \texttt{channels.bluebubbles.groups}：单群组配置（\texttt{requireMention} 等）。
  \item \texttt{channels.bluebubbles.sendReadReceipts}：发送已读回执（默认：\texttt{true}）。
  \item \texttt{channels.bluebubbles.blockStreaming}：启用分块流式传输（默认：\texttt{false}；流式回复必需）。
  \item \texttt{channels.bluebubbles.textChunkLimit}：出站分块大小（字符）（默认：4000）。
  \item \texttt{channels.bluebubbles.chunkMode}：\texttt{length}（默认）仅在超过 \texttt{textChunkLimit} 时分割；\texttt{newline} 在长度分块前先按空行（段落边界）分割。
  \item \texttt{channels.bluebubbles.mediaMaxMb}：入站媒体上限（MB）（默认：8）。
  \item \texttt{channels.bluebubbles.historyLimit}：上下文的最大群组消息数（0 表示禁用）。
  \item \texttt{channels.bluebubbles.dmHistoryLimit}：私信历史限制。
  \item \texttt{channels.bluebubbles.actions}：启用/禁用特定操作。
  \item \texttt{channels.bluebubbles.accounts}：多账户配置。
\end{itemize}


相关全局选项：

\begin{itemize}
  \item \texttt{agents.list[].groupChat.mentionPatterns}（或 \texttt{messages.groupChat.mentionPatterns}）。
  \item \texttt{messages.responsePrefix}。
\end{itemize}


\subsection{地址 / 投递目标}


优先使用 \texttt{chat\_guid} 以获得稳定的路由：

\begin{itemize}
  \item \texttt{chat\_guid:iMessage;-;+15555550123}（群组推荐）
  \item \texttt{chat\_id:123}
  \item \texttt{chat\_identifier:...}
  \item 直接句柄：\texttt{+15555550123}、\texttt{user@example.com}
  \item 如果直接句柄没有现有的私信聊天，OpenClaw 将通过 \texttt{POST /api/v1/chat/new} 创建一个。这需要启用 BlueBubbles Private API。
\end{itemize}


\subsection{安全性}


\begin{itemize}
  \item Webhook 请求通过比较 \texttt{guid}/\texttt{password} 查询参数或头部与 \texttt{channels.bluebubbles.password} 进行身份验证。来自 \texttt{localhost} 的请求也会被接受。
  \item 保持 API 密码和 webhook 端点的机密性（将它们视为凭证）。
  \item localhost 信任意味着同主机的反向代理可能无意中绕过密码验证。如果你使用代理 Gateway 网关，请在代理处要求身份验证并配置 \texttt{gateway.trustedProxies}。参见 \href{/gateway/security\#reverse-proxy-configuration}{Gateway 网关安全性}。
  \item 如果将 BlueBubbles 服务器暴露在局域网之外，请启用 HTTPS + 防火墙规则。
\end{itemize}


\subsection{故障排除}


\begin{itemize}
  \item 如果输入/已读事件停止工作，请检查 BlueBubbles webhook 日志并验证 Gateway 网关路径是否与 \texttt{channels.bluebubbles.webhookPath} 匹配。
  \item 配对码在一小时后过期；使用 \texttt{openclaw pairing list bluebubbles} 和 \texttt{openclaw pairing approve bluebubbles }。
  \item 回应需要 BlueBubbles private API（\texttt{POST /api/v1/message/react}）；确保服务器版本支持它。
  \item 编辑/撤回需要 macOS 13+ 和兼容的 BlueBubbles 服务器版本。在 macOS 26（Tahoe）上，由于 private API 变更，编辑功能目前不可用。
  \item 在 macOS 26（Tahoe）上群组图标更新可能不稳定：API 可能返回成功但新图标未同步。
  \item OpenClaw 会根据 BlueBubbles 服务器的 macOS 版本自动隐藏已知不可用的操作。如果在 macOS 26（Tahoe）上编辑仍然显示，请使用 \texttt{channels.bluebubbles.actions.edit=false} 手动禁用。
  \item 查看状态/健康信息：\texttt{openclaw status --all} 或 \texttt{openclaw status --deep}。
\end{itemize}


有关通用渠道工作流参考，请参阅\href{/channels}{渠道}和\href{/tools/plugin}{插件}指南。


\clearpage

% === 源文件: channels/broadcast-groups.md ===
\section{广播群组}


\textbf{状态：} 实验性功能
\textbf{版本：} 于 2026.1.9 版本新增

\subsection{概述}


广播群组允许多个智能体同时处理并响应同一条消息。这使你能够在单个 WhatsApp 群组或私信中创建协同工作的专业智能体团队——全部使用同一个手机号码。

当前范围：\textbf{仅限 WhatsApp}（web 渠道）。

广播群组在渠道白名单和群组激活规则之后进行评估。在 WhatsApp 群组中，这意味着广播会在 OpenClaw 正常回复时发生（例如：被提及时，具体取决于你的群组设置）。

\subsection{使用场景}


\subsubsection{1. 专业智能体团队}


部署多个具有原子化、专注职责的智能体：

\begin{lstlisting}
Group: "Development Team"
Agents:
  - CodeReviewer (reviews code snippets)
  - DocumentationBot (generates docs)
  - SecurityAuditor (checks for vulnerabilities)
  - TestGenerator (suggests test cases)
\end{lstlisting}


每个智能体处理相同的消息并提供其专业视角。

\subsubsection{2. 多语言支持}


\begin{lstlisting}
Group: "International Support"
Agents:
  - Agent_EN (responds in English)
  - Agent_DE (responds in German)
  - Agent_ES (responds in Spanish)
\end{lstlisting}


\subsubsection{3. 质量保证工作流}


\begin{lstlisting}
Group: "Customer Support"
Agents:
  - SupportAgent (provides answer)
  - QAAgent (reviews quality, only responds if issues found)
\end{lstlisting}


\subsubsection{4. 任务自动化}


\begin{lstlisting}
Group: "Project Management"
Agents:
  - TaskTracker (updates task database)
  - TimeLogger (logs time spent)
  - ReportGenerator (creates summaries)
\end{lstlisting}


\subsection{配置}


\subsubsection{基本设置}


添加一个顶层 \texttt{broadcast} 部分（与 \texttt{bindings} 同级）。键为 WhatsApp peer id：

\begin{itemize}
  \item 群聊：群组 JID（例如 \texttt{120363403215116621@g.us}）
  \item 私信：E.164 格式的电话号码（例如 \texttt{+15551234567}）
\end{itemize}


\begin{lstlisting}[language=json]
{
  "broadcast": {
    "120363403215116621@g.us": ["alfred", "baerbel", "assistant3"]
  }
}
\end{lstlisting}


\textbf{结果：} 当 OpenClaw 在此聊天中回复时，将运行所有三个智能体。

\subsubsection{处理策略}


控制智能体如何处理消息：

\paragraph{并行（默认）}


所有智能体同时处理：

\begin{lstlisting}[language=json]
{
  "broadcast": {
    "strategy": "parallel",
    "120363403215116621@g.us": ["alfred", "baerbel"]
  }
}
\end{lstlisting}


\paragraph{顺序}


智能体按顺序处理（后一个等待前一个完成）：

\begin{lstlisting}[language=json]
{
  "broadcast": {
    "strategy": "sequential",
    "120363403215116621@g.us": ["alfred", "baerbel"]
  }
}
\end{lstlisting}


\subsubsection{完整示例}


\begin{lstlisting}[language=json]
{
  "agents": {
    "list": [
      {
        "id": "code-reviewer",
        "name": "Code Reviewer",
        "workspace": "/path/to/code-reviewer",
        "sandbox": { "mode": "all" }
      },
      {
        "id": "security-auditor",
        "name": "Security Auditor",
        "workspace": "/path/to/security-auditor",
        "sandbox": { "mode": "all" }
      },
      {
        "id": "docs-generator",
        "name": "Documentation Generator",
        "workspace": "/path/to/docs-generator",
        "sandbox": { "mode": "all" }
      }
    ]
  },
  "broadcast": {
    "strategy": "parallel",
    "120363403215116621@g.us": ["code-reviewer", "security-auditor", "docs-generator"],
    "120363424282127706@g.us": ["support-en", "support-de"],
    "+15555550123": ["assistant", "logger"]
  }
}
\end{lstlisting}


\subsection{工作原理}


\subsubsection{消息流程}


\begin{enumerate}
  \item \textbf{接收消息} 到达 WhatsApp 群组
  \item \textbf{广播检查}：系统检查 peer ID 是否在 \texttt{broadcast} 中
  \item \textbf{如果在广播列表中}：
\end{enumerate}

\begin{itemize}
  \item 所有列出的智能体处理该消息
  \item 每个智能体有自己的会话键和隔离的上下文
  \item 智能体并行处理（默认）或顺序处理
\end{itemize}

\begin{enumerate}
  \item \textbf{如果不在广播列表中}：
\end{enumerate}

\begin{itemize}
  \item 应用正常路由（第一个匹配的绑定）
\end{itemize}


注意：广播群组不会绕过渠道白名单或群组激活规则（提及/命令等）。它们只改变消息符合处理条件时\textit{运行哪些智能体}。

\subsubsection{会话隔离}


广播群组中的每个智能体完全独立维护：

\begin{itemize}
  \item \textbf{会话键}（\texttt{agent:alfred:whatsapp:group:120363...} vs \texttt{agent:baerbel:whatsapp:group:120363...}）
  \item \textbf{对话历史}（智能体看不到其他智能体的消息）
  \item \textbf{工作空间}（如果配置了则使用独立的沙箱）
  \item \textbf{工具访问权限}（不同的允许/拒绝列表）
  \item \textbf{记忆/上下文}（独立的 IDENTITY.md、SOUL.md 等）
  \item \textbf{群组上下文缓冲区}（用于上下文的最近群组消息）按 peer 共享，因此所有广播智能体在被触发时看到相同的上下文
\end{itemize}


这允许每个智能体拥有：

\begin{itemize}
  \item 不同的个性
  \item 不同的工具访问权限（例如只读 vs 读写）
  \item 不同的模型（例如 opus vs sonnet）
  \item 不同的已安装 Skills
\end{itemize}


\subsubsection{示例：隔离的会话}


在群组 \texttt{120363403215116621@g.us} 中，智能体为 \texttt{["alfred", "baerbel"]}：

\textbf{Alfred 的上下文：}

\begin{lstlisting}
Session: agent:alfred:whatsapp:group:120363403215116621@g.us
History: [user message, alfred's previous responses]
Workspace: /Users/pascal/openclaw-alfred/
Tools: read, write, exec
\end{lstlisting}


\textbf{Bärbel 的上下文：}

\begin{lstlisting}
Session: agent:baerbel:whatsapp:group:120363403215116621@g.us
History: [user message, baerbel's previous responses]
Workspace: /Users/pascal/openclaw-baerbel/
Tools: read only
\end{lstlisting}


\subsection{最佳实践}


\subsubsection{1. 保持智能体专注}


将每个智能体设计为具有单一、明确的职责：

\begin{lstlisting}[language=json]
{
  "broadcast": {
    "DEV_GROUP": ["formatter", "linter", "tester"]
  }
}
\end{lstlisting}


\emoji{✅} \textbf{好的做法：} 每个智能体只有一个任务
\emoji{❌} \textbf{不好的做法：} 一个通用的"dev-helper"智能体

\subsubsection{2. 使用描述性名称}


明确每个智能体的功能：

\begin{lstlisting}[language=json]
{
  "agents": {
    "security-scanner": { "name": "Security Scanner" },
    "code-formatter": { "name": "Code Formatter" },
    "test-generator": { "name": "Test Generator" }
  }
}
\end{lstlisting}


\subsubsection{3. 配置不同的工具访问权限}


只给智能体提供它们需要的工具：

\begin{lstlisting}[language=json]
{
  "agents": {
    "reviewer": {
      "tools": { "allow": ["read", "exec"] } // Read-only
    },
    "fixer": {
      "tools": { "allow": ["read", "write", "edit", "exec"] } // Read-write
    }
  }
}
\end{lstlisting}


\subsubsection{4. 监控性能}


当有多个智能体时，请考虑：

\begin{itemize}
  \item 使用 \texttt{"strategy": "parallel"}（默认）以提高速度
  \item 将广播群组限制在 5-10 个智能体
  \item 为较简单的智能体使用较快的模型
\end{itemize}


\subsubsection{5. 优雅地处理失败}


智能体独立失败。一个智能体的错误不会阻塞其他智能体：

\begin{lstlisting}
Message → [Agent A ✓, Agent B ✗ error, Agent C ✓]
Result: Agent A and C respond, Agent B logs error
\end{lstlisting}


\subsection{兼容性}


\subsubsection{提供商}


广播群组目前支持：

\begin{itemize}
  \item \emoji{✅} WhatsApp（已实现）
  \item \emoji{🚧} Telegram（计划中）
  \item \emoji{🚧} Discord（计划中）
  \item \emoji{🚧} Slack（计划中）
\end{itemize}


\subsubsection{路由}


广播群组与现有路由一起工作：

\begin{lstlisting}[language=json]
{
  "bindings": [
    {
      "match": { "channel": "whatsapp", "peer": { "kind": "group", "id": "GROUP_A" } },
      "agentId": "alfred"
    }
  ],
  "broadcast": {
    "GROUP_B": ["agent1", "agent2"]
  }
}
\end{lstlisting}


\begin{itemize}
  \item \texttt{GROUP\_A}：只有 alfred 响应（正常路由）
  \item \texttt{GROUP\_B}：agent1 和 agent2 都响应（广播）
\end{itemize}


\textbf{优先级：} \texttt{broadcast} 优先于 \texttt{bindings}。

\subsection{故障排除}


\subsubsection{智能体不响应}


\textbf{检查：}

\begin{enumerate}
  \item 智能体 ID 存在于 \texttt{agents.list} 中
  \item Peer ID 格式正确（例如 \texttt{120363403215116621@g.us}）
  \item 智能体不在拒绝列表中
\end{enumerate}


\textbf{调试：}

\begin{lstlisting}[language=bash]
tail -f ~/.openclaw/logs/gateway.log | grep broadcast
\end{lstlisting}


\subsubsection{只有一个智能体响应}


\textbf{原因：} Peer ID 可能在 \texttt{bindings} 中但不在 \texttt{broadcast} 中。

\textbf{修复：} 添加到广播配置或从绑定中移除。

\subsubsection{性能问题}


\textbf{如果智能体较多时速度较慢：}

\begin{itemize}
  \item 减少每个群组的智能体数量
  \item 使用较轻的模型（sonnet 而非 opus）
  \item 检查沙箱启动时间
\end{itemize}


\subsection{示例}


\subsubsection{示例 1：代码审查团队}


\begin{lstlisting}[language=json]
{
  "broadcast": {
    "strategy": "parallel",
    "120363403215116621@g.us": [
      "code-formatter",
      "security-scanner",
      "test-coverage",
      "docs-checker"
    ]
  },
  "agents": {
    "list": [
      {
        "id": "code-formatter",
        "workspace": "~/agents/formatter",
        "tools": { "allow": ["read", "write"] }
      },
      {
        "id": "security-scanner",
        "workspace": "~/agents/security",
        "tools": { "allow": ["read", "exec"] }
      },
      {
        "id": "test-coverage",
        "workspace": "~/agents/testing",
        "tools": { "allow": ["read", "exec"] }
      },
      { "id": "docs-checker", "workspace": "~/agents/docs", "tools": { "allow": ["read"] } }
    ]
  }
}
\end{lstlisting}


\textbf{用户发送：} 代码片段
\textbf{响应：}

\begin{itemize}
  \item code-formatter："修复了缩进并添加了类型提示"
  \item security-scanner："\emoji{⚠} 第 12 行存在 SQL 注入漏洞"
  \item test-coverage："覆盖率为 45\%，缺少错误情况的测试"
  \item docs-checker："函数 \texttt{process\_data} 缺少文档字符串"
\end{itemize}


\subsubsection{示例 2：多语言支持}


\begin{lstlisting}[language=json]
{
  "broadcast": {
    "strategy": "sequential",
    "+15555550123": ["detect-language", "translator-en", "translator-de"]
  },
  "agents": {
    "list": [
      { "id": "detect-language", "workspace": "~/agents/lang-detect" },
      { "id": "translator-en", "workspace": "~/agents/translate-en" },
      { "id": "translator-de", "workspace": "~/agents/translate-de" }
    ]
  }
}
\end{lstlisting}


\subsection{API 参考}


\subsubsection{配置模式}


\begin{lstlisting}[language=typescript]
interface OpenClawConfig {
  broadcast?: {
    strategy?: "parallel" | "sequential";
    [peerId: string]: string[];
  };
}
\end{lstlisting}


\subsubsection{字段}


\begin{itemize}
  \item \texttt{strategy}（可选）：如何处理智能体
  \item \texttt{"parallel"}（默认）：所有智能体同时处理
  \item \texttt{"sequential"}：智能体按数组顺序处理
  \item \texttt{[peerId]}：WhatsApp 群组 JID、E.164 号码或其他 peer ID
  \item 值：应处理消息的智能体 ID 数组
\end{itemize}


\subsection{限制}


\begin{enumerate}
  \item \textbf{最大智能体数：} 无硬性限制，但 10 个以上智能体可能会较慢
  \item \textbf{共享上下文：} 智能体看不到彼此的响应（设计如此）
  \item \textbf{消息顺序：} 并行响应可能以任意顺序到达
  \item \textbf{速率限制：} 所有智能体都计入 WhatsApp 速率限制
\end{enumerate}


\subsection{未来增强}


计划中的功能：

\begin{itemize}
  \item [ ] 共享上下文模式（智能体可以看到彼此的响应）
  \item [ ] 智能体协调（智能体可以相互发信号）
  \item [ ] 动态智能体选择（根据消息内容选择智能体）
  \item [ ] 智能体优先级（某些智能体先于其他智能体响应）
\end{itemize}


\subsection{另请参阅}


\begin{itemize}
  \item \href{/tools/multi-agent-sandbox-tools}{多智能体配置}
  \item \href{/channels/channel-routing}{路由配置}
  \item \href{/concepts/session}{会话管理}
\end{itemize}



\clearpage

% === 源文件: channels/channel-routing.md ===
\section{渠道与路由}


OpenClaw 将回复\textbf{路由回消息来源的渠道}。模型不会选择渠道；路由是确定性的，由主机配置控制。

\subsection{关键术语}


\begin{itemize}
  \item \textbf{渠道}：\texttt{whatsapp}、\texttt{telegram}、\texttt{discord}、\texttt{slack}、\texttt{signal}、\texttt{imessage}、\texttt{webchat}。
  \item \textbf{AccountId}：每个渠道的账户实例（在支持的情况下）。
  \item \textbf{AgentId}：隔离的工作区 + 会话存储（"大脑"）。
  \item \textbf{SessionKey}：用于存储上下文和控制并发的桶键。
\end{itemize}


\subsection{会话键格式（示例）}


私信会折叠到智能体的\textbf{主}会话：

\begin{itemize}
  \item \texttt{agent::}（默认：\texttt{agent:main:main}）
\end{itemize}


群组和渠道按渠道隔离：

\begin{itemize}
  \item 群组：\texttt{agent:::group:}
  \item 渠道/房间：\texttt{agent:::channel:}
\end{itemize}


线程：

\begin{itemize}
  \item Slack/Discord 线程会在基础键后追加 \texttt{:thread:}。
  \item Telegram 论坛主题在群组键中嵌入 \texttt{:topic:}。
\end{itemize}


示例：

\begin{itemize}
  \item \texttt{agent:main:telegram:group:-1001234567890:topic:42}
  \item \texttt{agent:main:discord:channel:123456:thread:987654}
\end{itemize}


\subsection{路由规则（如何选择智能体）}


路由为每条入站消息选择\textbf{一个智能体}：

\begin{enumerate}
  \item \textbf{精确对端匹配}（\texttt{bindings} 中的 \texttt{peer.kind} + \texttt{peer.id}）。
  \item \textbf{Guild 匹配}（Discord）通过 \texttt{guildId}。
  \item \textbf{Team 匹配}（Slack）通过 \texttt{teamId}。
  \item \textbf{账户匹配}（渠道上的 \texttt{accountId}）。
  \item \textbf{渠道匹配}（该渠道上的任意账户）。
  \item \textbf{默认智能体}（\texttt{agents.list[].default}，否则取列表第一项，兜底为 \texttt{main}）。
\end{enumerate}


匹配到的智能体决定使用哪个工作区和会话存储。

\subsection{广播组（运行多个智能体）}


广播组允许你为同一对端运行\textbf{多个智能体}，\textbf{在 OpenClaw 正常回复时}触发（例如：在 WhatsApp 群组中，经过提及/激活门控之后）。

配置：

\begin{lstlisting}[language=json]
{
  broadcast: {
    strategy: "parallel",
    "120363403215116621@g.us": ["alfred", "baerbel"],
    "+15555550123": ["support", "logger"],
  },
}
\end{lstlisting}


参见：\href{/channels/broadcast-groups}{广播组}。

\subsection{配置概览}


\begin{itemize}
  \item \texttt{agents.list}：命名的智能体定义（工作区、模型等）。
  \item \texttt{bindings}：将入站渠道/账户/对端映射到智能体。
\end{itemize}


示例：

\begin{lstlisting}[language=json]
{
  agents: {
    list: [{ id: "support", name: "Support", workspace: "~/.openclaw/workspace-support" }],
  },
  bindings: [
    { match: { channel: "slack", teamId: "T123" }, agentId: "support" },
    { match: { channel: "telegram", peer: { kind: "group", id: "-100123" } }, agentId: "support" },
  ],
}
\end{lstlisting}


\subsection{会话存储}


会话存储位于状态目录下（默认 \texttt{\textasciitilde{}/.openclaw}）：

\begin{itemize}
  \item \texttt{\textasciitilde{}/.openclaw/agents//sessions/sessions.json}
  \item JSONL 记录文件与存储位于同一目录
\end{itemize}


你可以通过 \texttt{session.store} 和 \texttt{\{agentId\}} 模板来覆盖存储路径。

\subsection{WebChat 行为}


WebChat 连接到\textbf{所选智能体}，并默认使用该智能体的主会话。因此，WebChat 让你可以在一个地方查看该智能体的跨渠道上下文。

\subsection{回复上下文}


入站回复包含：

\begin{itemize}
  \item \texttt{ReplyToId}、\texttt{ReplyToBody} 和 \texttt{ReplyToSender}（在可用时）。
  \item 引用的上下文会以 \texttt{[Replying to ...]} 块的形式追加到 \texttt{Body} 中。
\end{itemize}


这在所有渠道中保持一致。


\clearpage

% === 源文件: channels/discord.md ===
\section{Discord（Bot API）}


状态：已支持通过官方 Discord 机器人网关进行私信和服务器文字频道通信。

\subsection{快速设置（新手）}


\begin{enumerate}
  \item 创建 Discord 机器人并复制机器人令牌。
  \item 在 Discord 应用设置中启用 \textbf{Message Content Intent}（如果你计划使用允许列表或名称查找，还需启用 \textbf{Server Members Intent}）。
  \item 为 OpenClaw 设置令牌：
\end{enumerate}

\begin{itemize}
  \item 环境变量：\texttt{DISCORD\_BOT\_TOKEN=...}
  \item 或配置：\texttt{channels.discord.token: "..."}。
  \item 如果两者都设置，配置优先（环境变量回退仅适用于默认账户）。
\end{itemize}

\begin{enumerate}
  \item 使用消息权限邀请机器人到你的服务器（如果你只想使用私信，可以创建一个私人服务器）。
  \item 启动 Gateway 网关。
  \item 私信访问默认采用配对模式；首次联系时需批准配对码。
\end{enumerate}


最小配置：

\begin{lstlisting}[language=json]
{
  channels: {
    discord: {
      enabled: true,
      token: "YOUR_BOT_TOKEN",
    },
  },
}
\end{lstlisting}


\subsection{目标}


\begin{itemize}
  \item 通过 Discord 私信或服务器频道与 OpenClaw 对话。
  \item 直接聊天会合并到智能体的主会话（默认 \texttt{agent:main:main}）；服务器频道保持隔离为 \texttt{agent::discord:channel:}（显示名称使用 \texttt{discord:\#}）。
  \item 群组私信默认被忽略；通过 \texttt{channels.discord.dm.groupEnabled} 启用，并可选择通过 \texttt{channels.discord.dm.groupChannels} 进行限制。
  \item 保持路由确定性：回复始终返回到消息来源的渠道。
\end{itemize}


\subsection{工作原理}


\begin{enumerate}
  \item 创建 Discord 应用程序 → Bot，启用你需要的意图（私信 + 服务器消息 + 消息内容），并获取机器人令牌。
  \item 使用所需权限邀请机器人到你的服务器，以便在你想使用的地方读取/发送消息。
  \item 使用 \texttt{channels.discord.token} 配置 OpenClaw（或使用 \texttt{DISCORD\_BOT\_TOKEN} 作为回退）。
  \item 运行 Gateway 网关；当令牌可用（配置优先，环境变量回退）且 \texttt{channels.discord.enabled} 不为 \texttt{false} 时，它会自动启动 Discord 渠道。
\end{enumerate}

\begin{itemize}
  \item 如果你更喜欢使用环境变量，设置 \texttt{DISCORD\_BOT\_TOKEN}（配置块是可选的）。
\end{itemize}

\begin{enumerate}
  \item 直接聊天：发送时使用 \texttt{user:}（或 \texttt{<@id>} 提及）；所有对话都进入共享的 \texttt{main} 会话。纯数字 ID 是模糊的，会被拒绝。
  \item 服务器频道：发送时使用 \texttt{channel:}。默认需要提及，可以按服务器或按频道设置。
  \item 直接聊天：默认通过 \texttt{channels.discord.dm.policy} 进行安全保护（默认：\texttt{"pairing"}）。未知发送者会收到配对码（1 小时后过期）；通过 \texttt{openclaw pairing approve discord } 批准。
\end{enumerate}

\begin{itemize}
  \item 要保持旧的"对任何人开放"行为：设置 \texttt{channels.discord.dm.policy="open"} 和 \texttt{channels.discord.dm.allowFrom=["*"]}。
  \item 要使用硬编码允许列表：设置 \texttt{channels.discord.dm.policy="allowlist"} 并在 \texttt{channels.discord.dm.allowFrom} 中列出发送者。
  \item 要忽略所有私信：设置 \texttt{channels.discord.dm.enabled=false} 或 \texttt{channels.discord.dm.policy="disabled"}。
\end{itemize}

\begin{enumerate}
  \item 群组私信默认被忽略；通过 \texttt{channels.discord.dm.groupEnabled} 启用，并可选择通过 \texttt{channels.discord.dm.groupChannels} 进行限制。
  \item 可选服务器规则：设置 \texttt{channels.discord.guilds}，以服务器 ID（首选）或 slug 为键，并包含每个频道的规则。
  \item 可选原生命令：\texttt{commands.native} 默认为 \texttt{"auto"}（Discord/Telegram 开启，Slack 关闭）。使用 \texttt{channels.discord.commands.native: true|false|"auto"} 覆盖；\texttt{false} 会清除之前注册的命令。文本命令由 \texttt{commands.text} 控制，必须作为独立的 \texttt{/...} 消息发送。使用 \texttt{commands.useAccessGroups: false} 可跳过命令的访问组检查。
\end{enumerate}

\begin{itemize}
  \item 完整命令列表 + 配置：\href{/tools/slash-commands}{斜杠命令}
\end{itemize}

\begin{enumerate}
  \item 可选服务器上下文历史：设置 \texttt{channels.discord.historyLimit}（默认 20，回退到 \texttt{messages.groupChat.historyLimit}）以在回复提及时包含最近 N 条服务器消息作为上下文。设置 \texttt{0} 禁用。
  \item 表情反应：智能体可以通过 \texttt{discord} 工具触发表情反应（受 \texttt{channels.discord.actions.*} 控制）。
\end{enumerate}

\begin{itemize}
  \item 表情反应移除语义：参见 \href{/tools/reactions}{/tools/reactions}。
  \item \texttt{discord} 工具仅在当前渠道是 Discord 时暴露。
\end{itemize}

\begin{enumerate}
  \item 原生命令使用隔离的会话键（\texttt{agent::discord:slash:}）而不是共享的 \texttt{main} 会话。
\end{enumerate}


注意：名称 → ID 解析使用服务器成员搜索，需要 Server Members Intent；如果机器人无法搜索成员，请使用 ID 或 \texttt{<@id>} 提及。
注意：Slug 为小写，空格替换为 \texttt{-}。频道名称的 slug 不包含前导 \texttt{\#}。
注意：服务器上下文 \texttt{[from:]} 行包含 \texttt{author.tag} + \texttt{id}，便于进行可提及的回复。

\subsection{配置写入}


默认情况下，允许 Discord 写入由 \texttt{/config set|unset} 触发的配置更新（需要 \texttt{commands.config: true}）。

禁用方式：

\begin{lstlisting}[language=json]
{
  channels: { discord: { configWrites: false } },
}
\end{lstlisting}


\subsection{如何创建自己的机器人}


这是在服务器（guild）频道（如 \texttt{\#help}）中运行 OpenClaw 的"Discord 开发者门户"设置。

\subsubsection{1）创建 Discord 应用 + 机器人用户}


\begin{enumerate}
  \item Discord 开发者门户 → \textbf{Applications} → \textbf{New Application}
  \item 在你的应用中：
\end{enumerate}

\begin{itemize}
  \item \textbf{Bot} → \textbf{Add Bot}
  \item 复制 \textbf{Bot Token}（这是你放入 \texttt{DISCORD\_BOT\_TOKEN} 的内容）
\end{itemize}


\subsubsection{2）启用 OpenClaw 需要的网关意图}


Discord 会阻止"特权意图"，除非你明确启用它们。

在 \textbf{Bot} → \textbf{Privileged Gateway Intents} 中启用：

\begin{itemize}
  \item \textbf{Message Content Intent}（在大多数服务器中读取消息文本所必需；没有它你会看到"Used disallowed intents"或机器人会连接但不响应消息）
  \item \textbf{Server Members Intent}（推荐；服务器中的某些成员/用户查找和允许列表匹配需要）
\end{itemize}


你通常\textbf{不需要} \textbf{Presence Intent}。

\subsubsection{3）生成邀请 URL（OAuth2 URL Generator）}


在你的应用中：\textbf{OAuth2} → \textbf{URL Generator}

\textbf{Scopes}

\begin{itemize}
  \item \emoji{✅} \texttt{bot}
  \item \emoji{✅} \texttt{applications.commands}（原生命令所需）
\end{itemize}


\textbf{Bot Permissions}（最小基线）

\begin{itemize}
  \item \emoji{✅} View Channels
  \item \emoji{✅} Send Messages
  \item \emoji{✅} Read Message History
  \item \emoji{✅} Embed Links
  \item \emoji{✅} Attach Files
  \item \emoji{✅} Add Reactions（可选但推荐）
  \item \emoji{✅} Use External Emojis / Stickers（可选；仅当你需要时）
\end{itemize}


除非你在调试并完全信任机器人，否则避免使用 \textbf{Administrator}。

复制生成的 URL，打开它，选择你的服务器，然后安装机器人。

\subsubsection{4）获取 ID（服务器/用户/频道）}


Discord 到处使用数字 ID；OpenClaw 配置优先使用 ID。

\begin{enumerate}
  \item Discord（桌面/网页）→ \textbf{用户设置} → \textbf{高级} → 启用 \textbf{开发者模式}
  \item 右键点击：
\end{enumerate}

\begin{itemize}
  \item 服务器名称 → \textbf{复制服务器 ID}（服务器 ID）
  \item 频道（例如 \texttt{\#help}）→ \textbf{复制频道 ID}
  \item 你的用户 → \textbf{复制用户 ID}
\end{itemize}


\subsubsection{5）配置 OpenClaw}


\paragraph{令牌}


通过环境变量设置机器人令牌（服务器上推荐）：

\begin{itemize}
  \item \texttt{DISCORD\_BOT\_TOKEN=...}
\end{itemize}


或通过配置：

\begin{lstlisting}[language=json]
{
  channels: {
    discord: {
      enabled: true,
      token: "YOUR_BOT_TOKEN",
    },
  },
}
\end{lstlisting}


多账户支持：使用 \texttt{channels.discord.accounts}，每个账户有自己的令牌和可选的 \texttt{name}。参见 \href{/gateway/configuration\#telegramaccounts--discordaccounts--slackaccounts--signalaccounts--imessageaccounts}{\texttt{gateway/configuration}} 了解通用模式。

\paragraph{允许列表 + 频道路由}


示例"单服务器，只允许我，只允许 \#help"：

\begin{lstlisting}[language=json]
{
  channels: {
    discord: {
      enabled: true,
      dm: { enabled: false },
      guilds: {
        YOUR_GUILD_ID: {
          users: ["YOUR_USER_ID"],
          requireMention: true,
          channels: {
            help: { allow: true, requireMention: true },
          },
        },
      },
      retry: {
        attempts: 3,
        minDelayMs: 500,
        maxDelayMs: 30000,
        jitter: 0.1,
      },
    },
  },
}
\end{lstlisting}


注意：

\begin{itemize}
  \item \texttt{requireMention: true} 意味着机器人只在被提及时回复（推荐用于共享频道）。
  \item \texttt{agents.list[].groupChat.mentionPatterns}（或 \texttt{messages.groupChat.mentionPatterns}）对于服务器消息也算作提及。
  \item 多智能体覆盖：在 \texttt{agents.list[].groupChat.mentionPatterns} 上设置每个智能体的模式。
  \item 如果存在 \texttt{channels}，任何未列出的频道默认被拒绝。
  \item 使用 \texttt{"*"} 频道条目在所有频道应用默认值；显式频道条目覆盖通配符。
  \item 话题继承父频道配置（允许列表、\texttt{requireMention}、Skills、提示词等），除非你显式添加话题频道 ID。
  \item 机器人发送的消息默认被忽略；设置 \texttt{channels.discord.allowBots=true} 允许它们（自己的消息仍被过滤）。
  \item 警告：如果你允许回复其他机器人（\texttt{channels.discord.allowBots=true}），请使用 \texttt{requireMention}、\texttt{channels.discord.guilds.*.channels..users} 允许列表和/或在 \texttt{AGENTS.md} 和 \texttt{SOUL.md} 中设置明确的防护措施来防止机器人之间的回复循环。
\end{itemize}


\subsubsection{6）验证是否工作}


\begin{enumerate}
  \item 启动 Gateway 网关。
  \item 在你的服务器频道中发送：\texttt{@Krill hello}（或你的机器人名称）。
  \item 如果没有反应：查看下面的\textbf{故障排除}。
\end{enumerate}


\subsubsection{故障排除}


\begin{itemize}
  \item 首先：运行 \texttt{openclaw doctor} 和 \texttt{openclaw channels status --probe}（可操作的警告 + 快速审计）。
  \item \textbf{"Used disallowed intents"}：在开发者门户中启用 \textbf{Message Content Intent}（可能还需要 \textbf{Server Members Intent}），然后重启 Gateway 网关。
  \item \textbf{机器人连接但从不在服务器频道回复}：
  \item 缺少 \textbf{Message Content Intent}，或
  \item 机器人缺少频道权限（View/Send/Read History），或
  \item 你的配置需要提及但你没有提及它，或
  \item 你的服务器/频道允许列表拒绝了该频道/用户。
  \item \textbf{\texttt{requireMention: false} 但仍然没有回复}：
  \item \texttt{channels.discord.groupPolicy} 默认为 \textbf{allowlist}；将其设置为 \texttt{"open"} 或在 \texttt{channels.discord.guilds} 下添加服务器条目（可选择在 \texttt{channels.discord.guilds..channels} 下列出频道以进行限制）。
  \item 如果你只设置了 \texttt{DISCORD\_BOT\_TOKEN} 而从未创建 \texttt{channels.discord} 部分，运行时会将 \texttt{groupPolicy} 默认为 \texttt{open}。添加 \texttt{channels.discord.groupPolicy}、\texttt{channels.defaults.groupPolicy} 或服务器/频道允许列表来锁定它。
  \item \texttt{requireMention} 必须位于 \texttt{channels.discord.guilds}（或特定频道）下。顶层的 \texttt{channels.discord.requireMention} 会被忽略。
  \item \textbf{权限审计}（\texttt{channels status --probe}）只检查数字频道 ID。如果你使用 slug/名称作为 \texttt{channels.discord.guilds.*.channels} 键，审计无法验证权限。
  \item \textbf{私信不工作}：\texttt{channels.discord.dm.enabled=false}、\texttt{channels.discord.dm.policy="disabled"}，或者你尚未被批准（\texttt{channels.discord.dm.policy="pairing"}）。
  \item \textbf{Discord 中的执行审批}：Discord 支持私信中执行审批的\textbf{按钮 UI}（允许一次 / 始终允许 / 拒绝）。\texttt{/approve  ...} 仅用于转发的审批，不会解析 Discord 的按钮提示。如果你看到 \texttt{\emoji{❌} Failed to submit approval: Error: unknown approval id} 或 UI 从未出现，请检查：
  \item 你的配置中有 \texttt{channels.discord.execApprovals.enabled: true}。
  \item 你的 Discord 用户 ID 在 \texttt{channels.discord.execApprovals.approvers} 中列出（UI 仅发送给审批者）。
  \item 使用私信提示中的按钮（\textbf{Allow once}、\textbf{Always allow}、\textbf{Deny}）。
  \item 参见\href{/tools/exec-approvals}{执行审批}和\href{/tools/slash-commands}{斜杠命令}了解更广泛的审批和命令流程。
\end{itemize}


\subsection{功能和限制}


\begin{itemize}
  \item 支持私信和服务器文字频道（话题被视为独立频道；不支持语音）。
  \item 打字指示器尽力发送；消息分块使用 \texttt{channels.discord.textChunkLimit}（默认 2000），并按行数分割长回复（\texttt{channels.discord.maxLinesPerMessage}，默认 17）。
  \item 可选换行分块：设置 \texttt{channels.discord.chunkMode="newline"} 以在空行（段落边界）处分割，然后再进行长度分块。
  \item 支持文件上传，最大 \texttt{channels.discord.mediaMaxMb}（默认 8 MB）。
  \item 默认服务器回复需要提及，以避免嘈杂的机器人。
  \item 当消息引用另一条消息时，会注入回复上下文（引用内容 + ID）。
  \item 原生回复线程\textbf{默认关闭}；使用 \texttt{channels.discord.replyToMode} 和回复标签启用。
\end{itemize}


\subsection{重试策略}


出站 Discord API 调用在速率限制（429）时使用 Discord \texttt{retry\_after}（如果可用）进行重试，采用指数退避和抖动。通过 \texttt{channels.discord.retry} 配置。参见\href{/concepts/retry}{重试策略}。

\subsection{配置}


\begin{lstlisting}[language=json]
{
  channels: {
    discord: {
      enabled: true,
      token: "abc.123",
      groupPolicy: "allowlist",
      guilds: {
        "*": {
          channels: {
            general: { allow: true },
          },
        },
      },
      mediaMaxMb: 8,
      actions: {
        reactions: true,
        stickers: true,
        emojiUploads: true,
        stickerUploads: true,
        polls: true,
        permissions: true,
        messages: true,
        threads: true,
        pins: true,
        search: true,
        memberInfo: true,
        roleInfo: true,
        roles: false,
        channelInfo: true,
        channels: true,
        voiceStatus: true,
        events: true,
        moderation: false,
      },
      replyToMode: "off",
      dm: {
        enabled: true,
        policy: "pairing", // pairing | allowlist | open | disabled
        allowFrom: ["123456789012345678", "steipete"],
        groupEnabled: false,
        groupChannels: ["openclaw-dm"],
      },
      guilds: {
        "*": { requireMention: true },
        "123456789012345678": {
          slug: "friends-of-openclaw",
          requireMention: false,
          reactionNotifications: "own",
          users: ["987654321098765432", "steipete"],
          channels: {
            general: { allow: true },
            help: {
              allow: true,
              requireMention: true,
              users: ["987654321098765432"],
              skills: ["search", "docs"],
              systemPrompt: "Keep answers short.",
            },
          },
        },
      },
    },
  },
}
\end{lstlisting}


确认表情反应通过 \texttt{messages.ackReaction} + \texttt{messages.ackReactionScope} 全局控制。使用 \texttt{messages.removeAckAfterReply} 在机器人回复后清除确认表情反应。

\begin{itemize}
  \item \texttt{dm.enabled}：设置 \texttt{false} 忽略所有私信（默认 \texttt{true}）。
  \item \texttt{dm.policy}：私信访问控制（推荐 \texttt{pairing}）。\texttt{"open"} 需要 \texttt{dm.allowFrom=["*"]}。
  \item \texttt{dm.allowFrom}：私信允许列表（用户 ID 或名称）。用于 \texttt{dm.policy="allowlist"} 和 \texttt{dm.policy="open"} 验证。向导接受用户名，并在机器人可以搜索成员时将其解析为 ID。
  \item \texttt{dm.groupEnabled}：启用群组私信（默认 \texttt{false}）。
  \item \texttt{dm.groupChannels}：群组私信频道 ID 或 slug 的可选允许列表。
  \item \texttt{groupPolicy}：控制服务器频道处理（\texttt{open|disabled|allowlist}）；\texttt{allowlist} 需要频道允许列表。
  \item \texttt{guilds}：按服务器规则，以服务器 ID（首选）或 slug 为键。
  \item \texttt{guilds."*"}：当没有显式条目时应用的默认每服务器设置。
  \item \texttt{guilds..slug}：用于显示名称的可选友好 slug。
  \item \texttt{guilds..users}：可选的每服务器用户允许列表（ID 或名称）。
  \item \texttt{guilds..tools}：可选的每服务器工具策略覆盖（\texttt{allow}/\texttt{deny}/\texttt{alsoAllow}），在频道覆盖缺失时使用。
  \item \texttt{guilds..toolsBySender}：服务器级别的可选每发送者工具策略覆盖（在频道覆盖缺失时应用；支持 \texttt{"*"} 通配符）。
  \item \texttt{guilds..channels..allow}：当 \texttt{groupPolicy="allowlist"} 时允许/拒绝频道。
  \item \texttt{guilds..channels..requireMention}：频道的提及限制。
  \item \texttt{guilds..channels..tools}：可选的每频道工具策略覆盖（\texttt{allow}/\texttt{deny}/\texttt{alsoAllow}）。
  \item \texttt{guilds..channels..toolsBySender}：频道内的可选每发送者工具策略覆盖（支持 \texttt{"*"} 通配符）。
  \item \texttt{guilds..channels..users}：可选的每频道用户允许列表。
  \item \texttt{guilds..channels..skills}：Skills 过滤器（省略 = 所有 Skills，空 = 无）。
  \item \texttt{guilds..channels..systemPrompt}：频道的额外系统提示词（与频道主题组合）。
  \item \texttt{guilds..channels..enabled}：设置 \texttt{false} 禁用频道。
  \item \texttt{guilds..channels}：频道规则（键为频道 slug 或 ID）。
  \item \texttt{guilds..requireMention}：每服务器提及要求（可按频道覆盖）。
  \item \texttt{guilds..reactionNotifications}：表情反应系统事件模式（\texttt{off}、\texttt{own}、\texttt{all}、\texttt{allowlist}）。
  \item \texttt{textChunkLimit}：出站文本块大小（字符）。默认：2000。
  \item \texttt{chunkMode}：\texttt{length}（默认）仅在超过 \texttt{textChunkLimit} 时分割；\texttt{newline} 在空行（段落边界）处分割，然后再进行长度分块。
  \item \texttt{maxLinesPerMessage}：每条消息的软最大行数。默认：17。
  \item \texttt{mediaMaxMb}：限制保存到磁盘的入站媒体大小。
  \item \texttt{historyLimit}：回复提及时作为上下文包含的最近服务器消息数量（默认 20；回退到 \texttt{messages.groupChat.historyLimit}；\texttt{0} 禁用）。
  \item \texttt{dmHistoryLimit}：私信历史限制（用户轮次）。每用户覆盖：\texttt{dms[""].historyLimit}。
  \item \texttt{retry}：出站 Discord API 调用的重试策略（attempts、minDelayMs、maxDelayMs、jitter）。
  \item \texttt{pluralkit}：解析 PluralKit 代理消息，使系统成员显示为不同的发送者。
  \item \texttt{actions}：每操作工具门控；省略允许所有（设置 \texttt{false} 禁用）。
  \item \texttt{reactions}（涵盖表情反应 + 读取表情反应）
  \item \texttt{stickers}、\texttt{emojiUploads}、\texttt{stickerUploads}、\texttt{polls}、\texttt{permissions}、\texttt{messages}、\texttt{threads}、\texttt{pins}、\texttt{search}
  \item \texttt{memberInfo}、\texttt{roleInfo}、\texttt{channelInfo}、\texttt{voiceStatus}、\texttt{events}
  \item \texttt{channels}（创建/编辑/删除频道 + 类别 + 权限）
  \item \texttt{roles}（角色添加/移除，默认 \texttt{false}）
  \item \texttt{moderation}（超时/踢出/封禁，默认 \texttt{false}）
  \item \texttt{execApprovals}：Discord 专用执行审批私信（按钮 UI）。支持 \texttt{enabled}、\texttt{approvers}、\texttt{agentFilter}、\texttt{sessionFilter}。
\end{itemize}


表情反应通知使用 \texttt{guilds..reactionNotifications}：

\begin{itemize}
  \item \texttt{off}：无表情反应事件。
  \item \texttt{own}：机器人自己消息上的表情反应（默认）。
  \item \texttt{all}：所有消息上的所有表情反应。
  \item \texttt{allowlist}：来自 \texttt{guilds..users} 的用户在所有消息上的表情反应（空列表禁用）。
\end{itemize}


\subsubsection{PluralKit（PK）支持}


启用 PK 查找，以便代理消息解析到底层系统 + 成员。启用后，OpenClaw 使用成员身份进行允许列表匹配，并将发送者标记为 \texttt{Member (PK:System)} 以避免意外的 Discord 提及。

\begin{lstlisting}[language=json]
{
  channels: {
    discord: {
      pluralkit: {
        enabled: true,
        token: "pk_live_...", // 可选；私有系统需要
      },
    },
  },
}
\end{lstlisting}


允许列表注意事项（启用 PK 时）：

\begin{itemize}
  \item 在 \texttt{dm.allowFrom}、\texttt{guilds..users} 或每频道 \texttt{users} 中使用 \texttt{pk:}。
  \item 成员显示名称也按名称/slug 匹配。
  \item 查找使用\textbf{原始} Discord 消息 ID（代理前的消息），因此 PK API 只在其 30 分钟窗口内解析它。
  \item 如果 PK 查找失败（例如，没有令牌的私有系统），代理消息会被视为机器人消息并被丢弃，除非 \texttt{channels.discord.allowBots=true}。
\end{itemize}


\subsubsection{工具操作默认值}



\begin{table}[h]
\centering
\small
\begin{tabular}{|l|l|l|}
\hline
操作组 & 默认 & 说明 \\
\hline
reactions & 启用 & 表情反应 + 列出表情反应 + emojiList \\
\hline
stickers & 启用 & 发送贴纸 \\
\hline
emojiUploads & 启用 & 上传表情 \\
\hline
stickerUploads & 启用 & 上传贴纸 \\
\hline
polls & 启用 & 创建投票 \\
\hline
permissions & 启用 & 频道权限快照 \\
\hline
messages & 启用 & 读取/发送/编辑/删除 \\
\hline
threads & 启用 & 创建/列出/回复 \\
\hline
pins & 启用 & 置顶/取消置顶/列出 \\
\hline
search & 启用 & 消息搜索（预览功能） \\
\hline
memberInfo & 启用 & 成员信息 \\
\hline
roleInfo & 启用 & 角色列表 \\
\hline
channelInfo & 启用 & 频道信息 + 列表 \\
\hline
channels & 启用 & 频道/类别管理 \\
\hline
voiceStatus & 启用 & 语音状态查询 \\
\hline
events & 启用 & 列出/创建预定事件 \\
\hline
roles & 禁用 & 角色添加/移除 \\
\hline
moderation & 禁用 & 超时/踢出/封禁 \\
\hline
\end{tabular}
\end{table}



\begin{itemize}
  \item \texttt{replyToMode}：\texttt{off}（默认）、\texttt{first} 或 \texttt{all}。仅在模型包含回复标签时适用。
\end{itemize}


\subsection{回复标签}


要请求线程回复，模型可以在其输出中包含一个标签：

\begin{itemize}
  \item \texttt{[[reply\_to\_current]]} — 回复触发的 Discord 消息。
  \item \texttt{[[reply\_to:]]} — 回复上下文/历史中的特定消息 ID。当前消息 ID 作为 \texttt{[message\_id: …]} 附加到提示词；历史条目已包含 ID。
\end{itemize}


行为由 \texttt{channels.discord.replyToMode} 控制：

\begin{itemize}
  \item \texttt{off}：忽略标签。
  \item \texttt{first}：只有第一个出站块/附件是回复。
  \item \texttt{all}：每个出站块/附件都是回复。
\end{itemize}


允许列表匹配注意事项：

\begin{itemize}
  \item \texttt{allowFrom}/\texttt{users}/\texttt{groupChannels} 接受 ID、名称、标签或像 \texttt{<@id>} 这样的提及。
  \item 支持 \texttt{discord:}/\texttt{user:}（用户）和 \texttt{channel:}（群组私信）等前缀。
  \item 使用 \texttt{*} 允许任何发送者/频道。
  \item 当存在 \texttt{guilds..channels} 时，未列出的频道默认被拒绝。
  \item 当省略 \texttt{guilds..channels} 时，允许列表中服务器的所有频道都被允许。
  \item 要\textbf{不允许任何频道}，设置 \texttt{channels.discord.groupPolicy: "disabled"}（或保持空允许列表）。
  \item 配置向导接受 \texttt{Guild/Channel} 名称（公开 + 私有）并在可能时将其解析为 ID。
  \item 启动时，OpenClaw 将允许列表中的频道/用户名称解析为 ID（当机器人可以搜索成员时）并记录映射；未解析的条目保持原样。
\end{itemize}


原生命令注意事项：

\begin{itemize}
  \item 注册的命令镜像 OpenClaw 的聊天命令。
  \item 原生命令遵循与私信/服务器消息相同的允许列表（\texttt{channels.discord.dm.allowFrom}、\texttt{channels.discord.guilds}、每频道规则）。
  \item 斜杠命令可能在 Discord UI 中对未在允许列表中的用户仍然可见；OpenClaw 在执行时强制执行允许列表并回复"未授权"。
\end{itemize}


\subsection{工具操作}


智能体可以使用以下操作调用 \texttt{discord}：

\begin{itemize}
  \item \texttt{react} / \texttt{reactions}（添加或列出表情反应）
  \item \texttt{sticker}、\texttt{poll}、\texttt{permissions}
  \item \texttt{readMessages}、\texttt{sendMessage}、\texttt{editMessage}、\texttt{deleteMessage}
  \item 读取/搜索/置顶工具负载包含规范化的 \texttt{timestampMs}（UTC 纪元毫秒）和 \texttt{timestampUtc} 以及原始 Discord \texttt{timestamp}。
  \item \texttt{threadCreate}、\texttt{threadList}、\texttt{threadReply}
  \item \texttt{pinMessage}、\texttt{unpinMessage}、\texttt{listPins}
  \item \texttt{searchMessages}、\texttt{memberInfo}、\texttt{roleInfo}、\texttt{roleAdd}、\texttt{roleRemove}、\texttt{emojiList}
  \item \texttt{channelInfo}、\texttt{channelList}、\texttt{voiceStatus}、\texttt{eventList}、\texttt{eventCreate}
  \item \texttt{timeout}、\texttt{kick}、\texttt{ban}
\end{itemize}


Discord 消息 ID 在注入的上下文中显示（\texttt{[discord message id: …]} 和历史行），以便智能体可以定位它们。
表情可以是 unicode（例如 \texttt{\emoji{✅}}）或自定义表情语法如 \texttt{<:party\_blob:1234567890>}。

\subsection{安全与运维}


\begin{itemize}
  \item 像对待密码一样对待机器人令牌；在受监督的主机上优先使用 \texttt{DISCORD\_BOT\_TOKEN} 环境变量，或锁定配置文件权限。
  \item 只授予机器人所需的权限（通常是读取/发送消息）。
  \item 如果机器人卡住或受到速率限制，在确认没有其他进程拥有 Discord 会话后重启 Gateway 网关（\texttt{openclaw gateway --force}）。
\end{itemize}



\clearpage

% === 源文件: channels/feishu.md ===
\section{飞书机器人}


状态：生产就绪，支持机器人私聊和群组。使用 WebSocket 长连接模式接收消息。

\hrulefill


\subsection{需要插件}


安装 Feishu 插件：

\begin{lstlisting}[language=bash]
openclaw plugins install @openclaw/feishu
\end{lstlisting}


本地 checkout（在 git 仓库内运行）：

\begin{lstlisting}[language=bash]
openclaw plugins install ./extensions/feishu
\end{lstlisting}


\hrulefill


\subsection{快速开始}


添加飞书渠道有两种方式：

\subsubsection{方式一：通过安装向导添加（推荐）}


如果您刚安装完 OpenClaw，可以直接运行向导，根据提示添加飞书：

\begin{lstlisting}[language=bash]
openclaw onboard
\end{lstlisting}


向导会引导您完成：

\begin{enumerate}
  \item 创建飞书应用并获取凭证
  \item 配置应用凭证
  \item 启动网关
\end{enumerate}


\emoji{✅} \textbf{完成配置后}，您可以使用以下命令检查网关状态：

\begin{itemize}
  \item \texttt{openclaw gateway status} - 查看网关运行状态
  \item \texttt{openclaw logs --follow} - 查看实时日志
\end{itemize}


\subsubsection{方式二：通过命令行添加}


如果您已经完成了初始安装，可以用以下命令添加飞书渠道：

\begin{lstlisting}[language=bash]
openclaw channels add
\end{lstlisting}


然后根据交互式提示选择 Feishu，输入 App ID 和 App Secret 即可。

\emoji{✅} \textbf{完成配置后}，您可以使用以下命令管理网关：

\begin{itemize}
  \item \texttt{openclaw gateway status} - 查看网关运行状态
  \item \texttt{openclaw gateway restart} - 重启网关以应用新配置
  \item \texttt{openclaw logs --follow} - 查看实时日志
\end{itemize}


\hrulefill


\subsection{第一步：创建飞书应用}


\subsubsection{1. 打开飞书开放平台}


访问 \href{https://open.feishu.cn/app}{飞书开放平台}，使用飞书账号登录。

Lark（国际版）请使用 https://open.larksuite.com/app，并在配置中设置 \texttt{domain: "lark"}。

\subsubsection{2. 创建应用}


\begin{enumerate}
  \item 点击 \textbf{创建企业自建应用}
  \item 填写应用名称和描述
  \item 选择应用图标
\end{enumerate}


% [图片: 创建企业自建应用]

\subsubsection{3. 获取应用凭证}


在应用的 \textbf{凭证与基础信息} 页面，复制：

\begin{itemize}
  \item \textbf{App ID}（格式如 \texttt{cli\_xxx}）
  \item \textbf{App Secret}
\end{itemize}


\emoji{❗} \textbf{重要}：请妥善保管 App Secret，不要分享给他人。

% [图片: 获取应用凭证]

\subsubsection{4. 配置应用权限}


在 \textbf{权限管理} 页面，点击 \textbf{批量导入} 按钮，粘贴以下 JSON 配置一键导入所需权限：

\begin{lstlisting}[language=json]
{
  "scopes": {
    "tenant": [
      "aily:file:read",
      "aily:file:write",
      "application:application.app_message_stats.overview:readonly",
      "application:application:self_manage",
      "application:bot.menu:write",
      "cardkit:card:write",
      "contact:user.employee_id:readonly",
      "corehr:file:download",
      "docs:document.content:read",
      "event:ip_list",
      "im:chat",
      "im:chat.access_event.bot_p2p_chat:read",
      "im:chat.members:bot_access",
      "im:message",
      "im:message.group_at_msg:readonly",
      "im:message.group_msg",
      "im:message.p2p_msg:readonly",
      "im:message:readonly",
      "im:message:send_as_bot",
      "im:resource",
      "sheets:spreadsheet",
      "wiki:wiki:readonly"
    ],
    "user": ["aily:file:read", "aily:file:write", "im:chat.access_event.bot_p2p_chat:read"]
  }
}
\end{lstlisting}


% [图片: 配置应用权限]

\subsubsection{5. 启用机器人能力}


在 \textbf{应用能力} > \textbf{机器人} 页面：

\begin{enumerate}
  \item 开启机器人能力
  \item 配置机器人名称
\end{enumerate}


% [图片: 启用机器人能力]

\subsubsection{6. 配置事件订阅}


\emoji{⚠} \textbf{重要提醒}：在配置事件订阅前，请务必确保已完成以下步骤：

\begin{enumerate}
  \item 运行 \texttt{openclaw channels add} 添加了 Feishu 渠道
  \item 网关处于启动状态（可通过 \texttt{openclaw gateway status} 检查状态）
\end{enumerate}


在 \textbf{事件订阅} 页面：

\begin{enumerate}
  \item 选择 \textbf{使用长连接接收事件}（WebSocket 模式）
  \item 添加事件：\texttt{im.message.receive\_v1}（接收消息）
\end{enumerate}


\emoji{⚠} \textbf{注意}：如果网关未启动或渠道未添加，长连接设置将保存失败。

% [图片: 配置事件订阅]

\subsubsection{7. 发布应用}


\begin{enumerate}
  \item 在 \textbf{版本管理与发布} 页面创建版本
  \item 提交审核并发布
  \item 等待管理员审批（企业自建应用通常自动通过）
\end{enumerate}


\hrulefill


\subsection{第二步：配置 OpenClaw}


\subsubsection{通过向导配置（推荐）}


运行以下命令，根据提示粘贴 App ID 和 App Secret：

\begin{lstlisting}[language=bash]
openclaw channels add
\end{lstlisting}


选择 \textbf{Feishu}，然后输入您在第一步获取的凭证即可。

\subsubsection{通过配置文件配置}


编辑 \texttt{\textasciitilde{}/.openclaw/openclaw.json}：

\begin{lstlisting}[language=json]
{
  channels: {
    feishu: {
      enabled: true,
      dmPolicy: "pairing",
      accounts: {
        main: {
          appId: "cli_xxx",
          appSecret: "xxx",
          botName: "我的AI助手",
        },
      },
    },
  },
}
\end{lstlisting}


若使用 \texttt{connectionMode: "webhook"}，需设置 \texttt{verificationToken}。飞书 Webhook 服务默认绑定 \texttt{127.0.0.1}；仅在需要不同监听地址时设置 \texttt{webhookHost}。

\paragraph{获取 Verification Token（仅 Webhook 模式）}


使用 Webhook 模式时，需在配置中设置 \texttt{channels.feishu.verificationToken}。获取方式：

\begin{enumerate}
  \item 在飞书开放平台打开您的应用
  \item 进入 \textbf{开发配置} → \textbf{事件与回调}
  \item 打开 \textbf{加密策略} 选项卡
  \item 复制 \textbf{Verification Token}（校验令牌）
\end{enumerate}


% [图片: Verification Token 位置]

\subsubsection{通过环境变量配置}


\begin{lstlisting}[language=bash]
export FEISHU_APP_ID="cli_xxx"
export FEISHU_APP_SECRET="xxx"
\end{lstlisting}


\subsubsection{Lark（国际版）域名}


如果您的租户在 Lark（国际版），请设置域名为 \texttt{lark}（或完整域名），可配置 \texttt{channels.feishu.domain} 或 \texttt{channels.feishu.accounts..domain}：

\begin{lstlisting}[language=json]
{
  channels: {
    feishu: {
      domain: "lark",
      accounts: {
        main: {
          appId: "cli_xxx",
          appSecret: "xxx",
        },
      },
    },
  },
}
\end{lstlisting}


\subsubsection{配额优化}


可通过以下可选配置减少飞书 API 调用：

\begin{itemize}
  \item \texttt{typingIndicator}（默认 \texttt{true}）：设为 \texttt{false} 时不发送“正在输入”状态。
  \item \texttt{resolveSenderNames}（默认 \texttt{true}）：设为 \texttt{false} 时不拉取发送者资料。
\end{itemize}


可在渠道级或账号级配置：

\begin{lstlisting}[language=json]
{
  channels: {
    feishu: {
      typingIndicator: false,
      resolveSenderNames: false,
      accounts: {
        main: {
          appId: "cli_xxx",
          appSecret: "xxx",
          typingIndicator: true,
          resolveSenderNames: false,
        },
      },
    },
  },
}
\end{lstlisting}


\hrulefill


\subsection{第三步：启动并测试}


\subsubsection{1. 启动网关}


\begin{lstlisting}[language=bash]
openclaw gateway
\end{lstlisting}


\subsubsection{2. 发送测试消息}


在飞书中找到您创建的机器人，发送一条消息。

\subsubsection{3. 配对授权}


默认情况下，机器人会回复一个 \textbf{配对码}。您需要批准此代码：

\begin{lstlisting}[language=bash]
openclaw pairing approve feishu <配对码>
\end{lstlisting}


批准后即可正常对话。

\hrulefill


\subsection{介绍}


\begin{itemize}
  \item \textbf{飞书机器人渠道}：由网关管理的飞书机器人
  \item \textbf{确定性路由}：回复始终返回飞书，模型不会选择渠道
  \item \textbf{会话隔离}：私聊共享主会话；群组独立隔离
  \item \textbf{WebSocket 连接}：使用飞书 SDK 的长连接模式，无需公网 URL
\end{itemize}


\hrulefill


\subsection{访问控制}


\subsubsection{私聊访问}


\begin{itemize}
  \item \textbf{默认}：\texttt{dmPolicy: "pairing"}，陌生用户会收到配对码
  \item \textbf{批准配对}：
\begin{lstlisting}[language=bash]
  openclaw pairing list feishu      # 查看待审批列表
  openclaw pairing approve feishu <CODE>  # 批准
\end{lstlisting}

  \item \textbf{白名单模式}：通过 \texttt{channels.feishu.allowFrom} 配置允许的用户 Open ID
\end{itemize}


\subsubsection{群组访问}


\textbf{1. 群组策略}（\texttt{channels.feishu.groupPolicy}）：

\begin{itemize}
  \item \texttt{"open"} = 允许群组中所有人（默认）
  \item \texttt{"allowlist"} = 仅允许 \texttt{groupAllowFrom} 中的群组
  \item \texttt{"disabled"} = 禁用群组消息
\end{itemize}


\textbf{2. @提及要求}（\texttt{channels.feishu.groups..requireMention}）：

\begin{itemize}
  \item \texttt{true} = 需要 @机器人才响应（默认）
  \item \texttt{false} = 无需 @也响应
\end{itemize}


\hrulefill


\subsection{群组配置示例}


\subsubsection{允许所有群组，需要 @提及（默认行为）}


\begin{lstlisting}[language=json]
{
  channels: {
    feishu: {
      groupPolicy: "open",
      // 默认 requireMention: true
    },
  },
}
\end{lstlisting}


\subsubsection{允许所有群组，无需 @提及}


需要为特定群组配置：

\begin{lstlisting}[language=json]
{
  channels: {
    feishu: {
      groups: {
        oc_xxx: { requireMention: false },
      },
    },
  },
}
\end{lstlisting}


\subsubsection{仅允许特定群组}


\begin{lstlisting}[language=json]
{
  channels: {
    feishu: {
      groupPolicy: "allowlist",
      // 群组 ID 格式为 oc_xxx
      groupAllowFrom: ["oc_xxx", "oc_yyy"],
    },
  },
}
\end{lstlisting}


\subsubsection{仅允许特定成员在群组中发信（发送者白名单）}


除群组白名单外，该群组内\textbf{所有消息}均按发送者 open\_id 校验：仅 \texttt{groups..allowFrom} 中列出的用户消息会被处理，其他成员的消息会被忽略（此为发送者级白名单，不仅针对 /reset、/new 等控制命令）。

\begin{lstlisting}[language=json]
{
  channels: {
    feishu: {
      groupPolicy: "allowlist",
      groupAllowFrom: ["oc_xxx"],
      groups: {
        oc_xxx: {
          // 用户 open_id 格式为 ou_xxx
          allowFrom: ["ou_user1", "ou_user2"],
        },
      },
    },
  },
}
\end{lstlisting}


\hrulefill


\subsection{获取群组/用户 ID}


\subsubsection{获取群组 ID（chat\_id）}


群组 ID 格式为 \texttt{oc\_xxx}，可以通过以下方式获取：

\textbf{方法一}（推荐）：

\begin{enumerate}
  \item 启动网关并在群组中 @机器人发消息
  \item 运行 \texttt{openclaw logs --follow} 查看日志中的 \texttt{chat\_id}
\end{enumerate}


\textbf{方法二}：
使用飞书 API 调试工具获取机器人所在群组列表。

\subsubsection{获取用户 ID（open\_id）}


用户 ID 格式为 \texttt{ou\_xxx}，可以通过以下方式获取：

\textbf{方法一}（推荐）：

\begin{enumerate}
  \item 启动网关并给机器人发消息
  \item 运行 \texttt{openclaw logs --follow} 查看日志中的 \texttt{open\_id}
\end{enumerate}


\textbf{方法二}：
查看配对请求列表，其中包含用户的 Open ID：

\begin{lstlisting}[language=bash]
openclaw pairing list feishu
\end{lstlisting}


\hrulefill


\subsection{常用命令}



\begin{table}[h]
\centering
\small
\begin{tabular}{|l|l|}
\hline
命令 & 说明 \\
\hline
\texttt{/status} & 查看机器人状态 \\
\hline
\texttt{/reset} & 重置对话会话 \\
\hline
\texttt{/model} & 查看/切换模型 \\
\hline
\end{tabular}
\end{table}



\begin{quote}
注意：飞书目前不支持原生命令菜单，命令需要以文本形式发送。
\end{quote}


\subsection{网关管理命令}


在配置和使用飞书渠道时，您可能需要使用以下网关管理命令：


\begin{table}[h]
\centering
\small
\begin{tabular}{|l|l|}
\hline
命令 & 说明 \\
\hline
\texttt{openclaw gateway status} & 查看网关运行状态 \\
\hline
\texttt{openclaw gateway install} & 安装/启动网关服务 \\
\hline
\texttt{openclaw gateway stop} & 停止网关服务 \\
\hline
\texttt{openclaw gateway restart} & 重启网关服务 \\
\hline
\texttt{openclaw logs --follow} & 实时查看日志输出 \\
\hline
\end{tabular}
\end{table}



\hrulefill


\subsection{故障排除}


\subsubsection{机器人在群组中不响应}


\begin{enumerate}
  \item 检查机器人是否已添加到群组
  \item 检查是否 @了机器人（默认需要 @提及）
  \item 检查 \texttt{groupPolicy} 是否为 \texttt{"disabled"}
  \item 查看日志：\texttt{openclaw logs --follow}
\end{enumerate}


\subsubsection{机器人收不到消息}


\begin{enumerate}
  \item 检查应用是否已发布并审批通过
  \item 检查事件订阅是否配置正确（\texttt{im.message.receive\_v1}）
  \item 检查是否选择了 \textbf{长连接} 模式
  \item 检查应用权限是否完整
  \item 检查网关是否正在运行：\texttt{openclaw gateway status}
  \item 查看实时日志：\texttt{openclaw logs --follow}
\end{enumerate}


\subsubsection{App Secret 泄露怎么办}


\begin{enumerate}
  \item 在飞书开放平台重置 App Secret
  \item 更新配置文件中的 App Secret
  \item 重启网关
\end{enumerate}


\subsubsection{发送消息失败}


\begin{enumerate}
  \item 检查应用是否有 \texttt{im:message:send\_as\_bot} 权限
  \item 检查应用是否已发布
  \item 查看日志获取详细错误信息
\end{enumerate}


\hrulefill


\subsection{高级配置}


\subsubsection{多账号配置}


如果需要管理多个飞书机器人，可配置 \texttt{defaultAccount} 指定出站未显式指定 \texttt{accountId} 时使用的账号：

\begin{lstlisting}[language=json]
{
  channels: {
    feishu: {
      defaultAccount: "main",
      accounts: {
        main: {
          appId: "cli_xxx",
          appSecret: "xxx",
          botName: "主机器人",
        },
        backup: {
          appId: "cli_yyy",
          appSecret: "yyy",
          botName: "备用机器人",
          enabled: false, // 暂时禁用
        },
      },
    },
  },
}
\end{lstlisting}


\subsubsection{消息限制}


\begin{itemize}
  \item \texttt{textChunkLimit}：出站文本分块大小（默认 2000 字符）
  \item \texttt{mediaMaxMb}：媒体上传/下载限制（默认 30MB）
\end{itemize}


\subsubsection{流式输出}


飞书支持通过交互式卡片实现流式输出，机器人会实时更新卡片内容显示生成进度。默认配置：

\begin{lstlisting}[language=json]
{
  channels: {
    feishu: {
      streaming: true, // 启用流式卡片输出（默认 true）
      blockStreaming: true, // 启用块级流式（默认 true）
    },
  },
}
\end{lstlisting}


如需禁用流式输出（等待完整回复后一次性发送），可设置 \texttt{streaming: false}。

\subsubsection{消息引用}


在群聊中，机器人的回复可以引用用户发送的原始消息，让对话上下文更加清晰。

配置选项：

\begin{lstlisting}[language=json]
{
  channels: {
    feishu: {
      // 账户级别配置（默认 "all"）
      replyToMode: "all",
      groups: {
        oc_xxx: {
          // 特定群组可以覆盖
          replyToMode: "first",
        },
      },
    },
  },
}
\end{lstlisting}


\texttt{replyToMode} 值说明：


\begin{table}[h]
\centering
\small
\begin{tabular}{|l|l|}
\hline
值 & 行为 \\
\hline
\texttt{"off"} & 不引用原消息（私聊默认值） \\
\hline
\texttt{"first"} & 仅在第一条回复时引用原消息 \\
\hline
\texttt{"all"} & 所有回复都引用原消息（群聊默认值） \\
\hline
\end{tabular}
\end{table}



\begin{quote}
注意：消息引用功能与流式卡片输出（\texttt{streaming: true}）不能同时使用。当启用流式输出时，回复会以卡片形式呈现，不会显示引用。
\end{quote}


\subsubsection{多 Agent 路由}


通过 \texttt{bindings} 配置，您可以用一个飞书机器人对接多个不同功能或性格的 Agent。系统会根据用户 ID 或群组 ID 自动将对话分发到对应的 Agent。

配置示例：

\begin{lstlisting}[language=json]
{
  agents: {
    list: [
      { id: "main" },
      {
        id: "clawd-fan",
        workspace: "/home/user/clawd-fan",
        agentDir: "/home/user/.openclaw/agents/clawd-fan/agent",
      },
      {
        id: "clawd-xi",
        workspace: "/home/user/clawd-xi",
        agentDir: "/home/user/.openclaw/agents/clawd-xi/agent",
      },
    ],
  },
  bindings: [
    {
      // 用户 A 的私聊 → main agent
      agentId: "main",
      match: {
        channel: "feishu",
        peer: { kind: "dm", id: "ou_28b31a88..." },
      },
    },
    {
      // 用户 B 的私聊 → clawd-fan agent
      agentId: "clawd-fan",
      match: {
        channel: "feishu",
        peer: { kind: "dm", id: "ou_0fe6b1c9..." },
      },
    },
    {
      // 某个群组 → clawd-xi agent
      agentId: "clawd-xi",
      match: {
        channel: "feishu",
        peer: { kind: "group", id: "oc_xxx..." },
      },
    },
  ],
}
\end{lstlisting}


匹配规则说明：


\begin{table}[h]
\centering
\small
\begin{tabular}{|l|l|}
\hline
字段 & 说明 \\
\hline
\texttt{agentId} & 目标 Agent 的 ID，需要在 \texttt{agents.list} 中定义 \\
\hline
\texttt{match.channel} & 渠道类型，这里固定为 \texttt{"feishu"} \\
\hline
\texttt{match.peer.kind} & 对话类型：\texttt{"dm"}（私聊）或 \texttt{"group"}（群组） \\
\hline
\texttt{match.peer.id} & 用户 Open ID（\texttt{ou\_xxx}）或群组 ID（\texttt{oc\_xxx}） \\
\hline
\end{tabular}
\end{table}



\begin{quote}
获取 ID 的方法：参见上文 \href{\#获取群组用户-id}{获取群组/用户 ID} 章节。
\end{quote}


\hrulefill


\subsection{配置参考}


完整配置请参考：\href{/gateway/configuration}{网关配置}

主要选项：


\begin{table}[h]
\centering
\small
\begin{tabular}{|l|l|l|}
\hline
配置项 & 说明 & 默认值 \\
\hline
\texttt{channels.feishu.enabled} & 启用/禁用渠道 & \texttt{true} \\
\hline
\texttt{channels.feishu.domain} & API 域名（\texttt{feishu} 或 \texttt{lark}） & \texttt{feishu} \\
\hline
\texttt{channels.feishu.connectionMode} & 事件传输模式（websocket/webhook） & \texttt{websocket} \\
\hline
\texttt{channels.feishu.defaultAccount} & 出站路由默认账号 ID & \texttt{default} \\
\hline
\texttt{channels.feishu.verificationToken} & Webhook 模式必填 & - \\
\hline
\texttt{channels.feishu.webhookPath} & Webhook 路由路径 & \texttt{/feishu/events} \\
\hline
\texttt{channels.feishu.webhookHost} & Webhook 监听地址 & \texttt{127.0.0.1} \\
\hline
\texttt{channels.feishu.webhookPort} & Webhook 监听端口 & \texttt{3000} \\
\hline
\texttt{channels.feishu.accounts..appId} & 应用 App ID & - \\
\hline
\texttt{channels.feishu.accounts..appSecret} & 应用 App Secret & - \\
\hline
\texttt{channels.feishu.accounts..domain} & 单账号 API 域名覆盖 & \texttt{feishu} \\
\hline
\texttt{channels.feishu.dmPolicy} & 私聊策略 & \texttt{pairing} \\
\hline
\texttt{channels.feishu.allowFrom} & 私聊白名单（open\_id 列表） & - \\
\hline
\texttt{channels.feishu.groupPolicy} & 群组策略 & \texttt{open} \\
\hline
\texttt{channels.feishu.groupAllowFrom} & 群组白名单 & - \\
\hline
\texttt{channels.feishu.groups..requireMention} & 是否需要 @提及 & \texttt{true} \\
\hline
\texttt{channels.feishu.groups..enabled} & 是否启用该群组 & \texttt{true} \\
\hline
\texttt{channels.feishu.textChunkLimit} & 消息分块大小 & \texttt{2000} \\
\hline
\texttt{channels.feishu.mediaMaxMb} & 媒体大小限制 & \texttt{30} \\
\hline
\texttt{channels.feishu.streaming} & 启用流式卡片输出 & \texttt{true} \\
\hline
\texttt{channels.feishu.blockStreaming} & 启用块级流式 & \texttt{true} \\
\hline
\end{tabular}
\end{table}



\hrulefill


\subsection{dmPolicy 策略说明}



\begin{table}[h]
\centering
\small
\begin{tabular}{|l|l|}
\hline
值 & 行为 \\
\hline
\texttt{"pairing"} & \textbf{默认}。未知用户收到配对码，管理员批准后才能对话 \\
\hline
\texttt{"allowlist"} & 仅 \texttt{allowFrom} 列表中的用户可对话，其他静默忽略 \\
\hline
\texttt{"open"} & 允许所有人对话（需在 allowFrom 中加 \texttt{"*"}） \\
\hline
\texttt{"disabled"} & 完全禁止私聊 \\
\hline
\end{tabular}
\end{table}



\hrulefill


\subsection{支持的消息类型}


\subsubsection{接收}


\begin{itemize}
  \item \emoji{✅} 文本消息
  \item \emoji{✅} 富文本（帖子）
  \item \emoji{✅} 图片
  \item \emoji{✅} 文件
  \item \emoji{✅} 音频
  \item \emoji{✅} 视频
  \item \emoji{✅} 表情包
\end{itemize}


\subsubsection{发送}


\begin{itemize}
  \item \emoji{✅} 文本消息
  \item \emoji{✅} 图片
  \item \emoji{✅} 文件
  \item \emoji{✅} 音频
  \item \emoji{⚠} 富文本（部分支持）
\end{itemize}



\clearpage

% === 源文件: channels/googlechat.md ===
\section{Google Chat（Chat API）}


状态：已支持通过 Google Chat API webhooks（仅 HTTP）使用私信和空间。

\subsection{快速设置（新手）}


\begin{enumerate}
  \item 创建一个 Google Cloud 项目并启用 \textbf{Google Chat API}。
\end{enumerate}

\begin{itemize}
  \item 前往：\href{https://console.cloud.google.com/apis/api/chat.googleapis.com/credentials}{Google Chat API Credentials}
  \item 如果 API 尚未启用，请启用它。
\end{itemize}

\begin{enumerate}
  \item 创建一个\textbf{服务账号}：
\end{enumerate}

\begin{itemize}
  \item 点击 \textbf{Create Credentials} > \textbf{Service Account}。
  \item 随意命名（例如 \texttt{openclaw-chat}）。
  \item 权限留空（点击 \textbf{Continue}）。
  \item 有访问权限的主账号留空（点击 \textbf{Done}）。
\end{itemize}

\begin{enumerate}
  \item 创建并下载 \textbf{JSON 密钥}：
\end{enumerate}

\begin{itemize}
  \item 在服务账号列表中，点击刚刚创建的账号。
  \item 前往 \textbf{Keys} 标签页。
  \item 点击 \textbf{Add Key} > \textbf{Create new key}。
  \item 选择 \textbf{JSON} 并点击 \textbf{Create}。
\end{itemize}

\begin{enumerate}
  \item 将下载的 JSON 文件存储在 Gateway 网关主机上（例如 \texttt{\textasciitilde{}/.openclaw/googlechat-service-account.json}）。
  \item 在 \href{https://console.cloud.google.com/apis/api/chat.googleapis.com/hangouts-chat}{Google Cloud Console Chat Configuration} 中创建一个 Google Chat 应用：
\end{enumerate}

\begin{itemize}
  \item 填写 \textbf{Application info}：
  \item \textbf{App name}：（例如 \texttt{OpenClaw}）
  \item \textbf{Avatar URL}：（例如 \texttt{https://openclaw.ai/logo.png}）
  \item \textbf{Description}：（例如 \texttt{Personal AI Assistant}）
  \item 启用 \textbf{Interactive features}。
  \item 在 \textbf{Functionality} 下，勾选 \textbf{Join spaces and group conversations}。
  \item 在 \textbf{Connection settings} 下，选择 \textbf{HTTP endpoint URL}。
  \item 在 \textbf{Triggers} 下，选择 \textbf{Use a common HTTP endpoint URL for all triggers} 并将其设置为你的 Gateway 网关公网 URL 后加 \texttt{/googlechat}。
  \item \_提示：运行 \texttt{openclaw status} 查看你的 Gateway 网关公网 URL。\_
  \item 在 \textbf{Visibility} 下，勾选 \textbf{Make this Chat app available to specific people and groups in \&lt;Your Domain\&gt;}。
  \item 在文本框中输入你的邮箱地址（例如 \texttt{user@example.com}）。
  \item 点击底部的 \textbf{Save}。
\end{itemize}

\begin{enumerate}
  \item \textbf{启用应用状态}：
\end{enumerate}

\begin{itemize}
  \item 保存后，\textbf{刷新页面}。
  \item 找到 \textbf{App status} 部分（通常在保存后位于顶部或底部附近）。
  \item 将状态更改为 \textbf{Live - available to users}。
  \item 再次点击 \textbf{Save}。
\end{itemize}

\begin{enumerate}
  \item 使用服务账号路径和 webhook audience 配置 OpenClaw：
\end{enumerate}

\begin{itemize}
  \item 环境变量：\texttt{GOOGLE\_CHAT\_SERVICE\_ACCOUNT\_FILE=/path/to/service-account.json}
  \item 或配置：\texttt{channels.googlechat.serviceAccountFile: "/path/to/service-account.json"}。
\end{itemize}

\begin{enumerate}
  \item 设置 webhook audience 类型和值（与你的 Chat 应用配置匹配）。
  \item 启动 Gateway 网关。Google Chat 将向你的 webhook 路径发送 POST 请求。
\end{enumerate}


\subsection{添加到 Google Chat}


Gateway 网关运行后，且你的邮箱已添加到可见性列表中：

\begin{enumerate}
  \item 前往 \href{https://chat.google.com/}{Google Chat}。
  \item 点击 \textbf{Direct Messages} 旁边的 \textbf{+}（加号）图标。
  \item 在搜索栏（通常用于添加联系人的位置）中，输入你在 Google Cloud Console 中配置的 \textbf{App name}。
\end{enumerate}

\begin{itemize}
  \item \textbf{注意}：该机器人\textit{不会}出现在"Marketplace"浏览列表中，因为它是私有应用。你必须按名称搜索。
\end{itemize}

\begin{enumerate}
  \item 从结果中选择你的机器人。
  \item 点击 \textbf{Add} 或 \textbf{Chat} 开始一对一对话。
  \item 发送"Hello"来触发助手！
\end{enumerate}


\subsection{公网 URL（仅 Webhook）}


Google Chat webhooks 需要一个公网 HTTPS 端点。为了安全起见，\textbf{只将 \texttt{/googlechat} 路径暴露到互联网}。将 OpenClaw 仪表板和其他敏感端点保留在你的私有网络上。

\subsubsection{方案 A：Tailscale Funnel（推荐）}


使用 Tailscale Serve 提供私有仪表板，使用 Funnel 提供公网 webhook 路径。这样可以保持 \texttt{/} 私有，同时只暴露 \texttt{/googlechat}。

\begin{enumerate}
  \item \textbf{检查你的 Gateway 网关绑定的地址：}
\end{enumerate}


\begin{lstlisting}[language=bash]
   ss -tlnp | grep 18789
\end{lstlisting}


记下 IP 地址（例如 \texttt{127.0.0.1}、\texttt{0.0.0.0} 或你的 Tailscale IP 如 \texttt{100.x.x.x}）。

\begin{enumerate}
  \item \textbf{仅将仪表板暴露给 tailnet（端口 8443）：}
\end{enumerate}


\begin{lstlisting}[language=bash]
   # 如果绑定到 localhost（127.0.0.1 或 0.0.0.0）：
   tailscale serve --bg --https 8443 http://127.0.0.1:18789

   # 如果仅绑定到 Tailscale IP（例如 100.106.161.80）：
   tailscale serve --bg --https 8443 http://100.106.161.80:18789
\end{lstlisting}


\begin{enumerate}
  \item \textbf{仅公开暴露 webhook 路径：}
\end{enumerate}


\begin{lstlisting}[language=bash]
   # 如果绑定到 localhost（127.0.0.1 或 0.0.0.0）：
   tailscale funnel --bg --set-path /googlechat http://127.0.0.1:18789/googlechat

   # 如果仅绑定到 Tailscale IP（例如 100.106.161.80）：
   tailscale funnel --bg --set-path /googlechat http://100.106.161.80:18789/googlechat
\end{lstlisting}


\begin{enumerate}
  \item \textbf{授权节点访问 Funnel：}
\end{enumerate}

如果出现提示，请访问输出中显示的授权 URL，以在你的 tailnet 策略中为此节点启用 Funnel。

\begin{enumerate}
  \item \textbf{验证配置：}
\begin{lstlisting}[language=bash]
   tailscale serve status
   tailscale funnel status
\end{lstlisting}

\end{enumerate}


你的公网 webhook URL 将是：
\texttt{https://..ts.net/googlechat}

你的私有仪表板仅限 tailnet 访问：
\texttt{https://..ts.net:8443/}

在 Google Chat 应用配置中使用公网 URL（不带 \texttt{:8443}）。

\begin{quote}
注意：此配置在重启后会保留。如需稍后移除，请运行 \texttt{tailscale funnel reset} 和 \texttt{tailscale serve reset}。
\end{quote}


\subsubsection{方案 B：反向代理（Caddy）}


如果你使用像 Caddy 这样的反向代理，只代理特定路径：

\begin{lstlisting}
your-domain.com {
    reverse_proxy /googlechat* localhost:18789
}
\end{lstlisting}


使用此配置，任何发往 \texttt{your-domain.com/} 的请求将被忽略或返回 404，而 \texttt{your-domain.com/googlechat} 会安全地路由到 OpenClaw。

\subsubsection{方案 C：Cloudflare Tunnel}


配置你的隧道入口规则，只路由 webhook 路径：

\begin{itemize}
  \item \textbf{路径}：\texttt{/googlechat} -> \texttt{http://localhost:18789/googlechat}
  \item \textbf{默认规则}：HTTP 404（未找到）
\end{itemize}


\subsection{工作原理}


\begin{enumerate}
  \item Google Chat 向 Gateway 网关发送 webhook POST 请求。每个请求都包含一个 \texttt{Authorization: Bearer } 头。
  \item OpenClaw 根据配置的 \texttt{audienceType} + \texttt{audience} 验证令牌：
\end{enumerate}

\begin{itemize}
  \item \texttt{audienceType: "app-url"} → audience 是你的 HTTPS webhook URL。
  \item \texttt{audienceType: "project-number"} → audience 是 Cloud 项目编号。
\end{itemize}

\begin{enumerate}
  \item 消息按空间路由：
\end{enumerate}

\begin{itemize}
  \item 私信使用会话键 \texttt{agent::googlechat:dm:}。
  \item 空间使用会话键 \texttt{agent::googlechat:group:}。
\end{itemize}

\begin{enumerate}
  \item 私信访问默认为配对模式。未知发送者会收到配对码；使用以下命令批准：
\end{enumerate}

\begin{itemize}
  \item \texttt{openclaw pairing approve googlechat }
\end{itemize}

\begin{enumerate}
  \item 群组空间默认需要 @提及。如果提及检测需要应用的用户名，请使用 \texttt{botUser}。
\end{enumerate}


\subsection{目标标识符}


使用这些标识符进行消息投递和允许列表：

\begin{itemize}
  \item 私信：\texttt{users/} 或 \texttt{users/}（接受邮箱地址）。
  \item 空间：\texttt{spaces/}。
\end{itemize}


\subsection{配置要点}


\begin{lstlisting}[language=json]
{
  channels: {
    googlechat: {
      enabled: true,
      serviceAccountFile: "/path/to/service-account.json",
      audienceType: "app-url",
      audience: "https://gateway.example.com/googlechat",
      webhookPath: "/googlechat",
      botUser: "users/1234567890", // 可选；帮助提及检测
      dm: {
        policy: "pairing",
        allowFrom: ["users/1234567890", "name@example.com"],
      },
      groupPolicy: "allowlist",
      groups: {
        "spaces/AAAA": {
          allow: true,
          requireMention: true,
          users: ["users/1234567890"],
          systemPrompt: "Short answers only.",
        },
      },
      actions: { reactions: true },
      typingIndicator: "message",
      mediaMaxMb: 20,
    },
  },
}
\end{lstlisting}


注意事项：

\begin{itemize}
  \item 服务账号凭证也可以通过 \texttt{serviceAccount}（JSON 字符串）内联传递。
  \item 如果未设置 \texttt{webhookPath}，默认 webhook 路径为 \texttt{/googlechat}。
  \item 当 \texttt{actions.reactions} 启用时，可通过 \texttt{reactions} 工具和 \texttt{channels action} 使用表情回应。
  \item \texttt{typingIndicator} 支持 \texttt{none}、\texttt{message}（默认）和 \texttt{reaction}（reaction 需要用户 OAuth）。
  \item 附件通过 Chat API 下载并存储在媒体管道中（大小受 \texttt{mediaMaxMb} 限制）。
\end{itemize}


\subsection{故障排除}


\subsubsection{405 Method Not Allowed}


如果 Google Cloud Logs Explorer 显示如下错误：

\begin{lstlisting}
status code: 405, reason phrase: HTTP error response: HTTP/1.1 405 Method Not Allowed
\end{lstlisting}


这意味着 webhook 处理程序未注册。常见原因：

\begin{enumerate}
  \item \textbf{渠道未配置}：配置中缺少 \texttt{channels.googlechat} 部分。使用以下命令验证：
\end{enumerate}


\begin{lstlisting}[language=bash]
   openclaw config get channels.googlechat
\end{lstlisting}


如果返回"Config path not found"，请添加配置（参见\href{\#配置要点}{配置要点}）。

\begin{enumerate}
  \item \textbf{插件未启用}：检查插件状态：
\end{enumerate}


\begin{lstlisting}[language=bash]
   openclaw plugins list | grep googlechat
\end{lstlisting}


如果显示"disabled"，请在配置中添加 \texttt{plugins.entries.googlechat.enabled: true}。

\begin{enumerate}
  \item \textbf{Gateway 网关未重启}：添加配置后，重启 Gateway 网关：
\begin{lstlisting}[language=bash]
   openclaw gateway restart
\end{lstlisting}

\end{enumerate}


验证渠道是否正在运行：

\begin{lstlisting}[language=bash]
openclaw channels status
# 应显示：Google Chat default: enabled, configured, ...
\end{lstlisting}


\subsubsection{其他问题}


\begin{itemize}
  \item 检查 \texttt{openclaw channels status --probe} 以查看认证错误或缺少 audience 配置。
  \item 如果没有收到消息，请确认 Chat 应用的 webhook URL 和事件订阅。
  \item 如果提及门控阻止了回复，请将 \texttt{botUser} 设置为应用的用户资源名称并验证 \texttt{requireMention}。
  \item 在发送测试消息时使用 \texttt{openclaw logs --follow} 查看请求是否到达 Gateway 网关。
\end{itemize}


相关文档：

\begin{itemize}
  \item \href{/gateway/configuration}{Gateway 网关配置}
  \item \href{/gateway/security}{安全}
  \item \href{/tools/reactions}{表情回应}
\end{itemize}



\clearpage

% === 源文件: channels/grammy.md ===
\section{grammY 集成（Telegram Bot API）}


\section{为什么选择 grammY}


\begin{itemize}
  \item 以 TS 为核心的 Bot API 客户端，内置长轮询 + webhook 辅助工具、中间件、错误处理和速率限制器。
  \item 媒体处理辅助工具比手动编写 fetch + FormData 更简洁；支持所有 Bot API 方法。
  \item 可扩展：通过自定义 fetch 支持代理，可选的会话中间件，类型安全的上下文。
\end{itemize}


\section{我们发布的内容}


\begin{itemize}
  \item \textbf{单一客户端路径：} 移除了基于 fetch 的实现；grammY 现在是唯一的 Telegram 客户端（发送 + Gateway 网关），默认启用 grammY throttler。
  \item \textbf{Gateway 网关：} \texttt{monitorTelegramProvider} 构建 grammY \texttt{Bot}，接入 mention/allowlist 网关控制，通过 \texttt{getFile}/\texttt{download} 下载媒体，并使用 \texttt{sendMessage/sendPhoto/sendVideo/sendAudio/sendDocument} 发送回复。通过 \texttt{webhookCallback} 支持长轮询或 webhook。
  \item \textbf{代理：} 可选的 \texttt{channels.telegram.proxy} 通过 grammY 的 \texttt{client.baseFetch} 使用 \texttt{undici.ProxyAgent}。
  \item \textbf{Webhook 支持：} \texttt{webhook-set.ts} 封装了 \texttt{setWebhook/deleteWebhook}；\texttt{webhook.ts} 托管回调，支持健康检查和优雅关闭。当设置了 \texttt{channels.telegram.webhookUrl} + \texttt{channels.telegram.webhookSecret} 时，Gateway 网关启用 webhook 模式（否则使用长轮询）。
  \item \textbf{会话：} 私聊折叠到智能体主会话（\texttt{agent::}）；群组使用 \texttt{agent::telegram:group:}；回复路由回同一渠道。
  \item \textbf{配置选项：} \texttt{channels.telegram.botToken}、\texttt{channels.telegram.dmPolicy}、\texttt{channels.telegram.groups}（allowlist + mention 默认值）、\texttt{channels.telegram.allowFrom}、\texttt{channels.telegram.groupAllowFrom}、\texttt{channels.telegram.groupPolicy}、\texttt{channels.telegram.mediaMaxMb}、\texttt{channels.telegram.linkPreview}、\texttt{channels.telegram.proxy}、\texttt{channels.telegram.webhookSecret}、\texttt{channels.telegram.webhookUrl}。
  \item \textbf{草稿流式传输：} 可选的 \texttt{channels.telegram.streamMode} 在私有话题聊天中使用 \texttt{sendMessageDraft}（Bot API 9.3+）。这与渠道分块流式传输是分开的。
  \item \textbf{测试：} grammY mock 覆盖了私信 + 群组 mention 网关控制和出站发送；欢迎添加更多媒体/webhook 测试用例。
\end{itemize}


待解决问题

\begin{itemize}
  \item 如果遇到 Bot API 429 错误，考虑使用可选的 grammY 插件（throttler）。
  \item 添加更多结构化媒体测试（贴纸、语音消息）。
  \item 使 webhook 监听端口可配置（目前固定为 8787，除非通过 Gateway 网关配置）。
\end{itemize}



\clearpage

% === 源文件: channels/group-messages.md ===
\section{群组消息（WhatsApp 网页渠道）}


目标：让 Clawd 留在 WhatsApp 群组中，仅在被提及时唤醒，并将该对话线程与个人私信会话分开。

注意：\texttt{agents.list[].groupChat.mentionPatterns} 现在也被 Telegram/Discord/Slack/iMessage 使用；本文档重点介绍 WhatsApp 特定的行为。对于多智能体设置，为每个智能体设置 \texttt{agents.list[].groupChat.mentionPatterns}（或使用 \texttt{messages.groupChat.mentionPatterns} 作为全局回退）。

\subsection{已实现的功能（2025-12-03）}


\begin{itemize}
  \item 激活模式：\texttt{mention}（默认）或 \texttt{always}。\texttt{mention} 需要被提及（通过 \texttt{mentionedJids} 的真实 WhatsApp @提及、正则表达式模式，或文本中任意位置的机器人 E.164 号码）。\texttt{always} 会在每条消息时唤醒智能体，但它应该只在能提供有意义价值时才回复；否则返回静默令牌 \texttt{NO\_REPLY}。默认值可在配置中设置（\texttt{channels.whatsapp.groups}），并可通过 \texttt{/activation} 为每个群组单独覆盖。当设置了 \texttt{channels.whatsapp.groups} 时，它同时充当群组允许列表（包含 \texttt{"*"} 以允许所有群组）。
  \item 群组策略：\texttt{channels.whatsapp.groupPolicy} 控制是否接受群组消息（\texttt{open|disabled|allowlist}）。\texttt{allowlist} 使用 \texttt{channels.whatsapp.groupAllowFrom}（回退：显式的 \texttt{channels.whatsapp.allowFrom}）。默认为 \texttt{allowlist}（在你添加发送者之前被阻止）。
  \item 独立群组会话：会话键格式为 \texttt{agent::whatsapp:group:}，因此 \texttt{/verbose on} 或 \texttt{/think high}（作为独立消息发送）等命令仅作用于该群组；个人私信状态不受影响。群组线程会跳过心跳。
  \item 上下文注入：\textbf{仅待处理}的群组消息（默认 50 条），即\textit{未}触发运行的消息，会以 \texttt{[Chat messages since your last reply - for context]} 为前缀注入，触发行在 \texttt{[Current message - respond to this]} 下。已在会话中的消息不会重复注入。
  \item 发送者显示：每个群组批次现在以 \texttt{[from: Sender Name (+E164)]} 结尾，让 Pi 知道是谁在说话。
  \item 阅后即焚/一次性查看：我们在提取文本/提及之前会先解包这些消息，因此其中的提及仍会触发。
  \item 群组系统提示：在群组会话的第一轮（以及每当 \texttt{/activation} 更改模式时），我们会向系统提示注入一段简短说明，如 \texttt{You are replying inside the WhatsApp group "". Group members: Alice (+44...), Bob (+43...), … Activation: trigger-only … Address the specific sender noted in the message context.} 如果元数据不可用，我们仍会告知智能体这是一个群聊。
\end{itemize}


\subsection{配置示例（WhatsApp）}


在 \texttt{\textasciitilde{}/.openclaw/openclaw.json} 中添加 \texttt{groupChat} 块，以便在 WhatsApp 剥离文本正文中的可视 \texttt{@} 时，显示名称提及仍能正常工作：

\begin{lstlisting}[language=json]
{
  channels: {
    whatsapp: {
      groups: {
        "*": { requireMention: true },
      },
    },
  },
  agents: {
    list: [
      {
        id: "main",
        groupChat: {
          historyLimit: 50,
          mentionPatterns: ["@?openclaw", "\\+?15555550123"],
        },
      },
    ],
  },
}
\end{lstlisting}


注意：

\begin{itemize}
  \item 正则表达式不区分大小写；它们涵盖了像 \texttt{@openclaw} 这样的显示名称提及，以及带或不带 \texttt{+}/空格的原始号码。
  \item 当有人点击联系人时，WhatsApp 仍会通过 \texttt{mentionedJids} 发送规范的提及，因此号码回退很少需要，但作为安全网很有用。
\end{itemize}


\subsubsection{激活命令（仅所有者）}


使用群聊命令：

\begin{itemize}
  \item \texttt{/activation mention}
  \item \texttt{/activation always}
\end{itemize}


只有所有者号码（来自 \texttt{channels.whatsapp.allowFrom}，或未设置时使用机器人自己的 E.164）可以更改此设置。在群组中发送 \texttt{/status} 作为独立消息以查看当前激活模式。

\subsection{使用方法}


\begin{enumerate}
  \item 将你的 WhatsApp 账号（运行 OpenClaw 的账号）添加到群组。
  \item 说 \texttt{@openclaw …}（或包含号码）。只有允许列表中的发送者才能触发，除非你设置 \texttt{groupPolicy: "open"}。
  \item 智能体提示将包含最近的群组上下文以及尾部的 \texttt{[from: …]} 标记，以便它能够回应正确的人。
  \item 会话级指令（\texttt{/verbose on}、\texttt{/think high}、\texttt{/new} 或 \texttt{/reset}、\texttt{/compact}）仅适用于该群组的会话；将它们作为独立消息发送以使其生效。你的个人私信会话保持独立。
\end{enumerate}


\subsection{测试/验证}


\begin{itemize}
  \item 手动冒烟测试：
  \item 在群组中发送 \texttt{@openclaw} 提及，确认收到引用发送者名称的回复。
  \item 发送第二次提及，验证历史记录块被包含，然后在下一轮清除。
  \item 检查 Gateway 网关日志（使用 \texttt{--verbose} 运行）以查看 \texttt{inbound web message} 条目，显示 \texttt{from: } 和 \texttt{[from: …]} 后缀。
\end{itemize}


\subsection{已知注意事项}


\begin{itemize}
  \item 群组有意跳过心跳以避免嘈杂的广播。
  \item 回声抑制使用组合的批次字符串；如果你发送两次相同的文本但没有提及，只有第一次会得到响应。
  \item 会话存储条目将在会话存储中显示为 \texttt{agent::whatsapp:group:}（默认为 \texttt{\textasciitilde{}/.openclaw/agents//sessions/sessions.json}）；缺失条目只是意味着该群组尚未触发运行。
  \item 群组中的输入指示器遵循 \texttt{agents.defaults.typingMode}（默认：未被提及时为 \texttt{message}）。
\end{itemize}



\clearpage

% === 源文件: channels/groups.md ===
\section{群组}


OpenClaw 在各平台上统一处理群聊：WhatsApp、Telegram、Discord、Slack、Signal、iMessage、Microsoft Teams。

\subsection{新手入门（2 分钟）}


OpenClaw"运行"在你自己的消息账户上。没有单独的 WhatsApp 机器人用户。如果\textbf{你}在一个群组中，OpenClaw 就可以看到该群组并在其中回复。

默认行为：

\begin{itemize}
  \item 群组受限（\texttt{groupPolicy: "allowlist"}）。
  \item 除非你明确禁用提及限制，否则回复需要 @ 提及。
\end{itemize}


解释：允许列表中的发送者可以通过提及来触发 OpenClaw。

\begin{quote}
简而言之

- \textbf{私信访问}由 \texttt{*.allowFrom} 控制。
- \textbf{群组访问}由 \texttt{\textit{.groupPolicy} + 允许列表（\texttt{}.groups}、\texttt{*.groupAllowFrom}）控制。
- \textbf{回复触发}由提及限制（\texttt{requireMention}、\texttt{/activation}）控制。
\end{quote}


快速流程（群消息会发生什么）：

\begin{lstlisting}
groupPolicy? disabled -> 丢弃
groupPolicy? allowlist -> 群组允许? 否 -> 丢弃
requireMention? 是 -> 被提及? 否 -> 仅存储为上下文
否则 -> 回复
\end{lstlisting}


% [图片: 群消息流程]

如果你想...

\begin{table}[h]
\centering
\small
\begin{tabular}{|l|l|}
\hline
目标 & 设置什么 \\
\hline
允许所有群组但仅在 @ 提及时回复 & \texttt{groups: \{ "*": \{ requireMention: true \} \}} \\
\hline
禁用所有群组回复 & \texttt{groupPolicy: "disabled"} \\
\hline
仅特定群组 & \texttt{groups: \{ "": \{ ... \} \}}（无 \texttt{"*"} 键） \\
\hline
仅你可以在群组中触发 & \texttt{groupPolicy: "allowlist"}、\texttt{groupAllowFrom: ["+1555..."]} \\
\hline
\end{tabular}
\end{table}



\subsection{会话键}


\begin{itemize}
  \item 群组会话使用 \texttt{agent:::group:} 会话键（房间/频道使用 \texttt{agent:::channel:}）。
  \item Telegram 论坛话题在群组 ID 后添加 \texttt{:topic:}，因此每个话题都有自己的会话。
  \item 私聊使用主会话（或按发送者配置时使用各自的会话）。
  \item 群组会话跳过心跳。
\end{itemize}


\subsection{模式：个人私信 + 公开群组（单智能体）}


是的——如果你的"个人"流量是\textbf{私信}而"公开"流量是\textbf{群组}，这种方式效果很好。

原因：在单智能体模式下，私信通常落在\textbf{主}会话键（\texttt{agent:main:main}）中，而群组始终使用\textbf{非主}会话键（\texttt{agent:main::group:}）。如果你启用 \texttt{mode: "non-main"} 的沙箱隔离，这些群组会话在 Docker 中运行，而你的主私信会话保持在主机上。

这给你一个智能体"大脑"（共享工作区 + 记忆），但两种执行姿态：

\begin{itemize}
  \item \textbf{私信}：完整工具（主机）
  \item \textbf{群组}：沙箱 + 受限工具（Docker）
\end{itemize}


\begin{quote}
如果你需要真正独立的工作区/角色（"个人"和"公开"绝不能混合），请使用第二个智能体 + 绑定。参见\href{/concepts/multi-agent}{多智能体路由}。
\end{quote}


示例（私信在主机上，群组沙箱隔离 + 仅消息工具）：

\begin{lstlisting}[language=json]
{
  agents: {
    defaults: {
      sandbox: {
        mode: "non-main", // 群组/频道是非主 -> 沙箱隔离
        scope: "session", // 最强隔离（每个群组/频道一个容器）
        workspaceAccess: "none",
      },
    },
  },
  tools: {
    sandbox: {
      tools: {
        // 如果 allow 非空，其他所有工具都被阻止（deny 仍然优先）。
        allow: ["group:messaging", "group:sessions"],
        deny: ["group:runtime", "group:fs", "group:ui", "nodes", "cron", "gateway"],
      },
    },
  },
}
\end{lstlisting}


想要"群组只能看到文件夹 X"而不是"无主机访问"？保持 \texttt{workspaceAccess: "none"} 并仅将允许的路径挂载到沙箱中：

\begin{lstlisting}[language=json]
{
  agents: {
    defaults: {
      sandbox: {
        mode: "non-main",
        scope: "session",
        workspaceAccess: "none",
        docker: {
          binds: [
            // hostPath:containerPath:mode
            "~/FriendsShared:/data:ro",
          ],
        },
      },
    },
  },
}
\end{lstlisting}


相关：

\begin{itemize}
  \item 配置键和默认值：\href{/gateway/configuration\#agentsdefaultssandbox}{Gateway 网关配置}
  \item 调试为什么工具被阻止：\href{/gateway/sandbox-vs-tool-policy-vs-elevated}{沙箱 vs 工具策略 vs 提权}
  \item 绑定挂载详情：\href{/gateway/sandboxing\#custom-bind-mounts}{沙箱隔离}
\end{itemize}


\subsection{显示标签}


\begin{itemize}
  \item UI 标签在可用时使用 \texttt{displayName}，格式为 \texttt{:}。
  \item \texttt{\#room} 保留用于房间/频道；群聊使用 \texttt{g-}（小写，空格 -> \texttt{-}，保留 \texttt{\#@+.\_-}）。
\end{itemize}


\subsection{群组策略}


控制每个渠道如何处理群组/房间消息：

\begin{lstlisting}[language=json]
{
  channels: {
    whatsapp: {
      groupPolicy: "disabled", // "open" | "disabled" | "allowlist"
      groupAllowFrom: ["+15551234567"],
    },
    telegram: {
      groupPolicy: "disabled",
      groupAllowFrom: ["123456789", "@username"],
    },
    signal: {
      groupPolicy: "disabled",
      groupAllowFrom: ["+15551234567"],
    },
    imessage: {
      groupPolicy: "disabled",
      groupAllowFrom: ["chat_id:123"],
    },
    msteams: {
      groupPolicy: "disabled",
      groupAllowFrom: ["user@org.com"],
    },
    discord: {
      groupPolicy: "allowlist",
      guilds: {
        GUILD_ID: { channels: { help: { allow: true } } },
      },
    },
    slack: {
      groupPolicy: "allowlist",
      channels: { "#general": { allow: true } },
    },
    matrix: {
      groupPolicy: "allowlist",
      groupAllowFrom: ["@owner:example.org"],
      groups: {
        "!roomId:example.org": { allow: true },
        "#alias:example.org": { allow: true },
      },
    },
  },
}
\end{lstlisting}



\begin{table}[h]
\centering
\small
\begin{tabular}{|l|l|}
\hline
策略 & 行为 \\
\hline
\texttt{"open"} & 群组绕过允许列表；提及限制仍然适用。 \\
\hline
\texttt{"disabled"} & 完全阻止所有群组消息。 \\
\hline
\texttt{"allowlist"} & 仅允许与配置的允许列表匹配的群组/房间。 \\
\hline
\end{tabular}
\end{table}



注意事项：

\begin{itemize}
  \item \texttt{groupPolicy} 与提及限制（需要 @ 提及）是分开的。
  \item WhatsApp/Telegram/Signal/iMessage/Microsoft Teams：使用 \texttt{groupAllowFrom}（回退：显式 \texttt{allowFrom}）。
  \item Discord：允许列表使用 \texttt{channels.discord.guilds..channels}。
  \item Slack：允许列表使用 \texttt{channels.slack.channels}。
  \item Matrix：允许列表使用 \texttt{channels.matrix.groups}（房间 ID、别名或名称）。使用 \texttt{channels.matrix.groupAllowFrom} 限制发送者；也支持每个房间的 \texttt{users} 允许列表。
  \item 群组私信单独控制（\texttt{channels.discord.dm.\textit{}、\texttt{channels.slack.dm.}}）。
  \item Telegram 允许列表可以匹配用户 ID（\texttt{"123456789"}、\texttt{"telegram:123456789"}、\texttt{"tg:123456789"}）或用户名（\texttt{"@alice"} 或 \texttt{"alice"}）；前缀不区分大小写。
  \item 默认为 \texttt{groupPolicy: "allowlist"}；如果你的群组允许列表为空，群组消息将被阻止。
\end{itemize}


快速心智模型（群组消息的评估顺序）：

\begin{enumerate}
  \item \texttt{groupPolicy}（open/disabled/allowlist）
  \item 群组允许列表（\texttt{\textit{.groups}、\texttt{}.groupAllowFrom}、渠道特定允许列表）
  \item 提及限制（\texttt{requireMention}、\texttt{/activation}）
\end{enumerate}


\subsection{提及限制（默认）}


群组消息需要提及，除非按群组覆盖。默认值位于 \texttt{\textit{.groups."}"} 下的每个子系统中。

回复机器人消息被视为隐式提及（当渠道支持回复元数据时）。这适用于 Telegram、WhatsApp、Slack、Discord 和 Microsoft Teams。

\begin{lstlisting}[language=json]
{
  channels: {
    whatsapp: {
      groups: {
        "*": { requireMention: true },
        "123@g.us": { requireMention: false },
      },
    },
    telegram: {
      groups: {
        "*": { requireMention: true },
        "123456789": { requireMention: false },
      },
    },
    imessage: {
      groups: {
        "*": { requireMention: true },
        "123": { requireMention: false },
      },
    },
  },
  agents: {
    list: [
      {
        id: "main",
        groupChat: {
          mentionPatterns: ["@openclaw", "openclaw", "\\+15555550123"],
          historyLimit: 50,
        },
      },
    ],
  },
}
\end{lstlisting}


注意事项：

\begin{itemize}
  \item \texttt{mentionPatterns} 是不区分大小写的正则表达式。
  \item 提供显式提及的平台仍然通过；模式是回退。
  \item 每个智能体覆盖：\texttt{agents.list[].groupChat.mentionPatterns}（当多个智能体共享一个群组时有用）。
  \item 提及限制仅在提及检测可行时强制执行（原生提及或 \texttt{mentionPatterns} 已配置）。
  \item Discord 默认值位于 \texttt{channels.discord.guilds."*"}（可按服务器/频道覆盖）。
  \item 群组历史上下文在渠道间统一包装，并且是\textbf{仅待处理}（由于提及限制而跳过的消息）；使用 \texttt{messages.groupChat.historyLimit} 作为全局默认值，使用 \texttt{channels..historyLimit}（或 \texttt{channels..accounts.*.historyLimit}）进行覆盖。设置 \texttt{0} 以禁用。
\end{itemize}


\subsection{群组/频道工具限制（可选）}


某些渠道配置支持限制\textbf{特定群组/房间/频道内}可用的工具。

\begin{itemize}
  \item \texttt{tools}：为整个群组允许/拒绝工具。
  \item \texttt{toolsBySender}：群组内的按发送者覆盖（键是发送者 ID/用户名/邮箱/电话号码，取决于渠道）。使用 \texttt{"*"} 作为通配符。
\end{itemize}


解析顺序（最具体的优先）：

\begin{enumerate}
  \item 群组/频道 \texttt{toolsBySender} 匹配
  \item 群组/频道 \texttt{tools}
  \item 默认（\texttt{"*"}）\texttt{toolsBySender} 匹配
  \item 默认（\texttt{"*"}）\texttt{tools}
\end{enumerate}


示例（Telegram）：

\begin{lstlisting}[language=json]
{
  channels: {
    telegram: {
      groups: {
        "*": { tools: { deny: ["exec"] } },
        "-1001234567890": {
          tools: { deny: ["exec", "read", "write"] },
          toolsBySender: {
            "123456789": { alsoAllow: ["exec"] },
          },
        },
      },
    },
  },
}
\end{lstlisting}


注意事项：

\begin{itemize}
  \item 群组/频道工具限制在全局/智能体工具策略之外额外应用（deny 仍然优先）。
  \item 某些渠道对房间/频道使用不同的嵌套结构（例如，Discord \texttt{guilds.\textit{.channels.}}、Slack \texttt{channels.\textit{}、MS Teams \texttt{teams.}.channels.*}）。
\end{itemize}


\subsection{群组允许列表}


当配置了 \texttt{channels.whatsapp.groups}、\texttt{channels.telegram.groups} 或 \texttt{channels.imessage.groups} 时，键作为群组允许列表。使用 \texttt{"*"} 允许所有群组，同时仍设置默认提及行为。

常见意图（复制/粘贴）：

\begin{enumerate}
  \item 禁用所有群组回复
\end{enumerate}


\begin{lstlisting}[language=json]
{
  channels: { whatsapp: { groupPolicy: "disabled" } },
}
\end{lstlisting}


\begin{enumerate}
  \item 仅允许特定群组（WhatsApp）
\end{enumerate}


\begin{lstlisting}[language=json]
{
  channels: {
    whatsapp: {
      groups: {
        "123@g.us": { requireMention: true },
        "456@g.us": { requireMention: false },
      },
    },
  },
}
\end{lstlisting}


\begin{enumerate}
  \item 允许所有群组但需要提及（显式）
\end{enumerate}


\begin{lstlisting}[language=json]
{
  channels: {
    whatsapp: {
      groups: { "*": { requireMention: true } },
    },
  },
}
\end{lstlisting}


\begin{enumerate}
  \item 仅所有者可以在群组中触发（WhatsApp）
\end{enumerate}


\begin{lstlisting}[language=json]
{
  channels: {
    whatsapp: {
      groupPolicy: "allowlist",
      groupAllowFrom: ["+15551234567"],
      groups: { "*": { requireMention: true } },
    },
  },
}
\end{lstlisting}


\subsection{激活（仅所有者）}


群组所有者可以切换每个群组的激活状态：

\begin{itemize}
  \item \texttt{/activation mention}
  \item \texttt{/activation always}
\end{itemize}


所有者由 \texttt{channels.whatsapp.allowFrom} 确定（未设置时为机器人自身的 E.164）。将命令作为独立消息发送。其他平台目前忽略 \texttt{/activation}。

\subsection{上下文字段}


群组入站负载设置：

\begin{itemize}
  \item \texttt{ChatType=group}
  \item \texttt{GroupSubject}（如果已知）
  \item \texttt{GroupMembers}（如果已知）
  \item \texttt{WasMentioned}（提及限制结果）
  \item Telegram 论坛话题还包括 \texttt{MessageThreadId} 和 \texttt{IsForum}。
\end{itemize}


智能体系统提示在新群组会话的第一轮包含群组介绍。它提醒模型像人类一样回复，避免 Markdown 表格，避免输入字面量 \texttt{\textbackslash\{\}n} 序列。

\subsection{iMessage 特定内容}


\begin{itemize}
  \item 路由或允许列表时优先使用 \texttt{chat\_id:}。
  \item 列出聊天：\texttt{imsg chats --limit 20}。
  \item 群组回复始终返回到相同的 \texttt{chat\_id}。
\end{itemize}


\subsection{WhatsApp 特定内容}


参见\href{/channels/group-messages}{群消息}了解 WhatsApp 专有行为（历史注入、提及处理详情）。


\clearpage

% === 源文件: channels/imessage.md ===
\section{iMessage (imsg)}


状态：外部 CLI 集成。Gateway 网关生成 \texttt{imsg rpc}（基于 stdio 的 JSON-RPC）。

\subsection{快速设置（新手）}


\begin{enumerate}
  \item 确保在此 Mac 上已登录"信息"。
  \item 安装 \texttt{imsg}：
\end{enumerate}

\begin{itemize}
  \item \texttt{brew install steipete/tap/imsg}
\end{itemize}

\begin{enumerate}
  \item 配置 OpenClaw 的 \texttt{channels.imessage.cliPath} 和 \texttt{channels.imessage.dbPath}。
  \item 启动 Gateway 网关并批准所有 macOS 提示（自动化 + 完全磁盘访问权限）。
\end{enumerate}


最小配置：

\begin{lstlisting}[language=json]
{
  channels: {
    imessage: {
      enabled: true,
      cliPath: "/usr/local/bin/imsg",
      dbPath: "/Users/<you>/Library/Messages/chat.db",
    },
  },
}
\end{lstlisting}


\subsection{简介}


\begin{itemize}
  \item 基于 macOS 上 \texttt{imsg} 的 iMessage 渠道。
  \item 确定性路由：回复始终返回到 iMessage。
  \item 私信共享智能体的主会话；群组是隔离的（\texttt{agent::imessage:group:}）。
  \item 如果多参与者会话以 \texttt{is\_group=false} 到达，你仍可使用 \texttt{channels.imessage.groups} 按 \texttt{chat\_id} 隔离（参见下方"类群组会话"）。
\end{itemize}


\subsection{配置写入}


默认情况下，iMessage 允许写入由 \texttt{/config set|unset} 触发的配置更新（需要 \texttt{commands.config: true}）。

禁用方式：

\begin{lstlisting}[language=json]
{
  channels: { imessage: { configWrites: false } },
}
\end{lstlisting}


\subsection{要求}


\begin{itemize}
  \item 已登录"信息"的 macOS。
  \item OpenClaw + \texttt{imsg} 的完全磁盘访问权限（访问"信息"数据库）。
  \item 发送时需要自动化权限。
  \item \texttt{channels.imessage.cliPath} 可以指向任何代理 stdin/stdout 的命令（例如，通过 SSH 连接到另一台 Mac 并运行 \texttt{imsg rpc} 的包装脚本）。
\end{itemize}


\subsection{设置（快速路径）}


\begin{enumerate}
  \item 确保在此 Mac 上已登录"信息"。
  \item 配置 iMessage 并启动 Gateway 网关。
\end{enumerate}


\subsubsection{专用机器人 macOS 用户（用于隔离身份）}


如果你希望机器人从\textbf{独立的 iMessage 身份}发送（并保持你的个人"信息"整洁），请使用专用 Apple ID + 专用 macOS 用户。

\begin{enumerate}
  \item 创建专用 Apple ID（例如：\texttt{my-cool-bot@icloud.com}）。
\end{enumerate}

\begin{itemize}
  \item Apple 可能需要电话号码进行验证 / 2FA。
\end{itemize}

\begin{enumerate}
  \item 创建 macOS 用户（例如：\texttt{openclawhome}）并登录。
  \item 在该 macOS 用户中打开"信息"并使用机器人 Apple ID 登录 iMessage。
  \item 启用远程登录（系统设置 → 通用 → 共享 → 远程登录）。
  \item 安装 \texttt{imsg}：
\end{enumerate}

\begin{itemize}
  \item \texttt{brew install steipete/tap/imsg}
\end{itemize}

\begin{enumerate}
  \item 设置 SSH 使 \texttt{ssh @localhost true} 无需密码即可工作。
  \item 将 \texttt{channels.imessage.accounts.bot.cliPath} 指向以机器人用户身份运行 \texttt{imsg} 的 SSH 包装脚本。
\end{enumerate}


首次运行注意事项：发送/接收可能需要在\textit{机器人 macOS 用户}中进行 GUI 批准（自动化 + 完全磁盘访问权限）。如果 \texttt{imsg rpc} 看起来卡住或退出，请登录该用户（屏幕共享很有帮助），运行一次 \texttt{imsg chats --limit 1} / \texttt{imsg send ...}，批准提示，然后重试。

示例包装脚本（\texttt{chmod +x}）。将 `` 替换为你的实际 macOS 用户名：

\begin{lstlisting}[language=bash]
#!/usr/bin/env bash
set -euo pipefail

# Run an interactive SSH once first to accept host keys:
#   ssh <bot-macos-user>@localhost true
exec /usr/bin/ssh -o BatchMode=yes -o ConnectTimeout=5 -T <bot-macos-user>@localhost \
  "/usr/local/bin/imsg" "$@"
\end{lstlisting}


示例配置：

\begin{lstlisting}[language=json]
{
  channels: {
    imessage: {
      enabled: true,
      accounts: {
        bot: {
          name: "Bot",
          enabled: true,
          cliPath: "/path/to/imsg-bot",
          dbPath: "/Users/<bot-macos-user>/Library/Messages/chat.db",
        },
      },
    },
  },
}
\end{lstlisting}


对于单账户设置，使用扁平选项（\texttt{channels.imessage.cliPath}、\texttt{channels.imessage.dbPath}）而不是 \texttt{accounts} 映射。

\subsubsection{远程/SSH 变体（可选）}


如果你想在另一台 Mac 上使用 iMessage，请将 \texttt{channels.imessage.cliPath} 设置为通过 SSH 在远程 macOS 主机上运行 \texttt{imsg} 的包装脚本。OpenClaw 只需要 stdio。

示例包装脚本：

\begin{lstlisting}[language=bash]
#!/usr/bin/env bash
exec ssh -T gateway-host imsg "$@"
\end{lstlisting}


\textbf{远程附件：} 当 \texttt{cliPath} 通过 SSH 指向远程主机时，"信息"数据库中的附件路径引用的是远程机器上的文件。OpenClaw 可以通过设置 \texttt{channels.imessage.remoteHost} 自动通过 SCP 获取这些文件：

\begin{lstlisting}[language=json]
{
  channels: {
    imessage: {
      cliPath: "~/imsg-ssh", // SSH wrapper to remote Mac
      remoteHost: "user@gateway-host", // for SCP file transfer
      includeAttachments: true,
    },
  },
}
\end{lstlisting}


如果未设置 \texttt{remoteHost}，OpenClaw 会尝试通过解析包装脚本中的 SSH 命令自动检测。建议显式配置以提高可靠性。

\paragraph{通过 Tailscale 连接远程 Mac（示例）}


如果 Gateway 网关运行在 Linux 主机/虚拟机上但 iMessage 必须运行在 Mac 上，Tailscale 是最简单的桥接方式：Gateway 网关通过 tailnet 与 Mac 通信，通过 SSH 运行 \texttt{imsg}，并通过 SCP 获取附件。

架构：

\begin{lstlisting}
┌──────────────────────────────┐          SSH (imsg rpc)          ┌──────────────────────────┐
│ Gateway host (Linux/VM)      │──────────────────────────────────▶│ Mac with Messages + imsg │
│ - openclaw gateway           │          SCP (attachments)        │ - Messages signed in     │
│ - channels.imessage.cliPath  │◀──────────────────────────────────│ - Remote Login enabled   │
└──────────────────────────────┘                                   └──────────────────────────┘
              ▲
              │ Tailscale tailnet (hostname or 100.x.y.z)
              ▼
        user@gateway-host
\end{lstlisting}


具体配置示例（Tailscale 主机名）：

\begin{lstlisting}[language=json]
{
  channels: {
    imessage: {
      enabled: true,
      cliPath: "~/.openclaw/scripts/imsg-ssh",
      remoteHost: "bot@mac-mini.tailnet-1234.ts.net",
      includeAttachments: true,
      dbPath: "/Users/bot/Library/Messages/chat.db",
    },
  },
}
\end{lstlisting}


示例包装脚本（\texttt{\textasciitilde{}/.openclaw/scripts/imsg-ssh}）：

\begin{lstlisting}[language=bash]
#!/usr/bin/env bash
exec ssh -T bot@mac-mini.tailnet-1234.ts.net imsg "$@"
\end{lstlisting}


注意事项：

\begin{itemize}
  \item 确保 Mac 已登录"信息"，并已启用远程登录。
  \item 使用 SSH 密钥使 \texttt{ssh bot@mac-mini.tailnet-1234.ts.net} 无需提示即可工作。
  \item \texttt{remoteHost} 应与 SSH 目标匹配，以便 SCP 可以获取附件。
\end{itemize}


多账户支持：使用 \texttt{channels.imessage.accounts} 配置每个账户及可选的 \texttt{name}。参见 \href{/gateway/configuration\#telegramaccounts--discordaccounts--slackaccounts--signalaccounts--imessageaccounts}{\texttt{gateway/configuration}} 了解共享模式。不要提交 \texttt{\textasciitilde{}/.openclaw/openclaw.json}（它通常包含令牌）。

\subsection{访问控制（私信 + 群组）}


私信：

\begin{itemize}
  \item 默认：\texttt{channels.imessage.dmPolicy = "pairing"}。
  \item 未知发送者会收到配对码；消息在批准前会被忽略（配对码在 1 小时后过期）。
  \item 批准方式：
  \item \texttt{openclaw pairing list imessage}
  \item \texttt{openclaw pairing approve imessage }
  \item 配对是 iMessage 私信的默认令牌交换方式。详情：\href{/channels/pairing}{配对}
\end{itemize}


群组：

\begin{itemize}
  \item \texttt{channels.imessage.groupPolicy = open | allowlist | disabled}。
  \item 设置 \texttt{allowlist} 时，\texttt{channels.imessage.groupAllowFrom} 控制谁可以在群组中触发。
  \item 提及检测使用 \texttt{agents.list[].groupChat.mentionPatterns}（或 \texttt{messages.groupChat.mentionPatterns}），因为 iMessage 没有原生提及元数据。
  \item 多智能体覆盖：在 \texttt{agents.list[].groupChat.mentionPatterns} 上设置每个智能体的模式。
\end{itemize}


\subsection{工作原理（行为）}


\begin{itemize}
  \item \texttt{imsg} 流式传输消息事件；Gateway 网关将它们规范化为共享渠道信封。
  \item 回复始终路由回相同的 chat id 或 handle。
\end{itemize}


\subsection{类群组会话（is\_group=false）}


某些 iMessage 会话可能有多个参与者，但根据"信息"存储聊天标识符的方式，仍以 \texttt{is\_group=false} 到达。

如果你在 \texttt{channels.imessage.groups} 下显式配置了 \texttt{chat\_id}，OpenClaw 会将该会话视为"群组"用于：

\begin{itemize}
  \item 会话隔离（独立的 \texttt{agent::imessage:group:} 会话键）
  \item 群组允许列表 / 提及检测行为
\end{itemize}


示例：

\begin{lstlisting}[language=json]
{
  channels: {
    imessage: {
      groupPolicy: "allowlist",
      groupAllowFrom: ["+15555550123"],
      groups: {
        "42": { requireMention: false },
      },
    },
  },
}
\end{lstlisting}


当你想为特定会话使用隔离的个性/模型时这很有用（参见\href{/concepts/multi-agent}{多智能体路由}）。关于文件系统隔离，参见\href{/gateway/sandboxing}{沙箱隔离}。

\subsection{媒体 + 限制}


\begin{itemize}
  \item 通过 \texttt{channels.imessage.includeAttachments} 可选附件摄取。
  \item 通过 \texttt{channels.imessage.mediaMaxMb} 设置媒体上限。
\end{itemize}


\subsection{限制}


\begin{itemize}
  \item 出站文本按 \texttt{channels.imessage.textChunkLimit} 分块（默认 4000）。
  \item 可选换行分块：设置 \texttt{channels.imessage.chunkMode="newline"} 在长度分块前按空行（段落边界）分割。
  \item 媒体上传受 \texttt{channels.imessage.mediaMaxMb} 限制（默认 16）。
\end{itemize}


\subsection{寻址 / 投递目标}


优先使用 \texttt{chat\_id} 进行稳定路由：

\begin{itemize}
  \item \texttt{chat\_id:123}（推荐）
  \item \texttt{chat\_guid:...}
  \item \texttt{chat\_identifier:...}
  \item 直接 handle：\texttt{imessage:+1555} / \texttt{sms:+1555} / \texttt{user@example.com}
\end{itemize}


列出聊天：

\begin{lstlisting}
imsg chats --limit 20
\end{lstlisting}


\subsection{配置参考（iMessage）}


完整配置：\href{/gateway/configuration}{配置}

提供商选项：

\begin{itemize}
  \item \texttt{channels.imessage.enabled}：启用/禁用渠道启动。
  \item \texttt{channels.imessage.cliPath}：\texttt{imsg} 路径。
  \item \texttt{channels.imessage.dbPath}："信息"数据库路径。
  \item \texttt{channels.imessage.remoteHost}：当 \texttt{cliPath} 指向远程 Mac 时用于 SCP 附件传输的 SSH 主机（例如 \texttt{user@gateway-host}）。如未设置则从 SSH 包装脚本自动检测。
  \item \texttt{channels.imessage.service}：\texttt{imessage | sms | auto}。
  \item \texttt{channels.imessage.region}：短信区域。
  \item \texttt{channels.imessage.dmPolicy}：\texttt{pairing | allowlist | open | disabled}（默认：pairing）。
  \item \texttt{channels.imessage.allowFrom}：私信允许列表（handle、邮箱、E.164 号码或 \texttt{chat\_id:\textit{}）。\texttt{open} 需要 \texttt{"}"}。iMessage 没有用户名；使用 handle 或聊天目标。
  \item \texttt{channels.imessage.groupPolicy}：\texttt{open | allowlist | disabled}（默认：allowlist）。
  \item \texttt{channels.imessage.groupAllowFrom}：群组发送者允许列表。
  \item \texttt{channels.imessage.historyLimit} / \texttt{channels.imessage.accounts.*.historyLimit}：作为上下文包含的最大群组消息数（0 禁用）。
  \item \texttt{channels.imessage.dmHistoryLimit}：私信历史限制（用户轮次）。每用户覆盖：\texttt{channels.imessage.dms[""].historyLimit}。
  \item \texttt{channels.imessage.groups}：每群组默认值 + 允许列表（使用 \texttt{"*"} 作为全局默认值）。
  \item \texttt{channels.imessage.includeAttachments}：将附件摄取到上下文。
  \item \texttt{channels.imessage.mediaMaxMb}：入站/出站媒体上限（MB）。
  \item \texttt{channels.imessage.textChunkLimit}：出站分块大小（字符）。
  \item \texttt{channels.imessage.chunkMode}：\texttt{length}（默认）或 \texttt{newline} 在长度分块前按空行（段落边界）分割。
\end{itemize}


相关全局选项：

\begin{itemize}
  \item \texttt{agents.list[].groupChat.mentionPatterns}（或 \texttt{messages.groupChat.mentionPatterns}）。
  \item \texttt{messages.responsePrefix}。
\end{itemize}



\clearpage

% === 源文件: channels/index.md ===
\section{聊天渠道}


OpenClaw 可以在你已经使用的任何聊天应用上与你交流。每个渠道通过 Gateway 网关连接。
所有渠道都支持文本；媒体和表情回应的支持因渠道而异。

\subsection{支持的渠道}


\begin{itemize}
  \item \href{/channels/bluebubbles}{BlueBubbles} — \textbf{推荐用于 iMessage}；使用 BlueBubbles macOS 服务器 REST API，功能完整（编辑、撤回、特效、回应、群组管理——编辑功能在 macOS 26 Tahoe 上目前不可用）。
  \item \href{/channels/discord}{Discord} — Discord Bot API + Gateway；支持服务器、频道和私信。
  \item \href{/channels/feishu}{飞书} — 飞书（Lark）机器人（插件，需单独安装）。
  \item \href{/channels/googlechat}{Google Chat} — 通过 HTTP webhook 的 Google Chat API 应用。
  \item \href{/channels/imessage}{iMessage（旧版）} — 通过 imsg CLI 的旧版 macOS 集成（已弃用，新设置请使用 BlueBubbles）。
  \item \href{/channels/line}{LINE} — LINE Messaging API 机器人（插件，需单独安装）。
  \item \href{/channels/matrix}{Matrix} — Matrix 协议（插件，需单独安装）。
  \item \href{/channels/mattermost}{Mattermost} — Bot API + WebSocket；频道、群组、私信（插件，需单独安装）。
  \item \href{/channels/msteams}{Microsoft Teams} — Bot Framework；企业支持（插件，需单独安装）。
  \item \href{/channels/nextcloud-talk}{Nextcloud Talk} — 通过 Nextcloud Talk 的自托管聊天（插件，需单独安装）。
  \item \href{/channels/nostr}{Nostr} — 通过 NIP-04 的去中心化私信（插件，需单独安装）。
  \item \href{/channels/signal}{Signal} — signal-cli；注重隐私。
  \item \href{/channels/slack}{Slack} — Bolt SDK；工作区应用。
  \item \href{/channels/telegram}{Telegram} — 通过 grammY 使用 Bot API；支持群组。
  \item \href{/channels/tlon}{Tlon} — 基于 Urbit 的消息应用（插件，需单独安装）。
  \item \href{/channels/twitch}{Twitch} — 通过 IRC 连接的 Twitch 聊天（插件，需单独安装）。
  \item \href{/web/webchat}{WebChat} — 基于 WebSocket 的 Gateway 网关 WebChat 界面。
  \item \href{/channels/whatsapp}{WhatsApp} — 最受欢迎；使用 Baileys，需要二维码配对。
  \item \href{/channels/zalo}{Zalo} — Zalo Bot API；越南流行的消息应用（插件，需单独安装）。
  \item \href{/channels/zalouser}{Zalo Personal} — 通过二维码登录的 Zalo 个人账号（插件，需单独安装）。
\end{itemize}


\subsection{注意事项}


\begin{itemize}
  \item 渠道可以同时运行；配置多个渠道后，OpenClaw 会按聊天进行路由。
  \item 最快的设置方式通常是 \textbf{Telegram}（简单的机器人令牌）。WhatsApp 需要二维码配对，
\end{itemize}

并在磁盘上存储更多状态。
\begin{itemize}
  \item 群组行为因渠道而异；参见\href{/channels/groups}{群组}。
  \item 为安全起见，私信配对和允许列表会被强制执行；参见\href{/gateway/security}{安全}。
  \item Telegram 内部机制：\href{/channels/grammy}{grammY 说明}。
  \item 故障排除：\href{/channels/troubleshooting}{渠道故障排除}。
  \item 模型提供商单独记录；参见\href{/providers/models}{模型提供商}。
\end{itemize}



\clearpage

% === 源文件: channels/line.md ===
\section{LINE（插件）}


LINE 通过 LINE Messaging API 连接到 OpenClaw。该插件作为 webhook 接收器在 Gateway 网关上运行，使用你的 channel access token + channel secret 进行身份验证。

状态：通过插件支持。支持私信、群聊、媒体、位置、Flex 消息、模板消息和快捷回复。不支持表情回应和话题回复。

\subsection{需要安装插件}


安装 LINE 插件：

\begin{lstlisting}[language=bash]
openclaw plugins install @openclaw/line
\end{lstlisting}


本地检出（从 git 仓库运行时）：

\begin{lstlisting}[language=bash]
openclaw plugins install ./extensions/line
\end{lstlisting}


\subsection{配置步骤}


\begin{enumerate}
  \item 创建 LINE Developers 账户并打开控制台：
\end{enumerate}

https://developers.line.biz/console/
\begin{enumerate}
  \item 创建（或选择）一个 Provider 并添加 \textbf{Messaging API} 渠道。
  \item 从渠道设置中复制 \textbf{Channel access token} 和 \textbf{Channel secret}。
  \item 在 Messaging API 设置中启用 \textbf{Use webhook}。
  \item 将 webhook URL 设置为你的 Gateway 网关端点（必须使用 HTTPS）：
\end{enumerate}


\begin{lstlisting}
https://gateway-host/line/webhook
\end{lstlisting}


Gateway 网关会响应 LINE 的 webhook 验证（GET）和入站事件（POST）。如果你需要自定义路径，请设置 \texttt{channels.line.webhookPath} 或 \texttt{channels.line.accounts..webhookPath} 并相应更新 URL。

\subsection{配置}


最小配置：

\begin{lstlisting}[language=json]
{
  channels: {
    line: {
      enabled: true,
      channelAccessToken: "LINE_CHANNEL_ACCESS_TOKEN",
      channelSecret: "LINE_CHANNEL_SECRET",
      dmPolicy: "pairing",
    },
  },
}
\end{lstlisting}


环境变量（仅限默认账户）：

\begin{itemize}
  \item \texttt{LINE\_CHANNEL\_ACCESS\_TOKEN}
  \item \texttt{LINE\_CHANNEL\_SECRET}
\end{itemize}


Token/secret 文件：

\begin{lstlisting}[language=json]
{
  channels: {
    line: {
      tokenFile: "/path/to/line-token.txt",
      secretFile: "/path/to/line-secret.txt",
    },
  },
}
\end{lstlisting}


多账户配置：

\begin{lstlisting}[language=json]
{
  channels: {
    line: {
      accounts: {
        marketing: {
          channelAccessToken: "...",
          channelSecret: "...",
          webhookPath: "/line/marketing",
        },
      },
    },
  },
}
\end{lstlisting}


\subsection{访问控制}


私信默认使用配对模式。未知发送者会收到配对码，其消息在获得批准前会被忽略。

\begin{lstlisting}[language=bash]
openclaw pairing list line
openclaw pairing approve line <CODE>
\end{lstlisting}


允许列表和策略：

\begin{itemize}
  \item \texttt{channels.line.dmPolicy}：\texttt{pairing | allowlist | open | disabled}
  \item \texttt{channels.line.allowFrom}：私信的允许列表 LINE 用户 ID
  \item \texttt{channels.line.groupPolicy}：\texttt{allowlist | open | disabled}
  \item \texttt{channels.line.groupAllowFrom}：群组的允许列表 LINE 用户 ID
  \item 单群组覆盖：\texttt{channels.line.groups..allowFrom}
\end{itemize}


LINE ID 区分大小写。有效 ID 格式如下：

\begin{itemize}
  \item 用户：\texttt{U} + 32 位十六进制字符
  \item 群组：\texttt{C} + 32 位十六进制字符
  \item 房间：\texttt{R} + 32 位十六进制字符
\end{itemize}


\subsection{消息行为}


\begin{itemize}
  \item 文本按 5000 字符分块。
  \item Markdown 格式会被移除；代码块和表格会尽可能转换为 Flex 卡片。
  \item 流式响应会被缓冲；智能体处理时，LINE 会收到完整分块并显示加载动画。
  \item 媒体下载受 \texttt{channels.line.mediaMaxMb} 限制（默认 10）。
\end{itemize}


\subsection{渠道数据（富消息）}


使用 \texttt{channelData.line} 发送快捷回复、位置、Flex 卡片或模板消息。

\begin{lstlisting}[language=json]
{
  text: "Here you go",
  channelData: {
    line: {
      quickReplies: ["Status", "Help"],
      location: {
        title: "Office",
        address: "123 Main St",
        latitude: 35.681236,
        longitude: 139.767125,
      },
      flexMessage: {
        altText: "Status card",
        contents: {
          /* Flex payload */
        },
      },
      templateMessage: {
        type: "confirm",
        text: "Proceed?",
        confirmLabel: "Yes",
        confirmData: "yes",
        cancelLabel: "No",
        cancelData: "no",
      },
    },
  },
}
\end{lstlisting}


LINE 插件还提供 \texttt{/card} 命令用于 Flex 消息预设：

\begin{lstlisting}
/card info "Welcome" "Thanks for joining!"
\end{lstlisting}


\subsection{故障排除}


\begin{itemize}
  \item \textbf{Webhook 验证失败：} 确保 webhook URL 使用 HTTPS 且 \texttt{channelSecret} 与 LINE 控制台中的一致。
  \item \textbf{没有入站事件：} 确认 webhook 路径与 \texttt{channels.line.webhookPath} 匹配，且 Gateway 网关可从 LINE 访问。
  \item \textbf{媒体下载错误：} 如果媒体超过默认限制，请提高 \texttt{channels.line.mediaMaxMb}。
\end{itemize}



\clearpage

% === 源文件: channels/location.md ===
\section{渠道位置解析}


OpenClaw 将聊天渠道中分享的位置标准化为：

\begin{itemize}
  \item 附加到入站消息体的可读文本，以及
  \item 自动回复上下文负载中的结构化字段。
\end{itemize}


目前支持：

\begin{itemize}
  \item \textbf{Telegram}（位置图钉 + 地点 + 实时位置）
  \item \textbf{WhatsApp}（locationMessage + liveLocationMessage）
  \item \textbf{Matrix}（\texttt{m.location} 配合 \texttt{geo\_uri}）
\end{itemize}


\subsection{文本格式}


位置以友好的行格式呈现，不带括号：

\begin{itemize}
  \item 图钉：
  \item \texttt{\emoji{📍} 48.858844, 2.294351 ±12m}
  \item 命名地点：
  \item \texttt{\emoji{📍} Eiffel Tower — Champ de Mars, Paris (48.858844, 2.294351 ±12m)}
  \item 实时分享：
  \item \texttt{\emoji{🛰} Live location: 48.858844, 2.294351 ±12m}
\end{itemize}


如果渠道包含标题/评论，会附加在下一行：

\begin{lstlisting}
📍 48.858844, 2.294351 ±12m
Meet here
\end{lstlisting}


\subsection{上下文字段}


当存在位置信息时，以下字段会被添加到 \texttt{ctx} 中：

\begin{itemize}
  \item \texttt{LocationLat}（数字）
  \item \texttt{LocationLon}（数字）
  \item \texttt{LocationAccuracy}（数字，米；可选）
  \item \texttt{LocationName}（字符串；可选）
  \item \texttt{LocationAddress}（字符串；可选）
  \item \texttt{LocationSource}（\texttt{pin | place | live}）
  \item \texttt{LocationIsLive}（布尔值）
\end{itemize}


\subsection{渠道说明}


\begin{itemize}
  \item \textbf{Telegram}：地点映射到 \texttt{LocationName/LocationAddress}；实时位置使用 \texttt{live\_period}。
  \item \textbf{WhatsApp}：\texttt{locationMessage.comment} 和 \texttt{liveLocationMessage.caption} 作为标题行附加。
  \item \textbf{Matrix}：\texttt{geo\_uri} 解析为图钉位置；忽略海拔高度，\texttt{LocationIsLive} 始终为 false。
\end{itemize}



\clearpage

% === 源文件: channels/matrix.md ===
\section{Matrix（插件）}


Matrix 是一个开放的去中心化消息协议。OpenClaw 以 Matrix \textbf{用户}身份连接到任意主服务器，因此你需要为机器人创建一个 Matrix 账户。登录后，你可以直接私信机器人或邀请它加入房间（Matrix"群组"）。Beeper 也是一个有效的客户端选项，但它需要启用 E2EE。

状态：通过插件（@vector-im/matrix-bot-sdk）支持。支持私信、房间、话题、媒体、表情回应、投票（发送 + poll-start 作为文本）、位置和 E2EE（需要加密支持）。

\subsection{需要插件}


Matrix 作为插件提供，不包含在核心安装中。

通过 CLI 安装（npm 仓库）：

\begin{lstlisting}[language=bash]
openclaw plugins install @openclaw/matrix
\end{lstlisting}


本地检出（从 git 仓库运行时）：

\begin{lstlisting}[language=bash]
openclaw plugins install ./extensions/matrix
\end{lstlisting}


如果你在配置/新手引导期间选择 Matrix 并检测到 git 检出，OpenClaw 将自动提供本地安装路径。

详情：\href{/tools/plugin}{插件}

\subsection{设置}


\begin{enumerate}
  \item 安装 Matrix 插件：
\end{enumerate}

\begin{itemize}
  \item 从 npm：\texttt{openclaw plugins install @openclaw/matrix}
  \item 从本地检出：\texttt{openclaw plugins install ./extensions/matrix}
\end{itemize}

\begin{enumerate}
  \item 在主服务器上创建 Matrix 账户：
\end{enumerate}

\begin{itemize}
  \item 在 \href{https://matrix.org/ecosystem/hosting/}{https://matrix.org/ecosystem/hosting/} 浏览托管选项
  \item 或自行托管。
\end{itemize}

\begin{enumerate}
  \item 获取机器人账户的访问令牌：
\end{enumerate}

\begin{itemize}
  \item 在你的主服务器上使用 \texttt{curl} 调用 Matrix 登录 API：
\end{itemize}


\begin{lstlisting}[language=bash]
   curl --request POST \
     --url https://matrix.example.org/_matrix/client/v3/login \
     --header 'Content-Type: application/json' \
     --data '{
     "type": "m.login.password",
     "identifier": {
       "type": "m.id.user",
       "user": "your-user-name"
     },
     "password": "your-password"
   }'
\end{lstlisting}


\begin{itemize}
  \item 将 \texttt{matrix.example.org} 替换为你的主服务器 URL。
  \item 或设置 \texttt{channels.matrix.userId} + \texttt{channels.matrix.password}：OpenClaw 会调用相同的登录端点，将访问令牌存储在 \texttt{\textasciitilde{}/.openclaw/credentials/matrix/credentials.json}，并在下次启动时重用。
\end{itemize}


\begin{enumerate}
  \item 配置凭证：
\end{enumerate}

\begin{itemize}
  \item 环境变量：\texttt{MATRIX\_HOMESERVER}、\texttt{MATRIX\_ACCESS\_TOKEN}（或 \texttt{MATRIX\_USER\_ID} + \texttt{MATRIX\_PASSWORD}）
  \item 或配置：\texttt{channels.matrix.*}
  \item 如果两者都设置，配置优先。
  \item 使用访问令牌时：用户 ID 通过 \texttt{/whoami} 自动获取。
  \item 设置时，\texttt{channels.matrix.userId} 应为完整的 Matrix ID（示例：\texttt{@bot:example.org}）。
\end{itemize}

\begin{enumerate}
  \item 重启 Gateway 网关（或完成新手引导）。
  \item 从任何 Matrix 客户端（Element、Beeper 等；参见 https://matrix.org/ecosystem/clients/）与机器人开始私信或邀请它加入房间。Beeper 需要 E2EE，因此请设置 \texttt{channels.matrix.encryption: true} 并验证设备。
\end{enumerate}


最小配置（访问令牌，用户 ID 自动获取）：

\begin{lstlisting}[language=json]
{
  channels: {
    matrix: {
      enabled: true,
      homeserver: "https://matrix.example.org",
      accessToken: "syt_***",
      dm: { policy: "pairing" },
    },
  },
}
\end{lstlisting}


E2EE 配置（启用端到端加密）：

\begin{lstlisting}[language=json]
{
  channels: {
    matrix: {
      enabled: true,
      homeserver: "https://matrix.example.org",
      accessToken: "syt_***",
      encryption: true,
      dm: { policy: "pairing" },
    },
  },
}
\end{lstlisting}


\subsection{加密（E2EE）}


通过 Rust 加密 SDK \textbf{支持}端到端加密。

使用 \texttt{channels.matrix.encryption: true} 启用：

\begin{itemize}
  \item 如果加密模块加载成功，加密房间会自动解密。
  \item 发送到加密房间时，出站媒体会被加密。
  \item 首次连接时，OpenClaw 会向你的其他会话请求设备验证。
  \item 在另一个 Matrix 客户端（Element 等）中验证设备以启用密钥共享。
  \item 如果无法加载加密模块，E2EE 将被禁用，加密房间将无法解密；OpenClaw 会记录警告。
  \item 如果你看到缺少加密模块的错误（例如 \texttt{@matrix-org/matrix-sdk-crypto-nodejs-*}），请允许 \texttt{@matrix-org/matrix-sdk-crypto-nodejs} 的构建脚本并运行 \texttt{pnpm rebuild @matrix-org/matrix-sdk-crypto-nodejs}，或使用 \texttt{node node\_modules/@matrix-org/matrix-sdk-crypto-nodejs/download-lib.js} 获取二进制文件。
\end{itemize}


加密状态按账户 + 访问令牌存储在 \texttt{\textasciitilde{}/.openclaw/matrix/accounts//\_\_//crypto/}（SQLite 数据库）。同步状态存储在同目录的 \texttt{bot-storage.json} 中。如果访问令牌（设备）更改，将创建新的存储，机器人必须重新验证才能访问加密房间。

\textbf{设备验证：}
启用 E2EE 时，机器人将在启动时向你的其他会话请求验证。打开 Element（或其他客户端）并批准验证请求以建立信任。验证后，机器人可以解密加密房间中的消息。

\subsection{路由模型}


\begin{itemize}
  \item 回复始终返回到 Matrix。
  \item 私信共享智能体的主会话；房间映射到群组会话。
\end{itemize}


\subsection{访问控制（私信）}


\begin{itemize}
  \item 默认：\texttt{channels.matrix.dm.policy = "pairing"}。未知发送者会收到配对码。
  \item 通过以下方式批准：
  \item \texttt{openclaw pairing list matrix}
  \item \texttt{openclaw pairing approve matrix }
  \item 公开私信：\texttt{channels.matrix.dm.policy="open"} 加上 \texttt{channels.matrix.dm.allowFrom=["*"]}。
  \item \texttt{channels.matrix.dm.allowFrom} 仅接受完整 Matrix 用户 ID（例如 \texttt{@user:server}）。向导仅在目录搜索得到唯一精确匹配时将显示名称解析为用户 ID。
\end{itemize}


\subsection{房间（群组）}


\begin{itemize}
  \item 默认：\texttt{channels.matrix.groupPolicy = "allowlist"}（提及门控）。使用 \texttt{channels.defaults.groupPolicy} 在未设置时覆盖默认值。
  \item 使用 \texttt{channels.matrix.groups} 配置房间允许列表（房间 ID 或别名；名称仅在目录搜索得到唯一精确匹配时解析为 ID）：
\end{itemize}


\begin{lstlisting}[language=json]
{
  channels: {
    matrix: {
      groupPolicy: "allowlist",
      groups: {
        "!roomId:example.org": { allow: true },
        "#alias:example.org": { allow: true },
      },
      groupAllowFrom: ["@owner:example.org"],
    },
  },
}
\end{lstlisting}


\begin{itemize}
  \item \texttt{requireMention: false} 启用该房间的自动回复。
  \item \texttt{groups."*"} 可以设置跨房间的提及门控默认值。
  \item \texttt{groupAllowFrom} 限制哪些发送者可以在房间中触发机器人（需完整 Matrix 用户 ID）。
  \item 每个房间的 \texttt{users} 允许列表可以进一步限制特定房间内的发送者（需完整 Matrix 用户 ID）。
  \item 配置向导会提示输入房间允许列表（房间 ID、别名或名称），仅在精确且唯一匹配时解析名称。
  \item 启动时，OpenClaw 将允许列表中的房间/用户名称解析为 ID 并记录映射；未解析的条目不会参与允许列表匹配。
  \item 默认自动加入邀请；使用 \texttt{channels.matrix.autoJoin} 和 \texttt{channels.matrix.autoJoinAllowlist} 控制。
  \item 要\textbf{禁止所有房间}，设置 \texttt{channels.matrix.groupPolicy: "disabled"}（或保持空的允许列表）。
  \item 旧版键名：\texttt{channels.matrix.rooms}（与 \texttt{groups} 相同的结构）。
\end{itemize}


\subsection{话题}


\begin{itemize}
  \item 支持回复话题。
  \item \texttt{channels.matrix.threadReplies} 控制回复是否保持在话题中：
  \item \texttt{off}、\texttt{inbound}（默认）、\texttt{always}
  \item \texttt{channels.matrix.replyToMode} 控制不在话题中回复时的 reply-to 元数据：
  \item \texttt{off}（默认）、\texttt{first}、\texttt{all}
\end{itemize}


\subsection{功能}



\begin{table}[h]
\centering
\small
\begin{tabular}{|l|l|}
\hline
功能 & 状态 \\
\hline
私信 & \emoji{✅} 支持 \\
\hline
房间 & \emoji{✅} 支持 \\
\hline
话题 & \emoji{✅} 支持 \\
\hline
媒体 & \emoji{✅} 支持 \\
\hline
E2EE & \emoji{✅} 支持（需要加密模块） \\
\hline
表情回应 & \emoji{✅} 支持（通过工具发送/读取） \\
\hline
投票 & \emoji{✅} 支持发送；入站投票开始转换为文本（响应/结束被忽略） \\
\hline
位置 & \emoji{✅} 支持（geo URI；忽略海拔） \\
\hline
原生命令 & \emoji{✅} 支持 \\
\hline
\end{tabular}
\end{table}



\subsection{配置参考（Matrix）}


完整配置：\href{/gateway/configuration}{配置}

提供商选项：

\begin{itemize}
  \item \texttt{channels.matrix.enabled}：启用/禁用渠道启动。
  \item \texttt{channels.matrix.homeserver}：主服务器 URL。
  \item \texttt{channels.matrix.userId}：Matrix 用户 ID（使用访问令牌时可选）。
  \item \texttt{channels.matrix.accessToken}：访问令牌。
  \item \texttt{channels.matrix.password}：登录密码（令牌会被存储）。
  \item \texttt{channels.matrix.deviceName}：设备显示名称。
  \item \texttt{channels.matrix.encryption}：启用 E2EE（默认：false）。
  \item \texttt{channels.matrix.initialSyncLimit}：初始同步限制。
  \item \texttt{channels.matrix.threadReplies}：\texttt{off | inbound | always}（默认：inbound）。
  \item \texttt{channels.matrix.textChunkLimit}：出站文本分块大小（字符）。
  \item \texttt{channels.matrix.chunkMode}：\texttt{length}（默认）或 \texttt{newline} 在长度分块前按空行（段落边界）分割。
  \item \texttt{channels.matrix.dm.policy}：\texttt{pairing | allowlist | open | disabled}（默认：pairing）。
  \item \texttt{channels.matrix.dm.allowFrom}：私信允许列表（需完整 Matrix 用户 ID）。\texttt{open} 需要 \texttt{"*"}。向导在可能时将名称解析为 ID。
  \item \texttt{channels.matrix.groupPolicy}：\texttt{allowlist | open | disabled}（默认：allowlist）。
  \item \texttt{channels.matrix.groupAllowFrom}：群组消息的允许发送者列表（需完整 Matrix 用户 ID）。
  \item \texttt{channels.matrix.allowlistOnly}：强制私信 + 房间使用允许列表规则。
  \item \texttt{channels.matrix.groups}：群组允许列表 + 每个房间的设置映射。
  \item \texttt{channels.matrix.rooms}：旧版群组允许列表/配置。
  \item \texttt{channels.matrix.replyToMode}：话题/标签的 reply-to 模式。
  \item \texttt{channels.matrix.mediaMaxMb}：入站/出站媒体上限（MB）。
  \item \texttt{channels.matrix.autoJoin}：邀请处理（\texttt{always | allowlist | off}，默认：always）。
  \item \texttt{channels.matrix.autoJoinAllowlist}：自动加入的允许房间 ID/别名。
  \item \texttt{channels.matrix.actions}：每个操作的工具限制（reactions/messages/pins/memberInfo/channelInfo）。
\end{itemize}



\clearpage

% === 源文件: channels/mattermost.md ===
\section{Mattermost（插件）}


状态：通过插件支持（bot token + WebSocket 事件）。支持频道、群组和私信。
Mattermost 是一个可自托管的团队消息平台；有关产品详情和下载，请访问官方网站
\href{https://mattermost.com}{mattermost.com}。

\subsection{需要插件}


Mattermost 以插件形式提供，不包含在核心安装中。

通过 CLI 安装（npm 注册表）：

\begin{lstlisting}[language=bash]
openclaw plugins install @openclaw/mattermost
\end{lstlisting}


本地检出（从 git 仓库运行时）：

\begin{lstlisting}[language=bash]
openclaw plugins install ./extensions/mattermost
\end{lstlisting}


如果你在配置/新手引导期间选择 Mattermost 并检测到 git 检出，OpenClaw 会自动提供本地安装路径。

详情：\href{/tools/plugin}{插件}

\subsection{快速设置}


\begin{enumerate}
  \item 安装 Mattermost 插件。
  \item 创建 Mattermost bot 账户并复制 \textbf{bot token}。
  \item 复制 Mattermost \textbf{基础 URL}（例如 \texttt{https://chat.example.com}）。
  \item 配置 OpenClaw 并启动 Gateway 网关。
\end{enumerate}


最小配置：

\begin{lstlisting}[language=json]
{
  channels: {
    mattermost: {
      enabled: true,
      botToken: "mm-token",
      baseUrl: "https://chat.example.com",
      dmPolicy: "pairing",
    },
  },
}
\end{lstlisting}


\subsection{环境变量（默认账户）}


如果你偏好使用环境变量，请在 Gateway 网关主机上设置：

\begin{itemize}
  \item \texttt{MATTERMOST\_BOT\_TOKEN=...}
  \item \texttt{MATTERMOST\_URL=https://chat.example.com}
\end{itemize}


环境变量仅适用于\textbf{默认}账户（\texttt{default}）。其他账户必须使用配置值。

\subsection{聊天模式}


Mattermost 自动响应私信。频道行为由 \texttt{chatmode} 控制：

\begin{itemize}
  \item \texttt{oncall}（默认）：仅在频道中被 @提及时响应。
  \item \texttt{onmessage}：响应每条频道消息。
  \item \texttt{onchar}：当消息以触发前缀开头时响应。
\end{itemize}


配置示例：

\begin{lstlisting}[language=json]
{
  channels: {
    mattermost: {
      chatmode: "onchar",
      oncharPrefixes: [">", "!"],
    },
  },
}
\end{lstlisting}


注意事项：

\begin{itemize}
  \item \texttt{onchar} 仍会响应显式 @提及。
  \item \texttt{channels.mattermost.requireMention} 对旧配置仍然有效，但推荐使用 \texttt{chatmode}。
\end{itemize}


\subsection{访问控制（私信）}


\begin{itemize}
  \item 默认：\texttt{channels.mattermost.dmPolicy = "pairing"}（未知发送者会收到配对码）。
  \item 通过以下方式批准：
  \item \texttt{openclaw pairing list mattermost}
  \item \texttt{openclaw pairing approve mattermost }
  \item 公开私信：\texttt{channels.mattermost.dmPolicy="open"} 加上 \texttt{channels.mattermost.allowFrom=["*"]}。
\end{itemize}


\subsection{频道（群组）}


\begin{itemize}
  \item 默认：\texttt{channels.mattermost.groupPolicy = "allowlist"}（提及限制）。
  \item 使用 \texttt{channels.mattermost.groupAllowFrom} 将发送者加入允许列表（用户 ID 或 \texttt{@username}）。
  \item 开放频道：\texttt{channels.mattermost.groupPolicy="open"}（提及限制）。
\end{itemize}


\subsection{出站投递目标}


在 \texttt{openclaw message send} 或 cron/webhooks 中使用这些目标格式：

\begin{itemize}
  \item \texttt{channel:} 用于频道
  \item \texttt{user:} 用于私信
  \item \texttt{@username} 用于私信（通过 Mattermost API 解析）
\end{itemize}


裸 ID 被视为频道。

\subsection{多账户}


Mattermost 支持在 \texttt{channels.mattermost.accounts} 下配置多个账户：

\begin{lstlisting}[language=json]
{
  channels: {
    mattermost: {
      accounts: {
        default: { name: "Primary", botToken: "mm-token", baseUrl: "https://chat.example.com" },
        alerts: { name: "Alerts", botToken: "mm-token-2", baseUrl: "https://alerts.example.com" },
      },
    },
  },
}
\end{lstlisting}


\subsection{故障排除}


\begin{itemize}
  \item 频道中无回复：确保 bot 在频道中并提及它（oncall），使用触发前缀（onchar），或设置 \texttt{chatmode: "onmessage"}。
  \item 认证错误：检查 bot token、基础 URL 以及账户是否已启用。
  \item 多账户问题：环境变量仅适用于 \texttt{default} 账户。
\end{itemize}



\clearpage

% === 源文件: channels/msteams.md ===
\section{Microsoft Teams（插件）}


\begin{quote}
"进入此地者，放弃一切希望。"
\end{quote}


更新时间：2026-01-21

状态：支持文本 + 私信附件；频道/群组文件发送需要 \texttt{sharePointSiteId} + Graph 权限（参见\href{\#sending-files-in-group-chats}{在群聊中发送文件}）。投票通过 Adaptive Cards 发送。

\subsection{需要插件}


Microsoft Teams 作为插件提供，不包含在核心安装中。

\textbf{破坏性变更（2026.1.15）：} MS Teams 已从核心移出。如果你使用它，必须安装插件。

原因说明：保持核心安装更轻量，并让 MS Teams 依赖项可以独立更新。

通过 CLI 安装（npm 注册表）：

\begin{lstlisting}[language=bash]
openclaw plugins install @openclaw/msteams
\end{lstlisting}


本地检出（从 git 仓库运行时）：

\begin{lstlisting}[language=bash]
openclaw plugins install ./extensions/msteams
\end{lstlisting}


如果你在配置/新手引导过程中选择 Teams 并检测到 git 检出，
OpenClaw 将自动提供本地安装路径。

详情：\href{/tools/plugin}{插件}

\subsection{快速设置（初学者）}


\begin{enumerate}
  \item 安装 Microsoft Teams 插件。
  \item 创建一个 \textbf{Azure Bot}（App ID + 客户端密钥 + 租户 ID）。
  \item 使用这些凭证配置 OpenClaw。
  \item 通过公共 URL 或隧道暴露 \texttt{/api/messages}（默认端口 3978）。
  \item 安装 Teams 应用包并启动 Gateway 网关。
\end{enumerate}


最小配置：

\begin{lstlisting}[language=json]
{
  channels: {
    msteams: {
      enabled: true,
      appId: "<APP_ID>",
      appPassword: "<APP_PASSWORD>",
      tenantId: "<TENANT_ID>",
      webhook: { port: 3978, path: "/api/messages" },
    },
  },
}
\end{lstlisting}


注意：群聊默认被阻止（\texttt{channels.msteams.groupPolicy: "allowlist"}）。要允许群组回复，请设置 \texttt{channels.msteams.groupAllowFrom}（或使用 \texttt{groupPolicy: "open"} 允许任何成员，需要提及才能触发）。

\subsection{目标}


\begin{itemize}
  \item 通过 Teams 私信、群聊或频道与 OpenClaw 交流。
  \item 保持路由确定性：回复始终返回到消息到达的渠道。
  \item 默认使用安全的渠道行为（除非另有配置，否则需要提及）。
\end{itemize}


\subsection{配置写入}


默认情况下，Microsoft Teams 允许通过 \texttt{/config set|unset} 触发的配置更新写入（需要 \texttt{commands.config: true}）。

禁用方式：

\begin{lstlisting}[language=json]
{
  channels: { msteams: { configWrites: false } },
}
\end{lstlisting}


\subsection{访问控制（私信 + 群组）}


\textbf{私信访问}

\begin{itemize}
  \item 默认：\texttt{channels.msteams.dmPolicy = "pairing"}。未知发送者在获得批准之前将被忽略。
  \item \texttt{channels.msteams.allowFrom} 接受 AAD 对象 ID、UPN 或显示名称。当凭证允许时，向导会通过 Microsoft Graph 将名称解析为 ID。
\end{itemize}


\textbf{群组访问}

\begin{itemize}
  \item 默认：\texttt{channels.msteams.groupPolicy = "allowlist"}（除非添加 \texttt{groupAllowFrom}，否则被阻止）。使用 \texttt{channels.defaults.groupPolicy} 在未设置时覆盖默认值。
  \item \texttt{channels.msteams.groupAllowFrom} 控制哪些发送者可以在群聊/频道中触发（回退到 \texttt{channels.msteams.allowFrom}）。
  \item 设置 \texttt{groupPolicy: "open"} 允许任何成员（默认仍需提及才能触发）。
  \item 要\textbf{不允许任何频道}，设置 \texttt{channels.msteams.groupPolicy: "disabled"}。
\end{itemize}


示例：

\begin{lstlisting}[language=json]
{
  channels: {
    msteams: {
      groupPolicy: "allowlist",
      groupAllowFrom: ["user@org.com"],
    },
  },
}
\end{lstlisting}


\textbf{团队 + 频道允许列表}

\begin{itemize}
  \item 通过在 \texttt{channels.msteams.teams} 下列出团队和频道来限定群组/频道回复的范围。
  \item 键可以是团队 ID 或名称；频道键可以是会话 ID 或名称。
  \item 当 \texttt{groupPolicy="allowlist"} 且存在团队允许列表时，仅接受列出的团队/频道（需要提及才能触发）。
  \item 配置向导接受 \texttt{Team/Channel} 条目并为你存储。
  \item 启动时，OpenClaw 将团队/频道和用户允许列表名称解析为 ID（当 Graph 权限允许时）
\end{itemize}

并记录映射；未解析的条目保持原样。

示例：

\begin{lstlisting}[language=json]
{
  channels: {
    msteams: {
      groupPolicy: "allowlist",
      teams: {
        "My Team": {
          channels: {
            General: { requireMention: true },
          },
        },
      },
    },
  },
}
\end{lstlisting}


\subsection{工作原理}


\begin{enumerate}
  \item 安装 Microsoft Teams 插件。
  \item 创建一个 \textbf{Azure Bot}（App ID + 密钥 + 租户 ID）。
  \item 构建一个引用机器人并包含以下 RSC 权限的 \textbf{Teams 应用包}。
  \item 将 Teams 应用上传/安装到团队中（或用于私信的个人范围）。
  \item 在 \texttt{\textasciitilde{}/.openclaw/openclaw.json}（或环境变量）中配置 \texttt{msteams} 并启动 Gateway 网关。
  \item Gateway 网关默认在 \texttt{/api/messages} 上监听 Bot Framework webhook 流量。
\end{enumerate}


\subsection{Azure Bot 设置（前提条件）}


在配置 OpenClaw 之前，你需要创建一个 Azure Bot 资源。

\subsubsection{步骤 1：创建 Azure Bot}


\begin{enumerate}
  \item 前往\href{https://portal.azure.com/\#create/Microsoft.AzureBot}{创建 Azure Bot}
  \item 填写\textbf{基本信息}选项卡：
\end{enumerate}



\begin{table}[h]
\centering
\small
\begin{tabular}{|l|l|}
\hline
字段 & 值 \\
\hline
\textbf{Bot handle} & 你的机器人名称，例如 \texttt{openclaw-msteams}（必须唯一） \\
\hline
\textbf{Subscription} & 选择你的 Azure 订阅 \\
\hline
\textbf{Resource group} & 新建或使用现有 \\
\hline
\textbf{Pricing tier} & \textbf{Free} 用于开发/测试 \\
\hline
\textbf{Type of App} & \textbf{Single Tenant}（推荐 - 见下方说明） \\
\hline
\textbf{Creation type} & \textbf{Create new Microsoft App ID} \\
\hline
\end{tabular}
\end{table}



\begin{quote}
\textbf{弃用通知：} 2025-07-31 之后已弃用创建新的多租户机器人。新机器人请使用 \textbf{Single Tenant}。
\end{quote}


\begin{enumerate}
  \item 点击 \textbf{Review + create} → \textbf{Create}（等待约 1-2 分钟）
\end{enumerate}


\subsubsection{步骤 2：获取凭证}


\begin{enumerate}
  \item 前往你的 Azure Bot 资源 → \textbf{Configuration}
  \item 复制 \textbf{Microsoft App ID} → 这是你的 \texttt{appId}
  \item 点击 \textbf{Manage Password} → 前往应用注册
  \item 在 \textbf{Certificates \& secrets} → \textbf{New client secret} → 复制 \textbf{Value} → 这是你的 \texttt{appPassword}
  \item 前往 \textbf{Overview} → 复制 \textbf{Directory (tenant) ID} → 这是你的 \texttt{tenantId}
\end{enumerate}


\subsubsection{步骤 3：配置消息端点}


\begin{enumerate}
  \item 在 Azure Bot → \textbf{Configuration}
  \item 将 \textbf{Messaging endpoint} 设置为你的 webhook URL：
\end{enumerate}

\begin{itemize}
  \item 生产环境：\texttt{https://your-domain.com/api/messages}
  \item 本地开发：使用隧道（见下方\href{\#local-development-tunneling}{本地开发}）
\end{itemize}


\subsubsection{步骤 4：启用 Teams 渠道}


\begin{enumerate}
  \item 在 Azure Bot → \textbf{Channels}
  \item 点击 \textbf{Microsoft Teams} → Configure → Save
  \item 接受服务条款
\end{enumerate}


\subsection{本地开发（隧道）}


Teams 无法访问 \texttt{localhost}。本地开发请使用隧道：

\textbf{选项 A：ngrok}

\begin{lstlisting}[language=bash]
ngrok http 3978
# 复制 https URL，例如 https://abc123.ngrok.io
# 将消息端点设置为：https://abc123.ngrok.io/api/messages
\end{lstlisting}


\textbf{选项 B：Tailscale Funnel}

\begin{lstlisting}[language=bash]
tailscale funnel 3978
# 使用你的 Tailscale funnel URL 作为消息端点
\end{lstlisting}


\subsection{Teams 开发者门户（替代方案）}


除了手动创建清单 ZIP，你可以使用 \href{https://dev.teams.microsoft.com/apps}{Teams 开发者门户}：

\begin{enumerate}
  \item 点击 \textbf{+ New app}
  \item 填写基本信息（名称、描述、开发者信息）
  \item 前往 \textbf{App features} → \textbf{Bot}
  \item 选择 \textbf{Enter a bot ID manually} 并粘贴你的 Azure Bot App ID
  \item 勾选范围：\textbf{Personal}、\textbf{Team}、\textbf{Group Chat}
  \item 点击 \textbf{Distribute} → \textbf{Download app package}
  \item 在 Teams 中：\textbf{Apps} → \textbf{Manage your apps} → \textbf{Upload a custom app} → 选择 ZIP
\end{enumerate}


这通常比手动编辑 JSON 清单更容易。

\subsection{测试机器人}


\textbf{选项 A：Azure Web Chat（先验证 webhook）}

\begin{enumerate}
  \item 在 Azure 门户 → 你的 Azure Bot 资源 → \textbf{Test in Web Chat}
  \item 发送一条消息 - 你应该看到响应
  \item 这确认你的 webhook 端点在 Teams 设置之前正常工作
\end{enumerate}


\textbf{选项 B：Teams（应用安装后）}

\begin{enumerate}
  \item 安装 Teams 应用（侧载或组织目录）
  \item 在 Teams 中找到机器人并发送私信
  \item 检查 Gateway 网关日志中的传入活动
\end{enumerate}


\subsection{设置（最小纯文本）}


\begin{enumerate}
  \item \textbf{安装 Microsoft Teams 插件}
\end{enumerate}

\begin{itemize}
  \item 从 npm：\texttt{openclaw plugins install @openclaw/msteams}
  \item 从本地检出：\texttt{openclaw plugins install ./extensions/msteams}
\end{itemize}


\begin{enumerate}
  \item \textbf{机器人注册}
\end{enumerate}

\begin{itemize}
  \item 创建一个 Azure Bot（见上文）并记录：
  \item App ID
  \item 客户端密钥（App password）
  \item 租户 ID（单租户）
\end{itemize}


\begin{enumerate}
  \item \textbf{Teams 应用清单}
\end{enumerate}

\begin{itemize}
  \item 包含一个 \texttt{bot} 条目，其中 \texttt{botId = }。
  \item 范围：\texttt{personal}、\texttt{team}、\texttt{groupChat}。
  \item \texttt{supportsFiles: true}（个人范围文件处理所需）。
  \item 添加 RSC 权限（见下文）。
  \item 创建图标：\texttt{outline.png}（32x32）和 \texttt{color.png}（192x192）。
  \item 将三个文件一起打包：\texttt{manifest.json}、\texttt{outline.png}、\texttt{color.png}。
\end{itemize}


\begin{enumerate}
  \item \textbf{配置 OpenClaw}
\end{enumerate}


\begin{lstlisting}[language=json]
   {
     "msteams": {
       "enabled": true,
       "appId": "<APP_ID>",
       "appPassword": "<APP_PASSWORD>",
       "tenantId": "<TENANT_ID>",
       "webhook": { "port": 3978, "path": "/api/messages" }
     }
   }
\end{lstlisting}


你也可以使用环境变量代替配置键：
\begin{itemize}
  \item \texttt{MSTEAMS\_APP\_ID}
  \item \texttt{MSTEAMS\_APP\_PASSWORD}
  \item \texttt{MSTEAMS\_TENANT\_ID}
\end{itemize}


\begin{enumerate}
  \item \textbf{机器人端点}
\end{enumerate}

\begin{itemize}
  \item 将 Azure Bot Messaging Endpoint 设置为：
  \item \texttt{https://:3978/api/messages}（或你选择的路径/端口）。
\end{itemize}


\begin{enumerate}
  \item \textbf{运行 Gateway 网关}
\end{enumerate}

\begin{itemize}
  \item 当插件已安装且 \texttt{msteams} 配置存在并有凭证时，Teams 渠道会自动启动。
\end{itemize}


\subsection{历史上下文}


\begin{itemize}
  \item \texttt{channels.msteams.historyLimit} 控制将多少条最近的频道/群组消息包含到提示中。
  \item 回退到 \texttt{messages.groupChat.historyLimit}。设置 \texttt{0} 禁用（默认 50）。
  \item 私信历史可以通过 \texttt{channels.msteams.dmHistoryLimit}（用户轮次）限制。每用户覆盖：\texttt{channels.msteams.dms[""].historyLimit}。
\end{itemize}


\subsection{当前 Teams RSC 权限（清单）}


这些是我们 Teams 应用清单中\textbf{现有的 resourceSpecific 权限}。它们仅适用于安装了应用的团队/聊天内部。

\textbf{对于频道（团队范围）：}

\begin{itemize}
  \item \texttt{ChannelMessage.Read.Group}（Application）- 无需 @提及即可接收所有频道消息
  \item \texttt{ChannelMessage.Send.Group}（Application）
  \item \texttt{Member.Read.Group}（Application）
  \item \texttt{Owner.Read.Group}（Application）
  \item \texttt{ChannelSettings.Read.Group}（Application）
  \item \texttt{TeamMember.Read.Group}（Application）
  \item \texttt{TeamSettings.Read.Group}（Application）
\end{itemize}


\textbf{对于群聊：}

\begin{itemize}
  \item \texttt{ChatMessage.Read.Chat}（Application）- 无需 @提及即可接收所有群聊消息
\end{itemize}


\subsection{Teams 清单示例（已脱敏）}


包含必需字段的最小有效示例。请替换 ID 和 URL。

\begin{lstlisting}[language=json]
{
  "$schema": "https://developer.microsoft.com/en-us/json-schemas/teams/v1.23/MicrosoftTeams.schema.json",
  "manifestVersion": "1.23",
  "version": "1.0.0",
  "id": "00000000-0000-0000-0000-000000000000",
  "name": { "short": "OpenClaw" },
  "developer": {
    "name": "Your Org",
    "websiteUrl": "https://example.com",
    "privacyUrl": "https://example.com/privacy",
    "termsOfUseUrl": "https://example.com/terms"
  },
  "description": { "short": "OpenClaw in Teams", "full": "OpenClaw in Teams" },
  "icons": { "outline": "outline.png", "color": "color.png" },
  "accentColor": "#5B6DEF",
  "bots": [
    {
      "botId": "11111111-1111-1111-1111-111111111111",
      "scopes": ["personal", "team", "groupChat"],
      "isNotificationOnly": false,
      "supportsCalling": false,
      "supportsVideo": false,
      "supportsFiles": true
    }
  ],
  "webApplicationInfo": {
    "id": "11111111-1111-1111-1111-111111111111"
  },
  "authorization": {
    "permissions": {
      "resourceSpecific": [
        { "name": "ChannelMessage.Read.Group", "type": "Application" },
        { "name": "ChannelMessage.Send.Group", "type": "Application" },
        { "name": "Member.Read.Group", "type": "Application" },
        { "name": "Owner.Read.Group", "type": "Application" },
        { "name": "ChannelSettings.Read.Group", "type": "Application" },
        { "name": "TeamMember.Read.Group", "type": "Application" },
        { "name": "TeamSettings.Read.Group", "type": "Application" },
        { "name": "ChatMessage.Read.Chat", "type": "Application" }
      ]
    }
  }
}
\end{lstlisting}


\subsubsection{清单注意事项（必填字段）}


\begin{itemize}
  \item \texttt{bots[].botId} \textbf{必须}与 Azure Bot App ID 匹配。
  \item \texttt{webApplicationInfo.id} \textbf{必须}与 Azure Bot App ID 匹配。
  \item \texttt{bots[].scopes} 必须包含你计划使用的界面（\texttt{personal}、\texttt{team}、\texttt{groupChat}）。
  \item \texttt{bots[].supportsFiles: true} 是个人范围文件处理所需的。
  \item \texttt{authorization.permissions.resourceSpecific} 如果你需要频道流量，必须包含频道读取/发送权限。
\end{itemize}


\subsubsection{更新现有应用}


要更新已安装的 Teams 应用（例如，添加 RSC 权限）：

\begin{enumerate}
  \item 使用新设置更新你的 \texttt{manifest.json}
  \item \textbf{增加 \texttt{version} 字段}（例如，\texttt{1.0.0} → \texttt{1.1.0}）
  \item \textbf{重新打包}清单和图标（\texttt{manifest.json}、\texttt{outline.png}、\texttt{color.png}）
  \item 上传新的 zip：
\end{enumerate}

\begin{itemize}
  \item \textbf{选项 A（Teams 管理中心）：} Teams 管理中心 → Teams apps → Manage apps → 找到你的应用 → Upload new version
  \item \textbf{选项 B（侧载）：} 在 Teams 中 → Apps → Manage your apps → Upload a custom app
\end{itemize}

\begin{enumerate}
  \item \textbf{对于团队频道：} 在每个团队中重新安装应用以使新权限生效
  \item \textbf{完全退出并重新启动 Teams}（不仅仅是关闭窗口）以清除缓存的应用元数据
\end{enumerate}


\subsection{功能：仅 RSC 与 Graph}


\subsubsection{仅使用 **Teams RSC**（应用已安装，无 Graph API 权限）}


可用：

\begin{itemize}
  \item 读取频道消息\textbf{文本}内容。
  \item 发送频道消息\textbf{文本}内容。
  \item 接收\textbf{个人（私信）}文件附件。
\end{itemize}


不可用：

\begin{itemize}
  \item 频道/群组\textbf{图片或文件内容}（负载仅包含 HTML 存根）。
  \item 下载存储在 SharePoint/OneDrive 中的附件。
  \item 读取消息历史（超出实时 webhook 事件）。
\end{itemize}


\subsubsection{使用 **Teams RSC + Microsoft Graph Application 权限**}


增加：

\begin{itemize}
  \item 下载托管内容（粘贴到消息中的图片）。
  \item 下载存储在 SharePoint/OneDrive 中的文件附件。
  \item 通过 Graph 读取频道/聊天消息历史。
\end{itemize}


\subsubsection{RSC 与 Graph API 对比}



\begin{table}[h]
\centering
\small
\begin{tabular}{|l|l|l|}
\hline
功能 & RSC 权限 & Graph API \\
\hline
\textbf{实时消息} & 是（通过 webhook） & 否（仅轮询） \\
\hline
\textbf{历史消息} & 否 & 是（可查询历史） \\
\hline
\textbf{设置复杂度} & 仅应用清单 & 需要管理员同意 + 令牌流程 \\
\hline
\textbf{离线工作} & 否（必须运行） & 是（随时查询） \\
\hline
\end{tabular}
\end{table}



\textbf{结论：} RSC 用于实时监听；Graph API 用于历史访问。要在离线时补上错过的消息，你需要带有 \texttt{ChannelMessage.Read.All} 的 Graph API（需要管理员同意）。

\subsection{启用 Graph 的媒体 + 历史（频道所需）}


如果你需要\textbf{频道}中的图片/文件或想要获取\textbf{消息历史}，你必须启用 Microsoft Graph 权限并授予管理员同意。

\begin{enumerate}
  \item 在 Entra ID（Azure AD）\textbf{App Registration} 中，添加 Microsoft Graph \textbf{Application 权限}：
\end{enumerate}

\begin{itemize}
  \item \texttt{ChannelMessage.Read.All}（频道附件 + 历史）
  \item \texttt{Chat.Read.All} 或 \texttt{ChatMessage.Read.All}（群聊）
\end{itemize}

\begin{enumerate}
  \item 为租户\textbf{授予管理员同意}。
  \item 提升 Teams 应用\textbf{清单版本}，重新上传，并\textbf{在 Teams 中重新安装应用}。
  \item \textbf{完全退出并重新启动 Teams} 以清除缓存的应用元数据。
\end{enumerate}


\subsection{已知限制}


\subsubsection{Webhook 超时}


Teams 通过 HTTP webhook 传递消息。如果处理时间过长（例如，LLM 响应缓慢），你可能会看到：

\begin{itemize}
  \item Gateway 网关超时
  \item Teams 重试消息（导致重复）
  \item 丢失的回复
\end{itemize}


OpenClaw 通过快速返回并主动发送回复来处理这个问题，但非常慢的响应仍可能导致问题。

\subsubsection{格式化}


Teams markdown 比 Slack 或 Discord 更有限：

\begin{itemize}
  \item 基本格式化有效：\textbf{粗体}、\_斜体\_、\texttt{代码}、链接
  \item 复杂的 markdown（表格、嵌套列表）可能无法正确渲染
  \item 支持 Adaptive Cards 用于投票和任意卡片发送（见下文）
\end{itemize}


\subsection{配置}


关键设置（共享渠道模式见 \texttt{/gateway/configuration}）：

\begin{itemize}
  \item \texttt{channels.msteams.enabled}：启用/禁用渠道。
  \item \texttt{channels.msteams.appId}、\texttt{channels.msteams.appPassword}、\texttt{channels.msteams.tenantId}：机器人凭证。
  \item \texttt{channels.msteams.webhook.port}（默认 \texttt{3978}）
  \item \texttt{channels.msteams.webhook.path}（默认 \texttt{/api/messages}）
  \item \texttt{channels.msteams.dmPolicy}：\texttt{pairing | allowlist | open | disabled}（默认：pairing）
  \item \texttt{channels.msteams.allowFrom}：私信允许列表（AAD 对象 ID、UPN 或显示名称）。当 Graph 访问可用时，向导在设置期间将名称解析为 ID。
  \item \texttt{channels.msteams.textChunkLimit}：出站文本分块大小。
  \item \texttt{channels.msteams.chunkMode}：\texttt{length}（默认）或 \texttt{newline} 在长度分块之前按空行（段落边界）分割。
  \item \texttt{channels.msteams.mediaAllowHosts}：入站附件主机允许列表（默认为 Microsoft/Teams 域名）。
  \item \texttt{channels.msteams.mediaAuthAllowHosts}：在媒体重试时附加 Authorization 头的允许列表（默认为 Graph + Bot Framework 主机）。
  \item \texttt{channels.msteams.requireMention}：在频道/群组中需要 @提及（默认 true）。
  \item \texttt{channels.msteams.replyStyle}：\texttt{thread | top-level}（见\href{\#reply-style-threads-vs-posts}{回复样式}）。
  \item \texttt{channels.msteams.teams..replyStyle}：每团队覆盖。
  \item \texttt{channels.msteams.teams..requireMention}：每团队覆盖。
  \item \texttt{channels.msteams.teams..tools}：当缺少频道覆盖时使用的默认每团队工具策略覆盖（\texttt{allow}/\texttt{deny}/\texttt{alsoAllow}）。
  \item \texttt{channels.msteams.teams..toolsBySender}：默认每团队每发送者工具策略覆盖（支持 \texttt{"*"} 通配符）。
  \item \texttt{channels.msteams.teams..channels..replyStyle}：每频道覆盖。
  \item \texttt{channels.msteams.teams..channels..requireMention}：每频道覆盖。
  \item \texttt{channels.msteams.teams..channels..tools}：每频道工具策略覆盖（\texttt{allow}/\texttt{deny}/\texttt{alsoAllow}）。
  \item \texttt{channels.msteams.teams..channels..toolsBySender}：每频道每发送者工具策略覆盖（支持 \texttt{"*"} 通配符）。
  \item \texttt{channels.msteams.sharePointSiteId}：用于群聊/频道文件上传的 SharePoint 站点 ID（见\href{\#sending-files-in-group-chats}{在群聊中发送文件}）。
\end{itemize}


\subsection{路由和会话}


\begin{itemize}
  \item 会话键遵循标准智能体格式（见 \href{/concepts/session}{/concepts/session}）：
  \item 私信共享主会话（\texttt{agent::}）。
  \item 频道/群组消息使用会话 ID：
  \item \texttt{agent::msteams:channel:}
  \item \texttt{agent::msteams:group:}
\end{itemize}


\subsection{回复样式：话题 vs 帖子}


Teams 最近在相同的底层数据模型上引入了两种频道 UI 样式：


\begin{table}[h]
\centering
\small
\begin{tabular}{|l|l|l|}
\hline
样式 & 描述 & 推荐的 \texttt{replyStyle} \\
\hline
\textbf{Posts}（经典） & 消息显示为卡片，下方有话题回复 & \texttt{thread}（默认） \\
\hline
\textbf{Threads}（类 Slack） & 消息线性流动，更像 Slack & \texttt{top-level} \\
\hline
\end{tabular}
\end{table}



\textbf{问题：} Teams API 不暴露频道使用的 UI 样式。如果你使用错误的 \texttt{replyStyle}：

\begin{itemize}
  \item 在 Threads 样式频道中使用 \texttt{thread} → 回复嵌套显示很别扭
  \item 在 Posts 样式频道中使用 \texttt{top-level} → 回复显示为单独的顶级帖子而不是在话题中
\end{itemize}


\textbf{解决方案：} 根据频道的设置方式为每个频道配置 \texttt{replyStyle}：

\begin{lstlisting}[language=json]
{
  "msteams": {
    "replyStyle": "thread",
    "teams": {
      "19:abc...@thread.tacv2": {
        "channels": {
          "19:xyz...@thread.tacv2": {
            "replyStyle": "top-level"
          }
        }
      }
    }
  }
}
\end{lstlisting}


\subsection{附件和图片}


\textbf{当前限制：}

\begin{itemize}
  \item \textbf{私信：} 图片和文件附件通过 Teams bot file API 工作。
  \item \textbf{频道/群组：} 附件存储在 M365 存储（SharePoint/OneDrive）中。webhook 负载仅包含 HTML 存根，而非实际文件字节。\textbf{需要 Graph API 权限}才能下载频道附件。
\end{itemize}


没有 Graph 权限，带图片的频道消息将作为纯文本接收（机器人无法访问图片内容）。
默认情况下，OpenClaw 仅从 Microsoft/Teams 主机名下载媒体。使用 \texttt{channels.msteams.mediaAllowHosts} 覆盖（使用 \texttt{["*"]} 允许任何主机）。
Authorization 头仅附加到 \texttt{channels.msteams.mediaAuthAllowHosts} 中的主机（默认为 Graph + Bot Framework 主机）。保持此列表严格（避免多租户后缀）。

\subsection{在群聊中发送文件}


机器人可以使用 FileConsentCard 流程在私信中发送文件（内置）。但是，\textbf{在群聊/频道中发送文件}需要额外设置：


\begin{table}[h]
\centering
\small
\begin{tabular}{|l|l|l|}
\hline
上下文 & 文件发送方式 & 所需设置 \\
\hline
\textbf{私信} & FileConsentCard → 用户接受 → 机器人上传 & 开箱即用 \\
\hline
\textbf{群聊/频道} & 上传到 SharePoint → 共享链接 & 需要 \texttt{sharePointSiteId} + Graph 权限 \\
\hline
\textbf{图片（任何上下文）} & Base64 编码内联 & 开箱即用 \\
\hline
\end{tabular}
\end{table}



\subsubsection{为什么群聊需要 SharePoint}


机器人没有个人 OneDrive 驱动器（\texttt{/me/drive} Graph API 端点对应用程序身份不起作用）。要在群聊/频道中发送文件，机器人上传到 \textbf{SharePoint 站点}并创建共享链接。

\subsubsection{设置}


\begin{enumerate}
  \item \textbf{在 Entra ID（Azure AD）→ App Registration 中添加 Graph API 权限}：
\end{enumerate}

\begin{itemize}
  \item \texttt{Sites.ReadWrite.All}（Application）- 上传文件到 SharePoint
  \item \texttt{Chat.Read.All}（Application）- 可选，启用每用户共享链接
\end{itemize}


\begin{enumerate}
  \item 为租户\textbf{授予管理员同意}。
\end{enumerate}


\begin{enumerate}
  \item \textbf{获取你的 SharePoint 站点 ID：}
\end{enumerate}


\begin{lstlisting}[language=bash]
   # 通过 Graph Explorer 或带有效令牌的 curl：
   curl -H "Authorization: Bearer $TOKEN" \
     "https://graph.microsoft.com/v1.0/sites/{hostname}:/{site-path}"

   # 示例：对于 "contoso.sharepoint.com/sites/BotFiles" 的站点
   curl -H "Authorization: Bearer $TOKEN" \
     "https://graph.microsoft.com/v1.0/sites/contoso.sharepoint.com:/sites/BotFiles"

   # 响应包含："id": "contoso.sharepoint.com,guid1,guid2"
\end{lstlisting}


\begin{enumerate}
  \item \textbf{配置 OpenClaw：}
\begin{lstlisting}[language=json]
   {
     channels: {
       msteams: {
         // ... 其他配置 ...
         sharePointSiteId: "contoso.sharepoint.com,guid1,guid2",
       },
     },
   }
\end{lstlisting}

\end{enumerate}


\subsubsection{共享行为}



\begin{table}[h]
\centering
\small
\begin{tabular}{|l|l|}
\hline
权限 & 共享行为 \\
\hline
仅 \texttt{Sites.ReadWrite.All} & 组织范围共享链接（组织中任何人都可以访问） \\
\hline
\texttt{Sites.ReadWrite.All} + \texttt{Chat.Read.All} & 每用户共享链接（仅聊天成员可以访问） \\
\hline
\end{tabular}
\end{table}



每用户共享更安全，因为只有聊天参与者才能访问文件。如果缺少 \texttt{Chat.Read.All} 权限，机器人回退到组织范围共享。

\subsubsection{回退行为}



\begin{table}[h]
\centering
\small
\begin{tabular}{|l|l|}
\hline
场景 & 结果 \\
\hline
群聊 + 文件 + 已配置 \texttt{sharePointSiteId} & 上传到 SharePoint，发送共享链接 \\
\hline
群聊 + 文件 + 无 \texttt{sharePointSiteId} & 尝试 OneDrive 上传（可能失败），仅发送文本 \\
\hline
个人聊天 + 文件 & FileConsentCard 流程（无需 SharePoint 即可工作） \\
\hline
任何上下文 + 图片 & Base64 编码内联（无需 SharePoint 即可工作） \\
\hline
\end{tabular}
\end{table}



\subsubsection{文件存储位置}


上传的文件存储在配置的 SharePoint 站点默认文档库中的 \texttt{/OpenClawShared/} 文件夹中。

\subsection{投票（Adaptive Cards）}


OpenClaw 将 Teams 投票作为 Adaptive Cards 发送（没有原生 Teams 投票 API）。

\begin{itemize}
  \item CLI：\texttt{openclaw message poll --channel msteams --target conversation: ...}
  \item 投票由 Gateway 网关记录在 \texttt{\textasciitilde{}/.openclaw/msteams-polls.json} 中。
  \item Gateway 网关必须保持在线才能记录投票。
  \item 投票尚不自动发布结果摘要（如需要请检查存储文件）。
\end{itemize}


\subsection{Adaptive Cards（任意）}


使用 \texttt{message} 工具或 CLI 向 Teams 用户或会话发送任意 Adaptive Card JSON。

\texttt{card} 参数接受 Adaptive Card JSON 对象。当提供 \texttt{card} 时，消息文本是可选的。

\textbf{智能体工具：}

\begin{lstlisting}[language=json]
{
  "action": "send",
  "channel": "msteams",
  "target": "user:<id>",
  "card": {
    "type": "AdaptiveCard",
    "version": "1.5",
    "body": [{ "type": "TextBlock", "text": "Hello!" }]
  }
}
\end{lstlisting}


\textbf{CLI：}

\begin{lstlisting}[language=bash]
openclaw message send --channel msteams \
  --target "conversation:19:abc...@thread.tacv2" \
  --card '{"type":"AdaptiveCard","version":"1.5","body":[{"type":"TextBlock","text":"Hello!"}]}'
\end{lstlisting}


参见 \href{https://adaptivecards.io/}{Adaptive Cards 文档}了解卡片模式和示例。目标格式详情见下方\href{\#target-formats}{目标格式}。

\subsection{目标格式}


MSTeams 目标使用前缀来区分用户和会话：


\begin{table}[h]
\centering
\small
\begin{tabular}{|l|l|l|}
\hline
目标类型 & 格式 & 示例 \\
\hline
用户（按 ID） & \texttt{user:} & \texttt{user:40a1a0ed-4ff2-4164-a219-55518990c197} \\
\hline
用户（按名称） & \texttt{user:} & \texttt{user:John Smith}（需要 Graph API） \\
\hline
群组/频道 & \texttt{conversation:} & \texttt{conversation:19:abc123...@thread.tacv2} \\
\hline
群组/频道（原始） & `` & \texttt{19:abc123...@thread.tacv2}（如果包含 \texttt{@thread}） \\
\hline
\end{tabular}
\end{table}



\textbf{CLI 示例：}

\begin{lstlisting}[language=bash]
# 按 ID 发送给用户
openclaw message send --channel msteams --target "user:40a1a0ed-..." --message "Hello"

# 按显示名称发送给用户（触发 Graph API 查找）
openclaw message send --channel msteams --target "user:John Smith" --message "Hello"

# 发送到群聊或频道
openclaw message send --channel msteams --target "conversation:19:abc...@thread.tacv2" --message "Hello"

# 向会话发送 Adaptive Card
openclaw message send --channel msteams --target "conversation:19:abc...@thread.tacv2" \
  --card '{"type":"AdaptiveCard","version":"1.5","body":[{"type":"TextBlock","text":"Hello"}]}'
\end{lstlisting}


\textbf{智能体工具示例：}

\begin{lstlisting}[language=json]
{
  "action": "send",
  "channel": "msteams",
  "target": "user:John Smith",
  "message": "Hello!"
}
\end{lstlisting}


\begin{lstlisting}[language=json]
{
  "action": "send",
  "channel": "msteams",
  "target": "conversation:19:abc...@thread.tacv2",
  "card": {
    "type": "AdaptiveCard",
    "version": "1.5",
    "body": [{ "type": "TextBlock", "text": "Hello" }]
  }
}
\end{lstlisting}


注意：没有 \texttt{user:} 前缀时，名称默认解析为群组/团队。按显示名称定位人员时始终使用 \texttt{user:}。

\subsection{主动消息}


\begin{itemize}
  \item 主动消息仅在用户交互\textbf{之后}才可能，因为我们在那时存储会话引用。
  \item 有关 \texttt{dmPolicy} 和允许列表控制，请参见 \texttt{/gateway/configuration}。
\end{itemize}


\subsection{团队和频道 ID（常见陷阱）}


Teams URL 中的 \texttt{groupId} 查询参数\textbf{不是}用于配置的团队 ID。请从 URL 路径中提取 ID：

\textbf{团队 URL：}

\begin{lstlisting}
https://teams.microsoft.com/l/team/19%3ABk4j...%40thread.tacv2/conversations?groupId=...
                                    └────────────────────────────┘
                                    团队 ID（URL 解码此部分）
\end{lstlisting}


\textbf{频道 URL：}

\begin{lstlisting}
https://teams.microsoft.com/l/channel/19%3A15bc...%40thread.tacv2/ChannelName?groupId=...
                                      └─────────────────────────┘
                                      频道 ID（URL 解码此部分）
\end{lstlisting}


\textbf{用于配置：}

\begin{itemize}
  \item 团队 ID = \texttt{/team/} 后的路径段（URL 解码，例如 \texttt{19:Bk4j...@thread.tacv2}）
  \item 频道 ID = \texttt{/channel/} 后的路径段（URL 解码）
  \item \textbf{忽略} \texttt{groupId} 查询参数
\end{itemize}


\subsection{私有频道}


机器人在私有频道中的支持有限：


\begin{table}[h]
\centering
\small
\begin{tabular}{|l|l|l|}
\hline
功能 & 标准频道 & 私有频道 \\
\hline
机器人安装 & 是 & 有限 \\
\hline
实时消息（webhook） & 是 & 可能不工作 \\
\hline
RSC 权限 & 是 & 行为可能不同 \\
\hline
@提及 & 是 & 如果机器人可访问 \\
\hline
Graph API 历史 & 是 & 是（有权限） \\
\hline
\end{tabular}
\end{table}



\textbf{如果私有频道不工作的变通方法：}

\begin{enumerate}
  \item 使用标准频道进行机器人交互
  \item 使用私信 - 用户始终可以直接给机器人发消息
  \item 使用 Graph API 进行历史访问（需要 \texttt{ChannelMessage.Read.All}）
\end{enumerate}


\subsection{故障排除}


\subsubsection{常见问题}


\begin{itemize}
  \item \textbf{频道中图片不显示：} 缺少 Graph 权限或管理员同意。重新安装 Teams 应用并完全退出/重新打开 Teams。
  \item \textbf{频道中无响应：} 默认需要提及；设置 \texttt{channels.msteams.requireMention=false} 或按团队/频道配置。
  \item \textbf{版本不匹配（Teams 仍显示旧清单）：} 移除 + 重新添加应用并完全退出 Teams 以刷新。
  \item \textbf{来自 webhook 的 401 Unauthorized：} 在没有 Azure JWT 的情况下手动测试时属于预期情况 - 意味着端点可达但认证失败。使用 Azure Web Chat 正确测试。
\end{itemize}


\subsubsection{清单上传错误}


\begin{itemize}
  \item \textbf{"Icon file cannot be empty"：} 清单引用的图标文件为 0 字节。创建有效的 PNG 图标（\texttt{outline.png} 为 32x32，\texttt{color.png} 为 192x192）。
  \item \textbf{"webApplicationInfo.Id already in use"：} 应用仍安装在另一个团队/聊天中。先找到并卸载它，或等待 5-10 分钟让其传播。
  \item \textbf{上传时"Something went wrong"：} 改为通过 https://admin.teams.microsoft.com 上传，打开浏览器 DevTools（F12）→ Network 选项卡，检查响应正文中的实际错误。
  \item \textbf{侧载失败：} 尝试"Upload an app to your org's app catalog"而不是"Upload a custom app" - 这通常可以绕过侧载限制。
\end{itemize}


\subsubsection{RSC 权限不工作}


\begin{enumerate}
  \item 验证 \texttt{webApplicationInfo.id} 与你的机器人 App ID 完全匹配
  \item 重新上传应用并在团队/聊天中重新安装
  \item 检查你的组织管理员是否阻止了 RSC 权限
  \item 确认你使用的是正确的范围：团队使用 \texttt{ChannelMessage.Read.Group}，群聊使用 \texttt{ChatMessage.Read.Chat}
\end{enumerate}


\subsection{参考资料}


\begin{itemize}
  \item \href{https://learn.microsoft.com/en-us/azure/bot-service/bot-service-quickstart-registration}{创建 Azure Bot} - Azure Bot 设置指南
  \item \href{https://dev.teams.microsoft.com/apps}{Teams 开发者门户} - 创建/管理 Teams 应用
  \item \href{https://learn.microsoft.com/en-us/microsoftteams/platform/resources/schema/manifest-schema}{Teams 应用清单模式}
  \item \href{https://learn.microsoft.com/en-us/microsoftteams/platform/bots/how-to/conversations/channel-messages-with-rsc}{使用 RSC 接收频道消息}
  \item \href{https://learn.microsoft.com/en-us/microsoftteams/platform/graph-api/rsc/resource-specific-consent}{RSC 权限参考}
  \item \href{https://learn.microsoft.com/en-us/microsoftteams/platform/bots/how-to/bots-filesv4}{Teams 机器人文件处理}（频道/群组需要 Graph）
  \item \href{https://learn.microsoft.com/en-us/microsoftteams/platform/bots/how-to/conversations/send-proactive-messages}{主动消息}
\end{itemize}



\clearpage

% === 源文件: channels/nextcloud-talk.md ===
\section{Nextcloud Talk（插件）}


状态：通过插件支持（webhook 机器人）。支持私信、房间、表情回应和 Markdown 消息。

\subsection{需要插件}


Nextcloud Talk 以插件形式提供，不包含在核心安装包中。

通过 CLI 安装（npm 仓库）：

\begin{lstlisting}[language=bash]
openclaw plugins install @openclaw/nextcloud-talk
\end{lstlisting}


本地检出安装（从 git 仓库运行时）：

\begin{lstlisting}[language=bash]
openclaw plugins install ./extensions/nextcloud-talk
\end{lstlisting}


如果你在配置/新手引导过程中选择了 Nextcloud Talk，并且检测到 git 检出，
OpenClaw 将自动提供本地安装路径。

详情：\href{/tools/plugin}{插件}

\subsection{快速设置（新手）}


\begin{enumerate}
  \item 安装 Nextcloud Talk 插件。
  \item 在你的 Nextcloud 服务器上创建机器人：
\begin{lstlisting}[language=bash]
   ./occ talk:bot:install "OpenClaw" "<shared-secret>" "<webhook-url>" --feature reaction
\end{lstlisting}

  \item 在目标房间设置中启用机器人。
  \item 配置 OpenClaw：
\end{enumerate}

\begin{itemize}
  \item 配置项：\texttt{channels.nextcloud-talk.baseUrl} + \texttt{channels.nextcloud-talk.botSecret}
  \item 或环境变量：\texttt{NEXTCLOUD\_TALK\_BOT\_SECRET}（仅默认账户）
\end{itemize}

\begin{enumerate}
  \item 重启 Gateway 网关（或完成新手引导）。
\end{enumerate}


最小配置：

\begin{lstlisting}[language=json]
{
  channels: {
    "nextcloud-talk": {
      enabled: true,
      baseUrl: "https://cloud.example.com",
      botSecret: "shared-secret",
      dmPolicy: "pairing",
    },
  },
}
\end{lstlisting}


\subsection{注意事项}


\begin{itemize}
  \item 机器人无法主动发起私信。用户必须先向机器人发送消息。
  \item Webhook URL 必须可被 Gateway 网关访问；如果在代理后面，请设置 \texttt{webhookPublicUrl}。
  \item 机器人 API 不支持媒体上传；媒体以 URL 形式发送。
  \item Webhook 载荷无法区分私信和房间；设置 \texttt{apiUser} + \texttt{apiPassword} 以启用房间类型查询（否则私信将被视为房间）。
\end{itemize}


\subsection{访问控制（私信）}


\begin{itemize}
  \item 默认：\texttt{channels.nextcloud-talk.dmPolicy = "pairing"}。未知发送者将收到配对码。
  \item 批准方式：
  \item \texttt{openclaw pairing list nextcloud-talk}
  \item \texttt{openclaw pairing approve nextcloud-talk }
  \item 公开私信：\texttt{channels.nextcloud-talk.dmPolicy="open"} 加上 \texttt{channels.nextcloud-talk.allowFrom=["*"]}。
\end{itemize}


\subsection{房间（群组）}


\begin{itemize}
  \item 默认：\texttt{channels.nextcloud-talk.groupPolicy = "allowlist"}（需要提及触发）。
  \item 使用 \texttt{channels.nextcloud-talk.rooms} 设置房间白名单：
\end{itemize}


\begin{lstlisting}[language=json]
{
  channels: {
    "nextcloud-talk": {
      rooms: {
        "room-token": { requireMention: true },
      },
    },
  },
}
\end{lstlisting}


\begin{itemize}
  \item 如需禁止所有房间，保持白名单为空或设置 \texttt{channels.nextcloud-talk.groupPolicy="disabled"}。
\end{itemize}


\subsection{功能支持}



\begin{table}[h]
\centering
\small
\begin{tabular}{|l|l|}
\hline
功能 & 状态 \\
\hline
私信 & 支持 \\
\hline
房间 & 支持 \\
\hline
话题 & 不支持 \\
\hline
媒体 & 仅 URL \\
\hline
表情回应 & 支持 \\
\hline
原生命令 & 不支持 \\
\hline
\end{tabular}
\end{table}



\subsection{配置参考（Nextcloud Talk）}


完整配置：\href{/gateway/configuration}{配置}

提供商选项：

\begin{itemize}
  \item \texttt{channels.nextcloud-talk.enabled}：启用/禁用渠道启动。
  \item \texttt{channels.nextcloud-talk.baseUrl}：Nextcloud 实例 URL。
  \item \texttt{channels.nextcloud-talk.botSecret}：机器人共享密钥。
  \item \texttt{channels.nextcloud-talk.botSecretFile}：密钥文件路径。
  \item \texttt{channels.nextcloud-talk.apiUser}：用于房间查询的 API 用户（私信检测）。
  \item \texttt{channels.nextcloud-talk.apiPassword}：用于房间查询的 API/应用密码。
  \item \texttt{channels.nextcloud-talk.apiPasswordFile}：API 密码文件路径。
  \item \texttt{channels.nextcloud-talk.webhookPort}：webhook 监听端口（默认：8788）。
  \item \texttt{channels.nextcloud-talk.webhookHost}：webhook 主机（默认：0.0.0.0）。
  \item \texttt{channels.nextcloud-talk.webhookPath}：webhook 路径（默认：/nextcloud-talk-webhook）。
  \item \texttt{channels.nextcloud-talk.webhookPublicUrl}：外部可达的 webhook URL。
  \item \texttt{channels.nextcloud-talk.dmPolicy}：\texttt{pairing | allowlist | open | disabled}。
  \item \texttt{channels.nextcloud-talk.allowFrom}：私信白名单（用户 ID）。\texttt{open} 需要 \texttt{"*"}。
  \item \texttt{channels.nextcloud-talk.groupPolicy}：\texttt{allowlist | open | disabled}。
  \item \texttt{channels.nextcloud-talk.groupAllowFrom}：群组白名单（用户 ID）。
  \item \texttt{channels.nextcloud-talk.rooms}：每个房间的设置和白名单。
  \item \texttt{channels.nextcloud-talk.historyLimit}：群组历史记录限制（0 表示禁用）。
  \item \texttt{channels.nextcloud-talk.dmHistoryLimit}：私信历史记录限制（0 表示禁用）。
  \item \texttt{channels.nextcloud-talk.dms}：每个私信的覆盖设置（historyLimit）。
  \item \texttt{channels.nextcloud-talk.textChunkLimit}：出站文本分块大小（字符数）。
  \item \texttt{channels.nextcloud-talk.chunkMode}：\texttt{length}（默认）或 \texttt{newline}，在长度分块前按空行（段落边界）分割。
  \item \texttt{channels.nextcloud-talk.blockStreaming}：禁用此渠道的分块流式传输。
  \item \texttt{channels.nextcloud-talk.blockStreamingCoalesce}：分块流式传输合并调优。
  \item \texttt{channels.nextcloud-talk.mediaMaxMb}：入站媒体大小上限（MB）。
\end{itemize}



\clearpage

% === 源文件: channels/nostr.md ===
\section{Nostr}


\textbf{状态：} 可选插件（默认禁用）。

Nostr 是一个去中心化的社交网络协议。此渠道使 OpenClaw 能够通过 NIP-04 接收和回复加密私信（DMs）。

\subsection{安装（按需）}


\subsubsection{新手引导（推荐）}


\begin{itemize}
  \item 新手引导向导（\texttt{openclaw onboard}）和 \texttt{openclaw channels add} 会列出可选的渠道插件。
  \item 选择 Nostr 会提示你按需安装插件。
\end{itemize}


安装默认值：

\begin{itemize}
  \item \textbf{Dev 渠道 + git checkout 可用：} 使用本地插件路径。
  \item \textbf{Stable/Beta：} 从 npm 下载。
\end{itemize}


你可以随时在提示中覆盖选择。

\subsubsection{手动安装}


\begin{lstlisting}[language=bash]
openclaw plugins install @openclaw/nostr
\end{lstlisting}


使用本地 checkout（开发工作流）：

\begin{lstlisting}[language=bash]
openclaw plugins install --link <path-to-openclaw>/extensions/nostr
\end{lstlisting}


安装或启用插件后重启 Gateway 网关。

\subsection{快速设置}


\begin{enumerate}
  \item 生成 Nostr 密钥对（如需要）：
\end{enumerate}


\begin{lstlisting}[language=bash]
# 使用 nak
nak key generate
\end{lstlisting}


\begin{enumerate}
  \item 添加到配置：
\end{enumerate}


\begin{lstlisting}[language=json]
{
  "channels": {
    "nostr": {
      "privateKey": "${NOSTR_PRIVATE_KEY}"
    }
  }
}
\end{lstlisting}


\begin{enumerate}
  \item 导出密钥：
\end{enumerate}


\begin{lstlisting}[language=bash]
export NOSTR_PRIVATE_KEY="nsec1..."
\end{lstlisting}


\begin{enumerate}
  \item 重启 Gateway 网关。
\end{enumerate}


\subsection{配置参考}



\begin{table}[h]
\centering
\small
\begin{tabular}{|l|l|l|l|}
\hline
键 & 类型 & 默认值 & 描述 \\
\hline
\texttt{privateKey} & string & 必填 & \texttt{nsec} 或十六进制格式的私钥 \\
\hline
\texttt{relays} & string[] & \texttt{['wss://relay.damus.io', 'wss://nos.lol']} & 中继 URL（WebSocket） \\
\hline
\texttt{dmPolicy} & string & \texttt{pairing} & 私信访问策略 \\
\hline
\texttt{allowFrom} & string[] & \texttt{[]} & 允许的发送者公钥 \\
\hline
\texttt{enabled} & boolean & \texttt{true} & 启用/禁用渠道 \\
\hline
\texttt{name} & string & - & 显示名称 \\
\hline
\texttt{profile} & object & - & NIP-01 个人资料元数据 \\
\hline
\end{tabular}
\end{table}



\subsection{个人资料元数据}


个人资料数据作为 NIP-01 \texttt{kind:0} 事件发布。你可以从控制界面（Channels -> Nostr -> Profile）管理它，或直接在配置中设置。

示例：

\begin{lstlisting}[language=json]
{
  "channels": {
    "nostr": {
      "privateKey": "${NOSTR_PRIVATE_KEY}",
      "profile": {
        "name": "openclaw",
        "displayName": "OpenClaw",
        "about": "Personal assistant DM bot",
        "picture": "https://example.com/avatar.png",
        "banner": "https://example.com/banner.png",
        "website": "https://example.com",
        "nip05": "openclaw@example.com",
        "lud16": "openclaw@example.com"
      }
    }
  }
}
\end{lstlisting}


注意事项：

\begin{itemize}
  \item 个人资料 URL 必须使用 \texttt{https://}。
  \item 从中继导入会合并字段并保留本地覆盖。
\end{itemize}


\subsection{访问控制}


\subsubsection{私信策略}


\begin{itemize}
  \item \textbf{pairing}（默认）：未知发送者会收到配对码。
  \item \textbf{allowlist}：只有 \texttt{allowFrom} 中的公钥可以发送私信。
  \item \textbf{open}：公开接收私信（需要 \texttt{allowFrom: ["*"]}）。
  \item \textbf{disabled}：忽略接收的私信。
\end{itemize}


\subsubsection{允许列表示例}


\begin{lstlisting}[language=json]
{
  "channels": {
    "nostr": {
      "privateKey": "${NOSTR_PRIVATE_KEY}",
      "dmPolicy": "allowlist",
      "allowFrom": ["npub1abc...", "npub1xyz..."]
    }
  }
}
\end{lstlisting}


\subsection{密钥格式}


接受的格式：

\begin{itemize}
  \item \textbf{私钥：} \texttt{nsec...} 或 64 字符十六进制
  \item \textbf{公钥（\texttt{allowFrom}）：} \texttt{npub...} 或十六进制
\end{itemize}


\subsection{中继}


默认值：\texttt{relay.damus.io} 和 \texttt{nos.lol}。

\begin{lstlisting}[language=json]
{
  "channels": {
    "nostr": {
      "privateKey": "${NOSTR_PRIVATE_KEY}",
      "relays": ["wss://relay.damus.io", "wss://relay.primal.net", "wss://nostr.wine"]
    }
  }
}
\end{lstlisting}


提示：

\begin{itemize}
  \item 使用 2-3 个中继以实现冗余。
  \item 避免使用过多中继（延迟、重复）。
  \item 付费中继可以提高可靠性。
  \item 本地中继适合测试（\texttt{ws://localhost:7777}）。
\end{itemize}


\subsection{协议支持}



\begin{table}[h]
\centering
\small
\begin{tabular}{|l|l|l|}
\hline
NIP & 状态 & 描述 \\
\hline
NIP-01 & 已支持 & 基本事件格式 + 个人资料元数据 \\
\hline
NIP-04 & 已支持 & 加密私信（\texttt{kind:4}） \\
\hline
NIP-17 & 计划中 & 礼物包装私信 \\
\hline
NIP-44 & 计划中 & 版本化加密 \\
\hline
\end{tabular}
\end{table}



\subsection{测试}


\subsubsection{本地中继}


\begin{lstlisting}[language=bash]
# 启动 strfry
docker run -p 7777:7777 ghcr.io/hoytech/strfry
\end{lstlisting}


\begin{lstlisting}[language=json]
{
  "channels": {
    "nostr": {
      "privateKey": "${NOSTR_PRIVATE_KEY}",
      "relays": ["ws://localhost:7777"]
    }
  }
}
\end{lstlisting}


\subsubsection{手动测试}


\begin{enumerate}
  \item 从日志中记下机器人公钥（npub）。
  \item 打开 Nostr 客户端（Damus、Amethyst 等）。
  \item 向机器人公钥发送私信。
  \item 验证响应。
\end{enumerate}


\subsection{故障排除}


\subsubsection{未收到消息}


\begin{itemize}
  \item 验证私钥是否有效。
  \item 确保中继 URL 可访问并使用 \texttt{wss://}（本地使用 \texttt{ws://}）。
  \item 确认 \texttt{enabled} 不是 \texttt{false}。
  \item 检查 Gateway 网关日志中的中继连接错误。
\end{itemize}


\subsubsection{未发送响应}


\begin{itemize}
  \item 检查中继是否接受写入。
  \item 验证出站连接。
  \item 注意中继速率限制。
\end{itemize}


\subsubsection{重复响应}


\begin{itemize}
  \item 使用多个中继时属于正常现象。
  \item 消息按事件 ID 去重；只有首次投递会触发响应。
\end{itemize}


\subsection{安全}


\begin{itemize}
  \item 切勿提交私钥。
  \item 使用环境变量存储密钥。
  \item 生产环境机器人考虑使用 \texttt{allowlist}。
\end{itemize}


\subsection{限制（MVP）}


\begin{itemize}
  \item 仅支持私信（不支持群聊）。
  \item 不支持媒体附件。
  \item 仅支持 NIP-04（计划支持 NIP-17 礼物包装）。
\end{itemize}



\clearpage

% === 源文件: channels/pairing.md ===
\section{配对}


"配对"是 OpenClaw 的显式\textbf{所有者批准}步骤。它用于两个地方：

\begin{enumerate}
  \item \textbf{私信配对}（谁被允许与机器人对话）
  \item \textbf{节点配对}（哪些设备/节点被允许加入 Gateway 网关网络）
\end{enumerate}


安全上下文：\href{/gateway/security}{安全}

\subsection{1）私信配对（入站聊天访问）}


当渠道配置为私信策略 \texttt{pairing} 时，未知发送者会收到一个短代码，他们的消息\textbf{不会被处理}，直到你批准。

默认私信策略记录在：\href{/gateway/security}{安全}

配对代码：

\begin{itemize}
  \item 8 个字符，大写，无歧义字符（\texttt{0O1I}）。
  \item \textbf{1 小时后过期}。机器人仅在创建新请求时发送配对消息（大约每个发送者每小时一次）。
  \item 待处理的私信配对请求默认上限为\textbf{每个渠道 3 个}；在一个过期或被批准之前，额外的请求将被忽略。
\end{itemize}


\subsubsection{批准发送者}


\begin{lstlisting}[language=bash]
openclaw pairing list telegram
openclaw pairing approve telegram <CODE>
\end{lstlisting}


支持的渠道：\texttt{telegram}、\texttt{whatsapp}、\texttt{signal}、\texttt{imessage}、\texttt{discord}、\texttt{slack}。

\subsubsection{状态存储位置}


存储在 \texttt{\textasciitilde{}/.openclaw/credentials/} 下：

\begin{itemize}
  \item 待处理请求：\texttt{-pairing.json}
  \item 已批准允许列表存储：\texttt{-allowFrom.json}
\end{itemize}


将这些视为敏感信息（它们控制对你助手的访问）。

\subsection{2）节点设备配对（iOS/Android/macOS/无头节点）}


节点作为 \texttt{role: node} 的\textbf{设备}连接到 Gateway 网关。Gateway 网关创建一个必须被批准的设备配对请求。

\subsubsection{批准节点设备}


\begin{lstlisting}[language=bash]
openclaw devices list
openclaw devices approve <requestId>
openclaw devices reject <requestId>
\end{lstlisting}


\subsubsection{状态存储位置}


存储在 \texttt{\textasciitilde{}/.openclaw/devices/} 下：

\begin{itemize}
  \item \texttt{pending.json}（短期；待处理请求会过期）
  \item \texttt{paired.json}（已配对设备 + 令牌）
\end{itemize}


\subsubsection{说明}


\begin{itemize}
  \item 旧版 \texttt{node.pair.*} API（CLI：\texttt{openclaw nodes pending/approve}）是一个单独的 Gateway 网关拥有的配对存储。WS 节点仍然需要设备配对。
\end{itemize}


\subsection{相关文档}


\begin{itemize}
  \item 安全模型 + 提示注入：\href{/gateway/security}{安全}
  \item 安全更新（运行 doctor）：\href{/install/updating}{更新}
  \item 渠道配置：
  \item Telegram：\href{/channels/telegram}{Telegram}
  \item WhatsApp：\href{/channels/whatsapp}{WhatsApp}
  \item Signal：\href{/channels/signal}{Signal}
  \item iMessage：\href{/channels/imessage}{iMessage}
  \item Discord：\href{/channels/discord}{Discord}
  \item Slack：\href{/channels/slack}{Slack}
\end{itemize}



\clearpage

% === 源文件: channels/signal.md ===
\section{Signal (signal-cli)}


状态：外部 CLI 集成。Gateway 网关通过 HTTP JSON-RPC + SSE 与 \texttt{signal-cli} 通信。

\subsection{快速设置（初学者）}


\begin{enumerate}
  \item 为 bot 使用\textbf{单独的 Signal 号码}（推荐）。
  \item 安装 \texttt{signal-cli}（需要 Java）。
  \item 链接 bot 设备并启动守护进程：
\end{enumerate}

\begin{itemize}
  \item \texttt{signal-cli link -n "OpenClaw"}
\end{itemize}

\begin{enumerate}
  \item 配置 OpenClaw 并启动 Gateway 网关。
\end{enumerate}


最小配置：

\begin{lstlisting}[language=json]
{
  channels: {
    signal: {
      enabled: true,
      account: "+15551234567",
      cliPath: "signal-cli",
      dmPolicy: "pairing",
      allowFrom: ["+15557654321"],
    },
  },
}
\end{lstlisting}


\subsection{它是什么}


\begin{itemize}
  \item 通过 \texttt{signal-cli} 的 Signal 渠道（非嵌入式 libsignal）。
  \item 确定性路由：回复始终返回到 Signal。
  \item 私信共享智能体的主会话；群组是隔离的（\texttt{agent::signal:group:}）。
\end{itemize}


\subsection{配置写入}


默认情况下，Signal 允许写入由 \texttt{/config set|unset} 触发的配置更新（需要 \texttt{commands.config: true}）。

禁用方式：

\begin{lstlisting}[language=json]
{
  channels: { signal: { configWrites: false } },
}
\end{lstlisting}


\subsection{号码模型（重要）}


\begin{itemize}
  \item Gateway 网关连接到一个 \textbf{Signal 设备}（\texttt{signal-cli} 账户）。
  \item 如果你在\textbf{个人 Signal 账户}上运行 bot，它会忽略你自己的消息（循环保护）。
  \item 要实现"我发消息给 bot 然后它回复"，请使用\textbf{单独的 bot 号码}。
\end{itemize}


\subsection{设置（快速路径）}


\begin{enumerate}
  \item 安装 \texttt{signal-cli}（需要 Java）。
  \item 链接 bot 账户：
\end{enumerate}

\begin{itemize}
  \item \texttt{signal-cli link -n "OpenClaw"} 然后在 Signal 中扫描二维码。
\end{itemize}

\begin{enumerate}
  \item 配置 Signal 并启动 Gateway 网关。
\end{enumerate}


示例：

\begin{lstlisting}[language=json]
{
  channels: {
    signal: {
      enabled: true,
      account: "+15551234567",
      cliPath: "signal-cli",
      dmPolicy: "pairing",
      allowFrom: ["+15557654321"],
    },
  },
}
\end{lstlisting}


多账户支持：使用 \texttt{channels.signal.accounts} 配置每个账户及可选的 \texttt{name}。共享模式请参见 \href{/gateway/configuration\#telegramaccounts--discordaccounts--slackaccounts--signalaccounts--imessageaccounts}{\texttt{gateway/configuration}}。

\subsection{外部守护进程模式（httpUrl）}


如果你想自己管理 \texttt{signal-cli}（JVM 冷启动慢、容器初始化或共享 CPU），请单独运行守护进程并将 OpenClaw 指向它：

\begin{lstlisting}[language=json]
{
  channels: {
    signal: {
      httpUrl: "http://127.0.0.1:8080",
      autoStart: false,
    },
  },
}
\end{lstlisting}


这会跳过自动启动和 OpenClaw 内部的启动等待。对于自动启动时的慢启动，请设置 \texttt{channels.signal.startupTimeoutMs}。

\subsection{访问控制（私信 + 群组）}


私信：

\begin{itemize}
  \item 默认：\texttt{channels.signal.dmPolicy = "pairing"}。
  \item 未知发送者会收到配对码；消息在批准前会被忽略（配对码 1 小时后过期）。
  \item 通过以下方式批准：
  \item \texttt{openclaw pairing list signal}
  \item \texttt{openclaw pairing approve signal }
  \item 配对是 Signal 私信的默认令牌交换方式。详情：\href{/channels/pairing}{配对}
  \item 仅有 UUID 的发送者（来自 \texttt{sourceUuid}）在 \texttt{channels.signal.allowFrom} 中存储为 \texttt{uuid:}。
\end{itemize}


群组：

\begin{itemize}
  \item \texttt{channels.signal.groupPolicy = open | allowlist | disabled}。
  \item 当设置为 \texttt{allowlist} 时，\texttt{channels.signal.groupAllowFrom} 控制谁可以在群组中触发。
\end{itemize}


\subsection{工作原理（行为）}


\begin{itemize}
  \item \texttt{signal-cli} 作为守护进程运行；Gateway 网关通过 SSE 读取事件。
  \item 入站消息被规范化为共享渠道信封。
  \item 回复始终路由回同一号码或群组。
\end{itemize}


\subsection{媒体 + 限制}


\begin{itemize}
  \item 出站文本按 \texttt{channels.signal.textChunkLimit} 分块（默认 4000）。
  \item 可选换行分块：设置 \texttt{channels.signal.chunkMode="newline"} 在长度分块前按空行（段落边界）分割。
  \item 支持附件（从 \texttt{signal-cli} 获取 base64）。
  \item 默认媒体上限：\texttt{channels.signal.mediaMaxMb}（默认 8）。
  \item 使用 \texttt{channels.signal.ignoreAttachments} 跳过下载媒体。
  \item 群组历史上下文使用 \texttt{channels.signal.historyLimit}（或 \texttt{channels.signal.accounts.*.historyLimit}），回退到 \texttt{messages.groupChat.historyLimit}。设置 \texttt{0} 禁用（默认 50）。
\end{itemize}


\subsection{输入指示器 + 已读回执}


\begin{itemize}
  \item \textbf{输入指示器}：OpenClaw 通过 \texttt{signal-cli sendTyping} 发送输入信号，并在回复运行时刷新它们。
  \item \textbf{已读回执}：当 \texttt{channels.signal.sendReadReceipts} 为 true 时，OpenClaw 为允许的私信转发已读回执。
  \item Signal-cli 不暴露群组的已读回执。
\end{itemize}


\subsection{表情回应（message 工具）}


\begin{itemize}
  \item 使用 \texttt{message action=react} 配合 \texttt{channel=signal}。
  \item 目标：发送者 E.164 或 UUID（使用配对输出中的 \texttt{uuid:}；裸 UUID 也可以）。
  \item \texttt{messageId} 是你要回应的消息的 Signal 时间戳。
  \item 群组表情回应需要 \texttt{targetAuthor} 或 \texttt{targetAuthorUuid}。
\end{itemize}


示例：

\begin{lstlisting}
message action=react channel=signal target=uuid:123e4567-e89b-12d3-a456-426614174000 messageId=1737630212345 emoji=🔥
message action=react channel=signal target=+15551234567 messageId=1737630212345 emoji=🔥 remove=true
message action=react channel=signal target=signal:group:<groupId> targetAuthor=uuid:<sender-uuid> messageId=1737630212345 emoji=✅
\end{lstlisting}


配置：

\begin{itemize}
  \item \texttt{channels.signal.actions.reactions}：启用/禁用表情回应操作（默认 true）。
  \item \texttt{channels.signal.reactionLevel}：\texttt{off | ack | minimal | extensive}。
  \item \texttt{off}/\texttt{ack} 禁用智能体表情回应（message 工具 \texttt{react} 会报错）。
  \item \texttt{minimal}/\texttt{extensive} 启用智能体表情回应并设置指导级别。
  \item 每账户覆盖：\texttt{channels.signal.accounts..actions.reactions}、\texttt{channels.signal.accounts..reactionLevel}。
\end{itemize}


\subsection{投递目标（CLI/cron）}


\begin{itemize}
  \item 私信：\texttt{signal:+15551234567}（或纯 E.164）。
  \item UUID 私信：\texttt{uuid:}（或裸 UUID）。
  \item 群组：\texttt{signal:group:}。
  \item 用户名：\texttt{username:}（如果你的 Signal 账户支持）。
\end{itemize}


\subsection{配置参考（Signal）}


完整配置：\href{/gateway/configuration}{配置}

提供商选项：

\begin{itemize}
  \item \texttt{channels.signal.enabled}：启用/禁用渠道启动。
  \item \texttt{channels.signal.account}：bot 账户的 E.164。
  \item \texttt{channels.signal.cliPath}：\texttt{signal-cli} 的路径。
  \item \texttt{channels.signal.httpUrl}：完整守护进程 URL（覆盖 host/port）。
  \item \texttt{channels.signal.httpHost}、\texttt{channels.signal.httpPort}：守护进程绑定（默认 127.0.0.1:8080）。
  \item \texttt{channels.signal.autoStart}：自动启动守护进程（如果未设置 \texttt{httpUrl} 则默认 true）。
  \item \texttt{channels.signal.startupTimeoutMs}：启动等待超时（毫秒）（上限 120000）。
  \item \texttt{channels.signal.receiveMode}：\texttt{on-start | manual}。
  \item \texttt{channels.signal.ignoreAttachments}：跳过附件下载。
  \item \texttt{channels.signal.ignoreStories}：忽略来自守护进程的动态。
  \item \texttt{channels.signal.sendReadReceipts}：转发已读回执。
  \item \texttt{channels.signal.dmPolicy}：\texttt{pairing | allowlist | open | disabled}（默认：pairing）。
  \item \texttt{channels.signal.allowFrom}：私信允许列表（E.164 或 \texttt{uuid:}）。\texttt{open} 需要 \texttt{"*"}。Signal 没有用户名；使用电话/UUID id。
  \item \texttt{channels.signal.groupPolicy}：\texttt{open | allowlist | disabled}（默认：allowlist）。
  \item \texttt{channels.signal.groupAllowFrom}：群组发送者允许列表。
  \item \texttt{channels.signal.historyLimit}：作为上下文包含的最大群组消息数（0 禁用）。
  \item \texttt{channels.signal.dmHistoryLimit}：私信历史限制（用户轮次）。每用户覆盖：\texttt{channels.signal.dms[""].historyLimit}。
  \item \texttt{channels.signal.textChunkLimit}：出站分块大小（字符）。
  \item \texttt{channels.signal.chunkMode}：\texttt{length}（默认）或 \texttt{newline} 在长度分块前按空行（段落边界）分割。
  \item \texttt{channels.signal.mediaMaxMb}：入站/出站媒体上限（MB）。
\end{itemize}


相关全局选项：

\begin{itemize}
  \item \texttt{agents.list[].groupChat.mentionPatterns}（Signal 不支持原生提及）。
  \item \texttt{messages.groupChat.mentionPatterns}（全局回退）。
  \item \texttt{messages.responsePrefix}。
\end{itemize}



\clearpage

% === 源文件: channels/slack.md ===
\section{Slack}


\subsection{Socket 模式（默认）}


\subsubsection{快速设置（新手）}


\begin{enumerate}
  \item 创建一个 Slack 应用并启用 \textbf{Socket Mode}。
  \item 创建一个 \textbf{App Token}（\texttt{xapp-...}）和 \textbf{Bot Token}（\texttt{xoxb-...}）。
  \item 为 OpenClaw 设置令牌并启动 Gateway 网关。
\end{enumerate}


最小配置：

\begin{lstlisting}[language=json]
{
  channels: {
    slack: {
      enabled: true,
      appToken: "xapp-...",
      botToken: "xoxb-...",
    },
  },
}
\end{lstlisting}


\subsubsection{设置}


\begin{enumerate}
  \item 在 https://api.slack.com/apps 创建一个 Slack 应用（从头开始）。
  \item \textbf{Socket Mode} → 开启。然后前往 \textbf{Basic Information} → \textbf{App-Level Tokens} → \textbf{Generate Token and Scopes}，添加 \texttt{connections:write} 权限范围。复制 \textbf{App Token}（\texttt{xapp-...}）。
  \item \textbf{OAuth \& Permissions} → 添加 bot token 权限范围（使用下面的 manifest）。点击 \textbf{Install to Workspace}。复制 \textbf{Bot User OAuth Token}（\texttt{xoxb-...}）。
  \item 可选：\textbf{OAuth \& Permissions} → 添加 \textbf{User Token Scopes}（参见下面的只读列表）。重新安装应用并复制 \textbf{User OAuth Token}（\texttt{xoxp-...}）。
  \item \textbf{Event Subscriptions} → 启用事件并订阅：
\end{enumerate}

\begin{itemize}
  \item \texttt{message.*}（包括编辑/删除/线程广播）
  \item \texttt{app\_mention}
  \item \texttt{reaction\_added}、\texttt{reaction\_removed}
  \item \texttt{member\_joined\_channel}、\texttt{member\_left\_channel}
  \item \texttt{channel\_rename}
  \item \texttt{pin\_added}、\texttt{pin\_removed}
\end{itemize}

\begin{enumerate}
  \item 邀请机器人加入你希望它读取的频道。
  \item Slash Commands → 如果你使用 \texttt{channels.slack.slashCommand}，创建 \texttt{/openclaw}。如果启用原生命令，为每个内置命令添加一个斜杠命令（名称与 \texttt{/help} 相同）。除非你设置 \texttt{channels.slack.commands.native: true}，否则 Slack 默认关闭原生命令（全局 \texttt{commands.native} 是 \texttt{"auto"}，对 Slack 保持关闭）。
  \item App Home → 启用 \textbf{Messages Tab} 以便用户可以私信机器人。
\end{enumerate}


使用下面的 manifest 以保持权限范围和事件同步。

多账户支持：使用 \texttt{channels.slack.accounts} 配置每个账户的令牌和可选的 \texttt{name}。参见 \href{/gateway/configuration\#telegramaccounts--discordaccounts--slackaccounts--signalaccounts--imessageaccounts}{\texttt{gateway/configuration}} 了解共享模式。

\subsubsection{OpenClaw 配置（最小）}


通过环境变量设置令牌（推荐）：

\begin{itemize}
  \item \texttt{SLACK\_APP\_TOKEN=xapp-...}
  \item \texttt{SLACK\_BOT\_TOKEN=xoxb-...}
\end{itemize}


或通过配置：

\begin{lstlisting}[language=json]
{
  channels: {
    slack: {
      enabled: true,
      appToken: "xapp-...",
      botToken: "xoxb-...",
    },
  },
}
\end{lstlisting}


\subsubsection{用户令牌（可选）}


OpenClaw 可以使用 Slack 用户令牌（\texttt{xoxp-...}）进行读取操作（历史记录、置顶、表情回应、表情符号、成员信息）。默认情况下保持只读：当存在用户令牌时，读取优先使用用户令牌，而写入仍然使用 bot 令牌，除非你明确选择加入。即使设置了 \texttt{userTokenReadOnly: false}，当 bot 令牌可用时，写入仍然优先使用 bot 令牌。

用户令牌在配置文件中配置（不支持环境变量）。对于多账户，设置 \texttt{channels.slack.accounts..userToken}。

包含 bot + app + 用户令牌的示例：

\begin{lstlisting}[language=json]
{
  channels: {
    slack: {
      enabled: true,
      appToken: "xapp-...",
      botToken: "xoxb-...",
      userToken: "xoxp-...",
    },
  },
}
\end{lstlisting}


明确设置 userTokenReadOnly 的示例（允许用户令牌写入）：

\begin{lstlisting}[language=json]
{
  channels: {
    slack: {
      enabled: true,
      appToken: "xapp-...",
      botToken: "xoxb-...",
      userToken: "xoxp-...",
      userTokenReadOnly: false,
    },
  },
}
\end{lstlisting}


\paragraph{令牌使用}


\begin{itemize}
  \item 读取操作（历史记录、表情回应列表、置顶列表、表情符号列表、成员信息、搜索）在配置了用户令牌时优先使用用户令牌，否则使用 bot 令牌。
  \item 写入操作（发送/编辑/删除消息、添加/移除表情回应、置顶/取消置顶、文件上传）默认使用 bot 令牌。如果 \texttt{userTokenReadOnly: false} 且没有可用的 bot 令牌，OpenClaw 会回退到用户令牌。
\end{itemize}


\subsubsection{历史上下文}


\begin{itemize}
  \item \texttt{channels.slack.historyLimit}（或 \texttt{channels.slack.accounts.*.historyLimit}）控制将多少条最近的频道/群组消息包含到提示中。
  \item 回退到 \texttt{messages.groupChat.historyLimit}。设置为 \texttt{0} 以禁用（默认 50）。
\end{itemize}


\subsection{HTTP 模式（Events API）}


当你的 Gateway 网关可以通过 HTTPS 被 Slack 访问时（服务器部署的典型情况），使用 HTTP webhook 模式。
HTTP 模式使用 Events API + Interactivity + Slash Commands，共享一个请求 URL。

\subsubsection{设置}


\begin{enumerate}
  \item 创建一个 Slack 应用并\textbf{禁用 Socket Mode}（如果你只使用 HTTP 则可选）。
  \item \textbf{Basic Information} → 复制 \textbf{Signing Secret}。
  \item \textbf{OAuth \& Permissions} → 安装应用并复制 \textbf{Bot User OAuth Token}（\texttt{xoxb-...}）。
  \item \textbf{Event Subscriptions} → 启用事件并将 \textbf{Request URL} 设置为你的 Gateway 网关 webhook 路径（默认 \texttt{/slack/events}）。
  \item \textbf{Interactivity \& Shortcuts} → 启用并设置相同的 \textbf{Request URL}。
  \item \textbf{Slash Commands} → 为你的命令设置相同的 \textbf{Request URL}。
\end{enumerate}


示例请求 URL：
\texttt{https://gateway-host/slack/events}

\subsubsection{OpenClaw 配置（最小）}


\begin{lstlisting}[language=json]
{
  channels: {
    slack: {
      enabled: true,
      mode: "http",
      botToken: "xoxb-...",
      signingSecret: "your-signing-secret",
      webhookPath: "/slack/events",
    },
  },
}
\end{lstlisting}


多账户 HTTP 模式：设置 \texttt{channels.slack.accounts..mode = "http"} 并为每个账户提供唯一的 \texttt{webhookPath}，以便每个 Slack 应用可以指向自己的 URL。

\subsubsection{Manifest（可选）}


使用此 Slack 应用 manifest 快速创建应用（如果需要可以调整名称/命令）。如果你计划配置用户令牌，请包含用户权限范围。

\begin{lstlisting}[language=json]
{
  "display_information": {
    "name": "OpenClaw",
    "description": "Slack connector for OpenClaw"
  },
  "features": {
    "bot_user": {
      "display_name": "OpenClaw",
      "always_online": false
    },
    "app_home": {
      "messages_tab_enabled": true,
      "messages_tab_read_only_enabled": false
    },
    "slash_commands": [
      {
        "command": "/openclaw",
        "description": "Send a message to OpenClaw",
        "should_escape": false
      }
    ]
  },
  "oauth_config": {
    "scopes": {
      "bot": [
        "chat:write",
        "channels:history",
        "channels:read",
        "groups:history",
        "groups:read",
        "groups:write",
        "im:history",
        "im:read",
        "im:write",
        "mpim:history",
        "mpim:read",
        "mpim:write",
        "users:read",
        "app_mentions:read",
        "reactions:read",
        "reactions:write",
        "pins:read",
        "pins:write",
        "emoji:read",
        "commands",
        "files:read",
        "files:write"
      ],
      "user": [
        "channels:history",
        "channels:read",
        "groups:history",
        "groups:read",
        "im:history",
        "im:read",
        "mpim:history",
        "mpim:read",
        "users:read",
        "reactions:read",
        "pins:read",
        "emoji:read",
        "search:read"
      ]
    }
  },
  "settings": {
    "socket_mode_enabled": true,
    "event_subscriptions": {
      "bot_events": [
        "app_mention",
        "message.channels",
        "message.groups",
        "message.im",
        "message.mpim",
        "reaction_added",
        "reaction_removed",
        "member_joined_channel",
        "member_left_channel",
        "channel_rename",
        "pin_added",
        "pin_removed"
      ]
    }
  }
}
\end{lstlisting}


如果启用原生命令，为每个要公开的命令添加一个 \texttt{slash\_commands} 条目（与 \texttt{/help} 列表匹配）。使用 \texttt{channels.slack.commands.native} 覆盖。

\subsection{权限范围（当前 vs 可选）}


Slack 的 Conversations API 是按类型区分的：你只需要你实际接触的会话类型（channels、groups、im、mpim）的权限范围。概述参见 https://docs.slack.dev/apis/web-api/using-the-conversations-api/。

\subsubsection{Bot 令牌权限范围（必需）}


\begin{itemize}
  \item \texttt{chat:write}（通过 \texttt{chat.postMessage} 发送/更新/删除消息）
\end{itemize}

https://docs.slack.dev/reference/methods/chat.postMessage
\begin{itemize}
  \item \texttt{im:write}（通过 \texttt{conversations.open} 打开私信用于用户私信）
\end{itemize}

https://docs.slack.dev/reference/methods/conversations.open
\begin{itemize}
  \item \texttt{channels:history}、\texttt{groups:history}、\texttt{im:history}、\texttt{mpim:history}
\end{itemize}

https://docs.slack.dev/reference/methods/conversations.history
\begin{itemize}
  \item \texttt{channels:read}、\texttt{groups:read}、\texttt{im:read}、\texttt{mpim:read}
\end{itemize}

https://docs.slack.dev/reference/methods/conversations.info
\begin{itemize}
  \item \texttt{users:read}（用户查询）
\end{itemize}

https://docs.slack.dev/reference/methods/users.info
\begin{itemize}
  \item \texttt{reactions:read}、\texttt{reactions:write}（\texttt{reactions.get} / \texttt{reactions.add}）
\end{itemize}

https://docs.slack.dev/reference/methods/reactions.get
https://docs.slack.dev/reference/methods/reactions.add
\begin{itemize}
  \item \texttt{pins:read}、\texttt{pins:write}（\texttt{pins.list} / \texttt{pins.add} / \texttt{pins.remove}）
\end{itemize}

https://docs.slack.dev/reference/scopes/pins.read
https://docs.slack.dev/reference/scopes/pins.write
\begin{itemize}
  \item \texttt{emoji:read}（\texttt{emoji.list}）
\end{itemize}

https://docs.slack.dev/reference/scopes/emoji.read
\begin{itemize}
  \item \texttt{files:write}（通过 \texttt{files.uploadV2} 上传）
\end{itemize}

https://docs.slack.dev/messaging/working-with-files/\#upload

\subsubsection{用户令牌权限范围（可选，默认只读）}


如果你配置了 \texttt{channels.slack.userToken}，在 \textbf{User Token Scopes} 下添加这些。

\begin{itemize}
  \item \texttt{channels:history}、\texttt{groups:history}、\texttt{im:history}、\texttt{mpim:history}
  \item \texttt{channels:read}、\texttt{groups:read}、\texttt{im:read}、\texttt{mpim:read}
  \item \texttt{users:read}
  \item \texttt{reactions:read}
  \item \texttt{pins:read}
  \item \texttt{emoji:read}
  \item \texttt{search:read}
\end{itemize}


\subsubsection{目前不需要（但未来可能需要）}


\begin{itemize}
  \item \texttt{mpim:write}（仅当我们添加群组私信打开/私信启动时通过 \texttt{conversations.open}）
  \item \texttt{groups:write}（仅当我们添加私有频道管理时：创建/重命名/邀请/归档）
  \item \texttt{chat:write.public}（仅当我们想发布到机器人未加入的频道时）
\end{itemize}

https://docs.slack.dev/reference/scopes/chat.write.public
\begin{itemize}
  \item \texttt{users:read.email}（仅当我们需要从 \texttt{users.info} 获取邮箱字段时）
\end{itemize}

https://docs.slack.dev/changelog/2017-04-narrowing-email-access
\begin{itemize}
  \item \texttt{files:read}（仅当我们开始列出/读取文件元数据时）
\end{itemize}


\subsection{配置}


Slack 仅使用 Socket Mode（无 HTTP webhook 服务器）。提供两个令牌：

\begin{lstlisting}[language=json]
{
  "slack": {
    "enabled": true,
    "botToken": "xoxb-...",
    "appToken": "xapp-...",
    "groupPolicy": "allowlist",
    "dm": {
      "enabled": true,
      "policy": "pairing",
      "allowFrom": ["U123", "U456", "*"],
      "groupEnabled": false,
      "groupChannels": ["G123"],
      "replyToMode": "all"
    },
    "channels": {
      "C123": { "allow": true, "requireMention": true },
      "#general": {
        "allow": true,
        "requireMention": true,
        "users": ["U123"],
        "skills": ["search", "docs"],
        "systemPrompt": "Keep answers short."
      }
    },
    "reactionNotifications": "own",
    "reactionAllowlist": ["U123"],
    "replyToMode": "off",
    "actions": {
      "reactions": true,
      "messages": true,
      "pins": true,
      "memberInfo": true,
      "emojiList": true
    },
    "slashCommand": {
      "enabled": true,
      "name": "openclaw",
      "sessionPrefix": "slack:slash",
      "ephemeral": true
    },
    "textChunkLimit": 4000,
    "mediaMaxMb": 20
  }
}
\end{lstlisting}


令牌也可以通过环境变量提供：

\begin{itemize}
  \item \texttt{SLACK\_BOT\_TOKEN}
  \item \texttt{SLACK\_APP\_TOKEN}
\end{itemize}


确认表情回应通过 \texttt{messages.ackReaction} + \texttt{messages.ackReactionScope} 全局控制。使用 \texttt{messages.removeAckAfterReply} 在机器人回复后清除确认表情回应。

\subsection{限制}


\begin{itemize}
  \item 出站文本按 \texttt{channels.slack.textChunkLimit} 分块（默认 4000）。
  \item 可选的换行分块：设置 \texttt{channels.slack.chunkMode="newline"} 以在长度分块之前按空行（段落边界）分割。
  \item 媒体上传受 \texttt{channels.slack.mediaMaxMb} 限制（默认 20）。
\end{itemize}


\subsection{回复线程}


默认情况下，OpenClaw 在主频道回复。使用 \texttt{channels.slack.replyToMode} 控制自动线程：


\begin{table}[h]
\centering
\small
\begin{tabular}{|l|l|}
\hline
模式 & 行为 \\
\hline
\texttt{off} & \textbf{默认。} 在主频道回复。仅当触发消息已在线程中时才使用线程。 \\
\hline
\texttt{first} & 第一条回复进入线程（在触发消息下），后续回复进入主频道。适合保持上下文可见同时避免线程混乱。 \\
\hline
\texttt{all} & 所有回复都进入线程。保持对话集中但可能降低可见性。 \\
\hline
\end{tabular}
\end{table}



该模式适用于自动回复和智能体工具调用（\texttt{slack sendMessage}）。

\subsubsection{按聊天类型的线程}


你可以通过设置 \texttt{channels.slack.replyToModeByChatType} 为每种聊天类型配置不同的线程行为：

\begin{lstlisting}[language=json]
{
  channels: {
    slack: {
      replyToMode: "off", // 频道的默认值
      replyToModeByChatType: {
        direct: "all", // 私信始终使用线程
        group: "first", // 群组私信/MPIM 第一条回复使用线程
      },
    },
  },
}
\end{lstlisting}


支持的聊天类型：

\begin{itemize}
  \item \texttt{direct}：一对一私信（Slack \texttt{im}）
  \item \texttt{group}：群组私信 / MPIM（Slack \texttt{mpim}）
  \item \texttt{channel}：标准频道（公开/私有）
\end{itemize}


优先级：

\begin{enumerate}
  \item \texttt{replyToModeByChatType.}
  \item \texttt{replyToMode}
  \item 提供商默认值（\texttt{off}）
\end{enumerate}


当未设置聊天类型覆盖时，旧版 \texttt{channels.slack.dm.replyToMode} 仍可作为 \texttt{direct} 的回退。

示例：

仅对私信使用线程：

\begin{lstlisting}[language=json]
{
  channels: {
    slack: {
      replyToMode: "off",
      replyToModeByChatType: { direct: "all" },
    },
  },
}
\end{lstlisting}


对群组私信使用线程但保持频道在根级别：

\begin{lstlisting}[language=json]
{
  channels: {
    slack: {
      replyToMode: "off",
      replyToModeByChatType: { group: "first" },
    },
  },
}
\end{lstlisting}


让频道使用线程，保持私信在根级别：

\begin{lstlisting}[language=json]
{
  channels: {
    slack: {
      replyToMode: "first",
      replyToModeByChatType: { direct: "off", group: "off" },
    },
  },
}
\end{lstlisting}


\subsubsection{手动线程标签}


对于细粒度控制，在智能体响应中使用这些标签：

\begin{itemize}
  \item \texttt{[[reply\_to\_current]]} — 回复触发消息（开始/继续线程）。
  \item \texttt{[[reply\_to:]]} — 回复特定的消息 id。
\end{itemize}


\subsection{会话 + 路由}


\begin{itemize}
  \item 私信共享 \texttt{main} 会话（与 WhatsApp/Telegram 相同）。
  \item 频道映射到 \texttt{agent::slack:channel:} 会话。
  \item 斜杠命令使用 \texttt{agent::slack:slash:} 会话（前缀可通过 \texttt{channels.slack.slashCommand.sessionPrefix} 配置）。
  \item 如果 Slack 未提供 \texttt{channel\_type}，OpenClaw 会从频道 ID 前缀（\texttt{D}、\texttt{C}、\texttt{G}）推断并默认为 \texttt{channel} 以保持会话键稳定。
  \item 原生命令注册使用 \texttt{commands.native}（全局默认 \texttt{"auto"} → Slack 关闭），可以使用 \texttt{channels.slack.commands.native} 按工作空间覆盖。文本命令需要独立的 \texttt{/...} 消息，可以使用 \texttt{commands.text: false} 禁用。Slack 斜杠命令在 Slack 应用中管理，不会自动移除。使用 \texttt{commands.useAccessGroups: false} 绕过命令的访问组检查。
  \item 完整命令列表 + 配置：\href{/tools/slash-commands}{斜杠命令}
\end{itemize}


\subsection{私信安全（配对）}


\begin{itemize}
  \item 默认：\texttt{channels.slack.dm.policy="pairing"} — 未知的私信发送者会收到配对码（1 小时后过期）。
  \item 通过以下方式批准：\texttt{openclaw pairing approve slack }。
  \item 要允许任何人：设置 \texttt{channels.slack.dm.policy="open"} 和 \texttt{channels.slack.dm.allowFrom=["*"]}。
  \item \texttt{channels.slack.dm.allowFrom} 接受用户 ID、@用户名或邮箱（在令牌允许时启动时解析）。向导在设置期间接受用户名，并在令牌允许时将其解析为 ID。
\end{itemize}


\subsection{群组策略}


\begin{itemize}
  \item \texttt{channels.slack.groupPolicy} 控制频道处理（\texttt{open|disabled|allowlist}）。
  \item \texttt{allowlist} 要求频道列在 \texttt{channels.slack.channels} 中。
  \item 如果你只设置了 \texttt{SLACK\_BOT\_TOKEN}/\texttt{SLACK\_APP\_TOKEN} 而从未创建 \texttt{channels.slack} 部分，运行时默认将 \texttt{groupPolicy} 设为 \texttt{open}。添加 \texttt{channels.slack.groupPolicy}、\texttt{channels.defaults.groupPolicy} 或频道白名单来锁定它。
  \item 配置向导接受 \texttt{\#channel} 名称，并在可能时（公开 + 私有）将其解析为 ID；如果存在多个匹配，它优先选择活跃的频道。
  \item 启动时，OpenClaw 将白名单中的频道/用户名解析为 ID（在令牌允许时）并记录映射；未解析的条目按原样保留。
  \item 要\textbf{不允许任何频道}，设置 \texttt{channels.slack.groupPolicy: "disabled"}（或保留空白名单）。
\end{itemize}


频道选项（\texttt{channels.slack.channels.} 或 \texttt{channels.slack.channels.}）：

\begin{itemize}
  \item \texttt{allow}：当 \texttt{groupPolicy="allowlist"} 时允许/拒绝频道。
  \item \texttt{requireMention}：频道的提及门控。
  \item \texttt{tools}：可选的每频道工具策略覆盖（\texttt{allow}/\texttt{deny}/\texttt{alsoAllow}）。
  \item \texttt{toolsBySender}：频道内可选的每发送者工具策略覆盖（键为发送者 id/@用户名/邮箱；支持 \texttt{"*"} 通配符）。
  \item \texttt{allowBots}：允许此频道中机器人发送的消息（默认：false）。
  \item \texttt{users}：可选的每频道用户白名单。
  \item \texttt{skills}：Skills 过滤器（省略 = 所有 Skills，空 = 无）。
  \item \texttt{systemPrompt}：频道的额外系统提示（与主题/目的组合）。
  \item \texttt{enabled}：设置为 \texttt{false} 以禁用频道。
\end{itemize}


\subsection{投递目标}


与 cron/CLI 发送一起使用：

\begin{itemize}
  \item \texttt{user:} 用于私信
  \item \texttt{channel:} 用于频道
\end{itemize}


\subsection{工具操作}


Slack 工具操作可以通过 \texttt{channels.slack.actions.*} 进行门控：


\begin{table}[h]
\centering
\small
\begin{tabular}{|l|l|l|}
\hline
操作组 & 默认 & 说明 \\
\hline
reactions & 已启用 & 表情回应 + 列出表情回应 \\
\hline
messages & 已启用 & 读取/发送/编辑/删除 \\
\hline
pins & 已启用 & 置顶/取消置顶/列表 \\
\hline
memberInfo & 已启用 & 成员信息 \\
\hline
emojiList & 已启用 & 自定义表情符号列表 \\
\hline
\end{tabular}
\end{table}



\subsection{安全说明}


\begin{itemize}
  \item 写入默认使用 bot 令牌，因此状态更改操作保持在应用的机器人权限和身份范围内。
  \item 设置 \texttt{userTokenReadOnly: false} 允许在 bot 令牌不可用时使用用户令牌进行写入操作，这意味着操作以安装用户的访问权限运行。将用户令牌视为高权限，并保持操作门控和白名单严格。
  \item 如果你启用用户令牌写入，请确保用户令牌包含你期望的写入权限范围（\texttt{chat:write}、\texttt{reactions:write}、\texttt{pins:write}、\texttt{files:write}），否则这些操作将失败。
\end{itemize}


\subsection{说明}


\begin{itemize}
  \item 提及门控通过 \texttt{channels.slack.channels} 控制（将 \texttt{requireMention} 设置为 \texttt{true}）；\texttt{agents.list[].groupChat.mentionPatterns}（或 \texttt{messages.groupChat.mentionPatterns}）也算作提及。
  \item 多智能体覆盖：在 \texttt{agents.list[].groupChat.mentionPatterns} 上设置每智能体的模式。
  \item 表情回应通知遵循 \texttt{channels.slack.reactionNotifications}（在 \texttt{allowlist} 模式下使用 \texttt{reactionAllowlist}）。
  \item 默认忽略机器人发送的消息；通过 \texttt{channels.slack.allowBots} 或 \texttt{channels.slack.channels..allowBots} 启用。
  \item 警告：如果你允许回复其他机器人（\texttt{channels.slack.allowBots=true} 或 \texttt{channels.slack.channels..allowBots=true}），请使用 \texttt{requireMention}、\texttt{channels.slack.channels..users} 白名单和/或在 \texttt{AGENTS.md} 和 \texttt{SOUL.md} 中设置明确的防护措施来防止机器人之间的回复循环。
  \item 对于 Slack 工具，表情回应移除语义见 \href{/tools/reactions}{/tools/reactions}。
  \item 附件在允许且在大小限制内时会下载到媒体存储。
\end{itemize}



\clearpage

% === 源文件: channels/telegram.md ===
\section{Telegram（Bot API）}


状态：通过 grammY 支持机器人私信和群组，已可用于生产环境。默认使用长轮询；webhook 可选。

\subsection{快速设置（入门）}


\begin{enumerate}
  \item 通过 \textbf{@BotFather}（\href{https://t.me/BotFather}{直达链接}）创建机器人。确认用户名确实是 \texttt{@BotFather}，然后复制 token。
  \item 设置 token：
\end{enumerate}

\begin{itemize}
  \item 环境变量：\texttt{TELEGRAM\_BOT\_TOKEN=...}
  \item 或配置：\texttt{channels.telegram.botToken: "..."}。
  \item 如果两者都设置了，配置优先（环境变量回退仅适用于默认账户）。
\end{itemize}

\begin{enumerate}
  \item 启动 Gateway 网关。
  \item 私信访问默认使用配对模式；首次联系时需要批准配对码。
\end{enumerate}


最小配置：

\begin{lstlisting}[language=json]
{
  channels: {
    telegram: {
      enabled: true,
      botToken: "123:abc",
      dmPolicy: "pairing",
    },
  },
}
\end{lstlisting}


\subsection{这是什么}


\begin{itemize}
  \item 一个由 Gateway 网关拥有的 Telegram Bot API 渠道。
  \item 确定性路由：回复返回到 Telegram；模型不会选择渠道。
  \item 私信共享智能体的主会话；群组保持隔离（\texttt{agent::telegram:group:}）。
\end{itemize}


\subsection{设置（快速路径）}


\subsubsection{1）创建机器人 token（BotFather）}


\begin{enumerate}
  \item 打开 Telegram 并与 \textbf{@BotFather}（\href{https://t.me/BotFather}{直达链接}）对话。确认用户名确实是 \texttt{@BotFather}。
  \item 运行 \texttt{/newbot}，然后按照提示操作（名称 + 以 \texttt{bot} 结尾的用户名）。
  \item 复制 token 并安全保存。
\end{enumerate}


可选的 BotFather 设置：

\begin{itemize}
  \item \texttt{/setjoingroups} — 允许/拒绝将机器人添加到群组。
  \item \texttt{/setprivacy} — 控制机器人是否可以看到所有群组消息。
\end{itemize}


\subsubsection{2）配置 token（环境变量或配置文件）}


示例：

\begin{lstlisting}[language=json]
{
  channels: {
    telegram: {
      enabled: true,
      botToken: "123:abc",
      dmPolicy: "pairing",
      groups: { "*": { requireMention: true } },
    },
  },
}
\end{lstlisting}


环境变量选项：\texttt{TELEGRAM\_BOT\_TOKEN=...}（适用于默认账户）。
如果环境变量和配置都设置了，配置优先。

多账户支持：使用 \texttt{channels.telegram.accounts}，每个账户有独立的 token 和可选的 \texttt{name}。参见 \href{/gateway/configuration\#telegramaccounts--discordaccounts--slackaccounts--signalaccounts--imessageaccounts}{\texttt{gateway/configuration}} 了解共享模式。

\begin{enumerate}
  \item 启动 Gateway 网关。当 token 解析成功时 Telegram 启动（配置优先，环境变量回退）。
  \item 私信访问默认为配对模式。机器人首次被联系时批准配对码。
  \item 对于群组：添加机器人，决定隐私/管理员行为（见下文），然后设置 \texttt{channels.telegram.groups} 来控制提及门控和允许列表。
\end{enumerate}


\subsection{Token + 隐私 + 权限（Telegram 端）}


\subsubsection{Token 创建（BotFather）}


\begin{itemize}
  \item \texttt{/newbot} 创建机器人并返回 token（请保密）。
  \item 如果 token 泄露，通过 @BotFather 撤销/重新生成，并更新你的配置。
\end{itemize}


\subsubsection{群组消息可见性（隐私模式）}


Telegram 机器人默认启用\textbf{隐私模式}，这会限制它们接收哪些群组消息。
如果你的机器人必须看到\textit{所有}群组消息，有两个选项：

\begin{itemize}
  \item 使用 \texttt{/setprivacy} 禁用隐私模式\textbf{或}
  \item 将机器人添加为群组\textbf{管理员}（管理员机器人可以接收所有消息）。
\end{itemize}


\textbf{注意：} 当你切换隐私模式时，Telegram 要求将机器人从每个群组中移除并重新添加，更改才能生效。

\subsubsection{群组权限（管理员权限）}


管理员状态在群组内设置（Telegram UI）。管理员机器人始终接收所有群组消息，因此如果需要完全可见性，请使用管理员身份。

\subsection{工作原理（行为）}


\begin{itemize}
  \item 入站消息被规范化为共享渠道信封，包含回复上下文和媒体占位符。
  \item 群组回复默认需要提及（原生 @提及或 \texttt{agents.list[].groupChat.mentionPatterns} / \texttt{messages.groupChat.mentionPatterns}）。
  \item 多智能体覆盖：在 \texttt{agents.list[].groupChat.mentionPatterns} 上设置每个智能体的模式。
  \item 回复始终路由回同一个 Telegram 聊天。
  \item 长轮询使用 grammY runner，每个聊天按顺序处理；总体并发受 \texttt{agents.defaults.maxConcurrent} 限制。
  \item Telegram Bot API 不支持已读回执；没有 \texttt{sendReadReceipts} 选项。
\end{itemize}


\subsection{草稿流式传输}


OpenClaw 可以在 Telegram 私信中使用 \texttt{sendMessageDraft} 流式传输部分回复。

要求：

\begin{itemize}
  \item 在 @BotFather 中为机器人启用线程模式（论坛话题模式）。
  \item 仅限私聊线程（Telegram 在入站消息中包含 \texttt{message\_thread\_id}）。
  \item \texttt{channels.telegram.streamMode} 未设置为 \texttt{"off"}（默认：\texttt{"partial"}，\texttt{"block"} 启用分块草稿更新）。
\end{itemize}


草稿流式传输仅限私信；Telegram 在群组或频道中不支持此功能。

\subsection{格式化（Telegram HTML）}


\begin{itemize}
  \item 出站 Telegram 文本使用 \texttt{parse\_mode: "HTML"}（Telegram 支持的标签子集）。
  \item 类 Markdown 输入被渲染为 \textbf{Telegram 安全 HTML}（粗体/斜体/删除线/代码/链接）；块级元素被扁平化为带换行/项目符号的文本。
  \item 来自模型的原始 HTML 会被转义，以避免 Telegram 解析错误。
  \item 如果 Telegram 拒绝 HTML 负载，OpenClaw 会以纯文本重试相同的消息。
\end{itemize}


\subsection{命令（原生 + 自定义）}


OpenClaw 在启动时向 Telegram 的机器人菜单注册原生命令（如 \texttt{/status}、\texttt{/reset}、\texttt{/model}）。
你可以通过配置向菜单添加自定义命令：

\begin{lstlisting}[language=json]
{
  channels: {
    telegram: {
      customCommands: [
        { command: "backup", description: "Git 备份" },
        { command: "generate", description: "创建图片" },
      ],
    },
  },
}
\end{lstlisting}


\subsection{故障排除}


\begin{itemize}
  \item 日志中出现 \texttt{setMyCommands failed} 通常意味着到 \texttt{api.telegram.org} 的出站 HTTPS/DNS 被阻止。
  \item 如果你看到 \texttt{sendMessage} 或 \texttt{sendChatAction} 失败，检查 IPv6 路由和 DNS。
\end{itemize}


更多帮助：\href{/channels/troubleshooting}{渠道故障排除}。

注意：

\begin{itemize}
  \item 自定义命令\textbf{仅是菜单条目}；除非你在其他地方处理它们，否则 OpenClaw 不会实现它们。
  \item 命令名称会被规范化（去除前导 \texttt{/}，转为小写），必须匹配 \texttt{a-z}、\texttt{0-9}、\texttt{\_}（1-32 个字符）。
  \item 自定义命令\textbf{不能覆盖原生命令}。冲突会被忽略并记录日志。
  \item 如果禁用了 \texttt{commands.native}，则只注册自定义命令（如果没有则清空）。
\end{itemize}


\subsection{限制}


\begin{itemize}
  \item 出站文本按 \texttt{channels.telegram.textChunkLimit} 分块（默认 4000）。
  \item 可选的换行分块：设置 \texttt{channels.telegram.chunkMode="newline"} 在长度分块之前按空行（段落边界）分割。
  \item 媒体下载/上传受 \texttt{channels.telegram.mediaMaxMb} 限制（默认 5）。
  \item Telegram Bot API 请求在 \texttt{channels.telegram.timeoutSeconds} 后超时（通过 grammY 默认 500）。设置较低的值以避免长时间挂起。
  \item 群组历史上下文使用 \texttt{channels.telegram.historyLimit}（或 \texttt{channels.telegram.accounts.*.historyLimit}），回退到 \texttt{messages.groupChat.historyLimit}。设置 \texttt{0} 禁用（默认 50）。
  \item 私信历史可以用 \texttt{channels.telegram.dmHistoryLimit}（用户轮次）限制。每用户覆盖：\texttt{channels.telegram.dms[""].historyLimit}。
\end{itemize}


\subsection{群组激活模式}


默认情况下，机器人只响应群组中的提及（\texttt{@botname} 或 \texttt{agents.list[].groupChat.mentionPatterns} 中的模式）。要更改此行为：

\subsubsection{通过配置（推荐）}


\begin{lstlisting}[language=json]
{
  channels: {
    telegram: {
      groups: {
        "-1001234567890": { requireMention: false }, // 在此群组中始终响应
      },
    },
  },
}
\end{lstlisting}


\textbf{重要：} 设置 \texttt{channels.telegram.groups} 会创建一个\textbf{允许列表} - 只有列出的群组（或 \texttt{"*"}）会被接受。
论坛话题继承其父群组配置（allowFrom、requireMention、skills、prompts），除非你在 \texttt{channels.telegram.groups..topics.} 下添加每话题覆盖。

要允许所有群组并始终响应：

\begin{lstlisting}[language=json]
{
  channels: {
    telegram: {
      groups: {
        "*": { requireMention: false }, // 所有群组，始终响应
      },
    },
  },
}
\end{lstlisting}


要保持所有群组仅提及响应（默认行为）：

\begin{lstlisting}[language=json]
{
  channels: {
    telegram: {
      groups: {
        "*": { requireMention: true }, // 或完全省略 groups
      },
    },
  },
}
\end{lstlisting}


\subsubsection{通过命令（会话级别）}


在群组中发送：

\begin{itemize}
  \item \texttt{/activation always} - 响应所有消息
  \item \texttt{/activation mention} - 需要提及（默认）
\end{itemize}


\textbf{注意：} 命令只更新会话状态。要在重启后保持持久行为，请使用配置。

\subsubsection{获取群组聊天 ID}


将群组中的任何消息转发给 Telegram 上的 \texttt{@userinfobot} 或 \texttt{@getidsbot} 以查看聊天 ID（负数，如 \texttt{-1001234567890}）。

\textbf{提示：} 要获取你自己的用户 ID，私信机器人，它会回复你的用户 ID（配对消息），或者在命令启用后使用 \texttt{/whoami}。

\textbf{隐私注意：} \texttt{@userinfobot} 是第三方机器人。如果你更倾向于其他方式，将机器人添加到群组，发送一条消息，然后使用 \texttt{openclaw logs --follow} 读取 \texttt{chat.id}，或使用 Bot API \texttt{getUpdates}。

\subsection{配置写入}


默认情况下，Telegram 允许写入由渠道事件或 \texttt{/config set|unset} 触发的配置更新。

这发生在以下情况：

\begin{itemize}
  \item 群组升级为超级群组，Telegram 发出 \texttt{migrate\_to\_chat\_id}（聊天 ID 更改）。OpenClaw 可以自动迁移 \texttt{channels.telegram.groups}。
  \item 你在 Telegram 聊天中运行 \texttt{/config set} 或 \texttt{/config unset}（需要 \texttt{commands.config: true}）。
\end{itemize}


禁用方式：

\begin{lstlisting}[language=json]
{
  channels: { telegram: { configWrites: false } },
}
\end{lstlisting}


\subsection{话题（论坛超级群组）}


Telegram 论坛话题在每条消息中包含 \texttt{message\_thread\_id}。OpenClaw：

\begin{itemize}
  \item 将 \texttt{:topic:} 附加到 Telegram 群组会话键，使每个话题隔离。
  \item 发送输入指示器和回复时带上 \texttt{message\_thread\_id}，使响应保持在话题内。
  \item 通用话题（线程 id \texttt{1}）是特殊的：消息发送省略 \texttt{message\_thread\_id}（Telegram 会拒绝），但输入指示器仍然包含它。
  \item 在模板上下文中暴露 \texttt{MessageThreadId} + \texttt{IsForum} 用于路由/模板。
  \item 话题特定配置可在 \texttt{channels.telegram.groups..topics.} 下设置（skills、允许列表、自动回复、系统提示、禁用）。
  \item 话题配置继承群组设置（requireMention、允许列表、skills、提示、enabled），除非每话题覆盖。
\end{itemize}


私聊在某些边缘情况下可能包含 \texttt{message\_thread\_id}。OpenClaw 保持私信会话键不变，但在存在线程 id 时仍将其用于回复/草稿流式传输。

\subsection{内联按钮}


Telegram 支持带回调按钮的内联键盘。

\begin{lstlisting}[language=json]
{
  channels: {
    telegram: {
      capabilities: {
        inlineButtons: "allowlist",
      },
    },
  },
}
\end{lstlisting}


对于每账户配置：

\begin{lstlisting}[language=json]
{
  channels: {
    telegram: {
      accounts: {
        main: {
          capabilities: {
            inlineButtons: "allowlist",
          },
        },
      },
    },
  },
}
\end{lstlisting}


作用域：

\begin{itemize}
  \item \texttt{off} — 禁用内联按钮
  \item \texttt{dm} — 仅私信（群组目标被阻止）
  \item \texttt{group} — 仅群组（私信目标被阻止）
  \item \texttt{all} — 私信 + 群组
  \item \texttt{allowlist} — 私信 + 群组，但仅限 \texttt{allowFrom}/\texttt{groupAllowFrom} 允许的发送者（与控制命令规则相同）
\end{itemize}


默认：\texttt{allowlist}。
旧版：\texttt{capabilities: ["inlineButtons"]} = \texttt{inlineButtons: "all"}。

\subsubsection{发送按钮}


使用带 \texttt{buttons} 参数的消息工具：

\begin{lstlisting}[language=json]
{
  action: "send",
  channel: "telegram",
  to: "123456789",
  message: "选择一个选项：",
  buttons: [
    [
      { text: "是", callback_data: "yes" },
      { text: "否", callback_data: "no" },
    ],
    [{ text: "取消", callback_data: "cancel" }],
  ],
}
\end{lstlisting}


当用户点击按钮时，回调数据会以以下格式作为消息发送回智能体：
\texttt{callback\_data: value}

\subsubsection{配置选项}


Telegram 功能可以在两个级别配置（上面显示的对象形式；旧版字符串数组仍然支持）：

\begin{itemize}
  \item \texttt{channels.telegram.capabilities}：应用于所有 Telegram 账户的全局默认功能配置，除非被覆盖。
  \item \texttt{channels.telegram.accounts..capabilities}：每账户功能，覆盖该特定账户的全局默认值。
\end{itemize}


当所有 Telegram 机器人/账户应具有相同行为时使用全局设置。当不同机器人需要不同行为时使用每账户配置（例如，一个账户只处理私信，而另一个允许在群组中使用）。

\subsection{访问控制（私信 + 群组）}


\subsubsection{私信访问}


\begin{itemize}
  \item 默认：\texttt{channels.telegram.dmPolicy = "pairing"}。未知发送者收到配对码；在批准之前消息被忽略（配对码 1 小时后过期）。
  \item 批准方式：
  \item \texttt{openclaw pairing list telegram}
  \item \texttt{openclaw pairing approve telegram }
  \item 配对是 Telegram 私信使用的默认 token 交换。详情：\href{/channels/pairing}{配对}
  \item \texttt{channels.telegram.allowFrom} 接受数字用户 ID（推荐）或 \texttt{@username} 条目。这\textbf{不是}机器人用户名；使用人类发送者的 ID。向导接受 \texttt{@username} 并在可能时将其解析为数字 ID。
\end{itemize}


\paragraph{查找你的 Telegram 用户 ID}


更安全（无第三方机器人）：

\begin{enumerate}
  \item 启动 Gateway 网关并私信你的机器人。
  \item 运行 \texttt{openclaw logs --follow} 并查找 \texttt{from.id}。
\end{enumerate}


备选（官方 Bot API）：

\begin{enumerate}
  \item 私信你的机器人。
  \item 使用你的机器人 token 获取更新并读取 \texttt{message.from.id}：
\begin{lstlisting}[language=bash]
   curl "https://api.telegram.org/bot<bot_token>/getUpdates"
\end{lstlisting}

\end{enumerate}


第三方（隐私性较低）：

\begin{itemize}
  \item 私信 \texttt{@userinfobot} 或 \texttt{@getidsbot} 并使用返回的用户 id。
\end{itemize}


\subsubsection{群组访问}


两个独立的控制：

\textbf{1. 允许哪些群组}（通过 \texttt{channels.telegram.groups} 的群组允许列表）：

\begin{itemize}
  \item 无 \texttt{groups} 配置 = 允许所有群组
  \item 有 \texttt{groups} 配置 = 只允许列出的群组或 \texttt{"*"}
  \item 示例：\texttt{"groups": \{ "-1001234567890": \{\}, "*": \{\} \}} 允许所有群组
\end{itemize}


\textbf{2. 允许哪些发送者}（通过 \texttt{channels.telegram.groupPolicy} 的发送者过滤）：

\begin{itemize}
  \item \texttt{"open"} = 允许群组中的所有发送者发消息
  \item \texttt{"allowlist"} = 只有 \texttt{channels.telegram.groupAllowFrom} 中的发送者可以发消息
  \item \texttt{"disabled"} = 不接受任何群组消息
\end{itemize}

默认是 \texttt{groupPolicy: "allowlist"}（除非添加 \texttt{groupAllowFrom} 否则被阻止）。

大多数用户需要：\texttt{groupPolicy: "allowlist"} + \texttt{groupAllowFrom} + 在 \texttt{channels.telegram.groups} 中列出特定群组

\subsection{长轮询 vs webhook}


\begin{itemize}
  \item 默认：长轮询（不需要公共 URL）。
  \item Webhook 模式：设置 \texttt{channels.telegram.webhookUrl} 和 \texttt{channels.telegram.webhookSecret}（可选 \texttt{channels.telegram.webhookPath}）。
  \item 本地监听器绑定到 \texttt{0.0.0.0:8787}，默认服务于 \texttt{POST /telegram-webhook}。
  \item 如果你的公共 URL 不同，使用反向代理并将 \texttt{channels.telegram.webhookUrl} 指向公共端点。
\end{itemize}


\subsection{回复线程}


Telegram 通过标签支持可选的线程回复：

\begin{itemize}
  \item \texttt{[[reply\_to\_current]]} -- 回复触发消息。
  \item \texttt{[[reply\_to:]]} -- 回复特定消息 id。
\end{itemize}


通过 \texttt{channels.telegram.replyToMode} 控制：

\begin{itemize}
  \item \texttt{first}（默认）、\texttt{all}、\texttt{off}。
\end{itemize}


\subsection{音频消息（语音 vs 文件）}


Telegram 区分\textbf{语音备忘录}（圆形气泡）和\textbf{音频文件}（元数据卡片）。
OpenClaw 默认使用音频文件以保持向后兼容性。

要在智能体回复中强制使用语音备忘录气泡，在回复中的任何位置包含此标签：

\begin{itemize}
  \item \texttt{[[audio\_as\_voice]]} — 将音频作为语音备忘录而不是文件发送。
\end{itemize}


该标签会从发送的文本中去除。其他渠道会忽略此标签。

对于消息工具发送，设置 \texttt{asVoice: true} 并配合兼容语音的音频 \texttt{media} URL（当存在 media 时 \texttt{message} 是可选的）：

\begin{lstlisting}[language=json]
{
  action: "send",
  channel: "telegram",
  to: "123456789",
  media: "https://example.com/voice.ogg",
  asVoice: true,
}
\end{lstlisting}


\subsection{贴纸}


OpenClaw 支持接收和发送 Telegram 贴纸，并具有智能缓存功能。

\subsubsection{接收贴纸}


当用户发送贴纸时，OpenClaw 根据贴纸类型处理：

\begin{itemize}
  \item \textbf{静态贴纸（WEBP）：} 下载并通过视觉处理。贴纸在消息内容中显示为 `` 占位符。
  \item \textbf{动画贴纸（TGS）：} 跳过（Lottie 格式不支持处理）。
  \item \textbf{视频贴纸（WEBM）：} 跳过（视频格式不支持处理）。
\end{itemize}


接收贴纸时可用的模板上下文字段：

\begin{itemize}
  \item \texttt{Sticker} — 包含以下属性的对象：
  \item \texttt{emoji} — 与贴纸关联的表情符号
  \item \texttt{setName} — 贴纸集名称
  \item \texttt{fileId} — Telegram 文件 ID（用于发送相同贴纸）
  \item \texttt{fileUniqueId} — 用于缓存查找的稳定 ID
  \item \texttt{cachedDescription} — 可用时的缓存视觉描述
\end{itemize}


\subsubsection{贴纸缓存}


贴纸通过 AI 的视觉功能处理以生成描述。由于相同的贴纸经常重复发送，OpenClaw 缓存这些描述以避免冗余的 API 调用。

\textbf{工作原理：}

\begin{enumerate}
  \item \textbf{首次遇到：} 贴纸图像被发送给 AI 进行视觉分析。AI 生成描述（例如"一只卡通猫热情地挥手"）。
  \item \textbf{缓存存储：} 描述与贴纸的文件 ID、表情符号和集合名称一起保存。
  \item \textbf{后续遇到：} 当再次看到相同贴纸时，直接使用缓存的描述。图像不会发送给 AI。
\end{enumerate}


\textbf{缓存位置：} \texttt{\textasciitilde{}/.openclaw/telegram/sticker-cache.json}

\textbf{缓存条目格式：}

\begin{lstlisting}[language=json]
{
  "fileId": "CAACAgIAAxkBAAI...",
  "fileUniqueId": "AgADBAADb6cxG2Y",
  "emoji": "👋",
  "setName": "CoolCats",
  "description": "一只卡通猫热情地挥手",
  "cachedAt": "2026-01-15T10:30:00.000Z"
}
\end{lstlisting}


\textbf{优点：}

\begin{itemize}
  \item 通过避免对相同贴纸重复调用视觉 API 来降低 API 成本
  \item 缓存贴纸响应更快（无视觉处理延迟）
  \item 基于缓存描述启用贴纸搜索功能
\end{itemize}


缓存在接收贴纸时自动填充。无需手动缓存管理。

\subsubsection{发送贴纸}


智能体可以使用 \texttt{sticker} 和 \texttt{sticker-search} 动作发送和搜索贴纸。这些默认禁用，必须在配置中启用：

\begin{lstlisting}[language=json]
{
  channels: {
    telegram: {
      actions: {
        sticker: true,
      },
    },
  },
}
\end{lstlisting}


\textbf{发送贴纸：}

\begin{lstlisting}[language=json]
{
  action: "sticker",
  channel: "telegram",
  to: "123456789",
  fileId: "CAACAgIAAxkBAAI...",
}
\end{lstlisting}


参数：

\begin{itemize}
  \item \texttt{fileId}（必需）— 贴纸的 Telegram 文件 ID。从接收贴纸时的 \texttt{Sticker.fileId} 获取，或从 \texttt{sticker-search} 结果获取。
  \item \texttt{replyTo}（可选）— 要回复的消息 ID。
  \item \texttt{threadId}（可选）— 论坛话题的消息线程 ID。
\end{itemize}


\textbf{搜索贴纸：}

智能体可以按描述、表情符号或集合名称搜索缓存的贴纸：

\begin{lstlisting}[language=json]
{
  action: "sticker-search",
  channel: "telegram",
  query: "猫 挥手",
  limit: 5,
}
\end{lstlisting}


返回缓存中匹配的贴纸：

\begin{lstlisting}[language=json]
{
  ok: true,
  count: 2,
  stickers: [
    {
      fileId: "CAACAgIAAxkBAAI...",
      emoji: "👋",
      description: "一只卡通猫热情地挥手",
      setName: "CoolCats",
    },
  ],
}
\end{lstlisting}


搜索在描述文本、表情符号字符和集合名称之间使用模糊匹配。

\textbf{带线程的示例：}

\begin{lstlisting}[language=json]
{
  action: "sticker",
  channel: "telegram",
  to: "-1001234567890",
  fileId: "CAACAgIAAxkBAAI...",
  replyTo: 42,
  threadId: 123,
}
\end{lstlisting}


\subsection{流式传输（草稿）}


Telegram 可以在智能体生成响应时流式传输\textbf{草稿气泡}。
OpenClaw 使用 Bot API \texttt{sendMessageDraft}（不是真实消息），然后将最终回复作为普通消息发送。

要求（Telegram Bot API 9.3+）：

\begin{itemize}
  \item \textbf{启用话题的私聊}（机器人的论坛话题模式）。
  \item 入站消息必须包含 \texttt{message\_thread\_id}（私有话题线程）。
  \item 群组/超级群组/频道的流式传输被忽略。
\end{itemize}


配置：

\begin{itemize}
  \item \texttt{channels.telegram.streamMode: "off" | "partial" | "block"}（默认：\texttt{partial}）
  \item \texttt{partial}：用最新的流式文本更新草稿气泡。
  \item \texttt{block}：以较大块（分块）更新草稿气泡。
  \item \texttt{off}：禁用草稿流式传输。
  \item 可选（仅用于 \texttt{streamMode: "block"}）：
  \item \texttt{channels.telegram.draftChunk: \{ minChars?, maxChars?, breakPreference? \}}
  \item 默认值：\texttt{minChars: 200}、\texttt{maxChars: 800}、\texttt{breakPreference: "paragraph"}（限制在 \texttt{channels.telegram.textChunkLimit} 内）。
\end{itemize}


注意：草稿流式传输与\textbf{分块流式传输}（渠道消息）不同。
分块流式传输默认关闭，如果你想要早期 Telegram 消息而不是草稿更新，需要 \texttt{channels.telegram.blockStreaming: true}。

推理流（仅限 Telegram）：

\begin{itemize}
  \item \texttt{/reasoning stream} 在回复生成时将推理流式传输到草稿气泡中，然后发送不带推理的最终答案。
  \item 如果 \texttt{channels.telegram.streamMode} 为 \texttt{off}，推理流被禁用。
\end{itemize}

更多上下文：\href{/concepts/streaming}{流式传输 + 分块}。

\subsection{重试策略}


出站 Telegram API 调用在遇到临时网络/429 错误时会以指数退避和抖动进行重试。通过 \texttt{channels.telegram.retry} 配置。参见\href{/concepts/retry}{重试策略}。

\subsection{智能体工具（消息 + 反应）}


\begin{itemize}
  \item 工具：\texttt{telegram}，使用 \texttt{sendMessage} 动作（\texttt{to}、\texttt{content}，可选 \texttt{mediaUrl}、\texttt{replyToMessageId}、\texttt{messageThreadId}）。
  \item 工具：\texttt{telegram}，使用 \texttt{react} 动作（\texttt{chatId}、\texttt{messageId}、\texttt{emoji}）。
  \item 工具：\texttt{telegram}，使用 \texttt{deleteMessage} 动作（\texttt{chatId}、\texttt{messageId}）。
  \item 反应移除语义：参见 \href{/tools/reactions}{/tools/reactions}。
  \item 工具门控：\texttt{channels.telegram.actions.reactions}、\texttt{channels.telegram.actions.sendMessage}、\texttt{channels.telegram.actions.deleteMessage}（默认：启用），以及 \texttt{channels.telegram.actions.sticker}（默认：禁用）。
\end{itemize}


\subsection{反应通知}


\textbf{反应工作原理：}
Telegram 反应作为\textbf{单独的 \texttt{message\_reaction} 事件}到达，而不是消息负载中的属性。当用户添加反应时，OpenClaw：

\begin{enumerate}
  \item 从 Telegram API 接收 \texttt{message\_reaction} 更新
  \item 将其转换为\textbf{系统事件}，格式为：\texttt{"Telegram reaction added: \{emoji\} by \{user\} on msg \{id\}"}
  \item 使用与常规消息\textbf{相同的会话键}将系统事件加入队列
  \item 当该对话中的下一条消息到达时，系统事件被排出并前置到智能体的上下文中
\end{enumerate}


智能体将反应视为对话历史中的\textbf{系统通知}，而不是消息元数据。

\textbf{配置：}

\begin{itemize}
  \item \texttt{channels.telegram.reactionNotifications}：控制哪些反应触发通知
  \item \texttt{"off"} — 忽略所有反应
  \item \texttt{"own"} — 当用户对机器人消息做出反应时通知（尽力而为；内存中）（默认）
  \item \texttt{"all"} — 通知所有反应
\end{itemize}


\begin{itemize}
  \item \texttt{channels.telegram.reactionLevel}：控制智能体的反应能力
  \item \texttt{"off"} — 智能体不能对消息做出反应
  \item \texttt{"ack"} — 机器人发送确认反应（处理时显示 \emoji{👀}）（默认）
  \item \texttt{"minimal"} — 智能体可以少量反应（指导：每 5-10 次交换 1 次）
  \item \texttt{"extensive"} — 智能体可以在适当时自由反应
\end{itemize}


\textbf{论坛群组：} 论坛群组中的反应包含 \texttt{message\_thread\_id}，使用类似 \texttt{agent:main:telegram:group:\{chatId\}:topic:\{threadId\}} 的会话键。这确保同一话题中的反应和消息保持在一起。

\textbf{示例配置：}

\begin{lstlisting}[language=json]
{
  channels: {
    telegram: {
      reactionNotifications: "all", // 查看所有反应
      reactionLevel: "minimal", // 智能体可以少量反应
    },
  },
}
\end{lstlisting}


\textbf{要求：}

\begin{itemize}
  \item Telegram 机器人必须在 \texttt{allowed\_updates} 中明确请求 \texttt{message\_reaction}（由 OpenClaw 自动配置）
  \item 对于 webhook 模式，反应包含在 webhook \texttt{allowed\_updates} 中
  \item 对于轮询模式，反应包含在 \texttt{getUpdates} \texttt{allowed\_updates} 中
\end{itemize}


\subsection{投递目标（CLI/cron）}


\begin{itemize}
  \item 使用聊天 id（\texttt{123456789}）或用户名（\texttt{@name}）作为目标。
  \item 示例：\texttt{openclaw message send --channel telegram --target 123456789 --message "hi"}。
\end{itemize}


\subsection{故障排除}


\textbf{机器人不响应群组中的非提及消息：}

\begin{itemize}
  \item 如果你设置了 \texttt{channels.telegram.groups.*.requireMention=false}，Telegram 的 Bot API \textbf{隐私模式}必须禁用。
  \item BotFather：\texttt{/setprivacy} → \textbf{Disable}（然后从群组中移除并重新添加机器人）
  \item \texttt{openclaw channels status} 在配置期望未提及群组消息时显示警告。
  \item \texttt{openclaw channels status --probe} 可以额外检查显式数字群组 ID 的成员资格（它无法审计通配符 \texttt{"*"} 规则）。
  \item 快速测试：\texttt{/activation always}（仅会话级别；使用配置以持久化）
\end{itemize}


\textbf{机器人完全看不到群组消息：}

\begin{itemize}
  \item 如果设置了 \texttt{channels.telegram.groups}，群组必须被列出或使用 \texttt{"*"}
  \item 在 @BotFather 中检查隐私设置 →"Group Privacy"应为 \textbf{OFF}
  \item 验证机器人确实是成员（不仅仅是没有读取权限的管理员）
  \item 检查 Gateway 网关日志：\texttt{openclaw logs --follow}（查找"skipping group message"）
\end{itemize}


\textbf{机器人响应提及但不响应 \texttt{/activation always}：}

\begin{itemize}
  \item \texttt{/activation} 命令更新会话状态但不持久化到配置
  \item 要持久化行为，将群组添加到 \texttt{channels.telegram.groups} 并设置 \texttt{requireMention: false}
\end{itemize}


\textbf{像 \texttt{/status} 这样的命令不起作用：}

\begin{itemize}
  \item 确保你的 Telegram 用户 ID 已授权（通过配对或 \texttt{channels.telegram.allowFrom}）
  \item 即使在 \texttt{groupPolicy: "open"} 的群组中，命令也需要授权
\end{itemize}


\textbf{长轮询在 Node 22+ 上立即中止（通常与代理/自定义 fetch 有关）：}

\begin{itemize}
  \item Node 22+ 对 \texttt{AbortSignal} 实例更严格；外部信号可以立即中止 \texttt{fetch} 调用。
  \item 升级到规范化中止信号的 OpenClaw 构建版本，或在可以升级之前在 Node 20 上运行 Gateway 网关。
\end{itemize}


\textbf{机器人启动后静默停止响应（或日志显示 \texttt{HttpError: Network request ... failed}）：}

\begin{itemize}
  \item 某些主机首先将 \texttt{api.telegram.org} 解析为 IPv6。如果你的服务器没有可用的 IPv6 出口，grammY 可能会卡在仅 IPv6 的请求上。
  \item 通过启用 IPv6 出口\textbf{或}强制 \texttt{api.telegram.org} 使用 IPv4 解析来修复（例如，使用 IPv4 A 记录添加 \texttt{/etc/hosts} 条目，或在你的 OS DNS 堆栈中优先使用 IPv4），然后重启 Gateway 网关。
  \item 快速检查：\texttt{dig +short api.telegram.org A} 和 \texttt{dig +short api.telegram.org AAAA} 确认 DNS 返回的内容。
\end{itemize}


\subsection{配置参考（Telegram）}


完整配置：\href{/gateway/configuration}{配置}

提供商选项：

\begin{itemize}
  \item \texttt{channels.telegram.enabled}：启用/禁用渠道启动。
  \item \texttt{channels.telegram.botToken}：机器人 token（BotFather）。
  \item \texttt{channels.telegram.tokenFile}：从文件路径读取 token。
  \item \texttt{channels.telegram.dmPolicy}：\texttt{pairing | allowlist | open | disabled}（默认：pairing）。
  \item \texttt{channels.telegram.allowFrom}：私信允许列表（id/用户名）。\texttt{open} 需要 \texttt{"*"}。
  \item \texttt{channels.telegram.groupPolicy}：\texttt{open | allowlist | disabled}（默认：allowlist）。
  \item \texttt{channels.telegram.groupAllowFrom}：群组发送者允许列表（id/用户名）。
  \item \texttt{channels.telegram.groups}：每群组默认值 + 允许列表（使用 \texttt{"*"} 作为全局默认值）。
  \item \texttt{channels.telegram.groups..requireMention}：提及门控默认值。
  \item \texttt{channels.telegram.groups..skills}：skill 过滤器（省略 = 所有 skills，空 = 无）。
  \item \texttt{channels.telegram.groups..allowFrom}：每群组发送者允许列表覆盖。
  \item \texttt{channels.telegram.groups..systemPrompt}：群组的额外系统提示。
  \item \texttt{channels.telegram.groups..enabled}：为 \texttt{false} 时禁用群组。
  \item \texttt{channels.telegram.groups..topics..*}：每话题覆盖（与群组相同的字段）。
  \item \texttt{channels.telegram.groups..topics..requireMention}：每话题提及门控覆盖。
  \item \texttt{channels.telegram.capabilities.inlineButtons}：\texttt{off | dm | group | all | allowlist}（默认：allowlist）。
  \item \texttt{channels.telegram.accounts..capabilities.inlineButtons}：每账户覆盖。
  \item \texttt{channels.telegram.replyToMode}：\texttt{off | first | all}（默认：\texttt{off}）。
  \item \texttt{channels.telegram.textChunkLimit}：出站分块大小（字符）。
  \item \texttt{channels.telegram.chunkMode}：\texttt{length}（默认）或 \texttt{newline} 在长度分块之前按空行（段落边界）分割。
  \item \texttt{channels.telegram.linkPreview}：切换出站消息的链接预览（默认：true）。
  \item \texttt{channels.telegram.streamMode}：\texttt{off | partial | block}（草稿流式传输）。
  \item \texttt{channels.telegram.mediaMaxMb}：入站/出站媒体上限（MB）。
  \item \texttt{channels.telegram.retry}：出站 Telegram API 调用的重试策略（attempts、minDelayMs、maxDelayMs、jitter）。
  \item \texttt{channels.telegram.network.autoSelectFamily}：覆盖 Node autoSelectFamily（true=启用，false=禁用）。在 Node 22 上默认禁用以避免 Happy Eyeballs 超时。
  \item \texttt{channels.telegram.proxy}：Bot API 调用的代理 URL（SOCKS/HTTP）。
  \item \texttt{channels.telegram.webhookUrl}：启用 webhook 模式（需要 \texttt{channels.telegram.webhookSecret}）。
  \item \texttt{channels.telegram.webhookSecret}：webhook 密钥（设置 webhookUrl 时必需）。
  \item \texttt{channels.telegram.webhookPath}：本地 webhook 路径（默认 \texttt{/telegram-webhook}）。
  \item \texttt{channels.telegram.actions.reactions}：门控 Telegram 工具反应。
  \item \texttt{channels.telegram.actions.sendMessage}：门控 Telegram 工具消息发送。
  \item \texttt{channels.telegram.actions.deleteMessage}：门控 Telegram 工具消息删除。
  \item \texttt{channels.telegram.actions.sticker}：门控 Telegram 贴纸动作 — 发送和搜索（默认：false）。
  \item \texttt{channels.telegram.reactionNotifications}：\texttt{off | own | all} — 控制哪些反应触发系统事件（未设置时默认：\texttt{own}）。
  \item \texttt{channels.telegram.reactionLevel}：\texttt{off | ack | minimal | extensive} — 控制智能体的反应能力（未设置时默认：\texttt{minimal}）。
\end{itemize}


相关全局选项：

\begin{itemize}
  \item \texttt{agents.list[].groupChat.mentionPatterns}（提及门控模式）。
  \item \texttt{messages.groupChat.mentionPatterns}（全局回退）。
  \item \texttt{commands.native}（默认为 \texttt{"auto"} → Telegram/Discord 开启，Slack 关闭）、\texttt{commands.text}、\texttt{commands.useAccessGroups}（命令行为）。使用 \texttt{channels.telegram.commands.native} 覆盖。
  \item \texttt{messages.responsePrefix}、\texttt{messages.ackReaction}、\texttt{messages.ackReactionScope}、\texttt{messages.removeAckAfterReply}。
\end{itemize}



\clearpage

% === 源文件: channels/tlon.md ===
\section{Tlon（插件）}


Tlon 是一个基于 Urbit 构建的去中心化即时通讯工具。OpenClaw 连接到你的 Urbit ship，可以响应私信和群聊消息。群组回复默认需要 @ 提及，并可通过允许列表进一步限制。

状态：通过插件支持。支持私信、群组提及、话题回复和纯文本媒体回退（URL 附加到说明文字）。不支持表情回应、投票和原生媒体上传。

\subsection{需要插件}


Tlon 作为插件提供，不包含在核心安装中。

通过 CLI 安装（npm 仓库）：

\begin{lstlisting}[language=bash]
openclaw plugins install @openclaw/tlon
\end{lstlisting}


本地检出（从 git 仓库运行时）：

\begin{lstlisting}[language=bash]
openclaw plugins install ./extensions/tlon
\end{lstlisting}


详情：\href{/tools/plugin}{插件}

\subsection{设置}


\begin{enumerate}
  \item 安装 Tlon 插件。
  \item 获取你的 ship URL 和登录代码。
  \item 配置 \texttt{channels.tlon}。
  \item 重启 Gateway 网关。
  \item 私信机器人或在群组频道中提及它。
\end{enumerate}


最小配置（单账户）：

\begin{lstlisting}[language=json]
{
  channels: {
    tlon: {
      enabled: true,
      ship: "~sampel-palnet",
      url: "https://your-ship-host",
      code: "lidlut-tabwed-pillex-ridrup",
    },
  },
}
\end{lstlisting}


\subsection{群组频道}


默认启用自动发现。你也可以手动固定频道：

\begin{lstlisting}[language=json]
{
  channels: {
    tlon: {
      groupChannels: ["chat/~host-ship/general", "chat/~host-ship/support"],
    },
  },
}
\end{lstlisting}


禁用自动发现：

\begin{lstlisting}[language=json]
{
  channels: {
    tlon: {
      autoDiscoverChannels: false,
    },
  },
}
\end{lstlisting}


\subsection{访问控制}


私信允许列表（空 = 允许全部）：

\begin{lstlisting}[language=json]
{
  channels: {
    tlon: {
      dmAllowlist: ["~zod", "~nec"],
    },
  },
}
\end{lstlisting}


群组授权（默认受限）：

\begin{lstlisting}[language=json]
{
  channels: {
    tlon: {
      defaultAuthorizedShips: ["~zod"],
      authorization: {
        channelRules: {
          "chat/~host-ship/general": {
            mode: "restricted",
            allowedShips: ["~zod", "~nec"],
          },
          "chat/~host-ship/announcements": {
            mode: "open",
          },
        },
      },
    },
  },
}
\end{lstlisting}


\subsection{投递目标（CLI/cron）}


与 \texttt{openclaw message send} 或 cron 投递一起使用：

\begin{itemize}
  \item 私信：\texttt{\textasciitilde{}sampel-palnet} 或 \texttt{dm/\textasciitilde{}sampel-palnet}
  \item 群组：\texttt{chat/\textasciitilde{}host-ship/channel} 或 \texttt{group:\textasciitilde{}host-ship/channel}
\end{itemize}


\subsection{注意事项}


\begin{itemize}
  \item 群组回复需要提及（例如 \texttt{\textasciitilde{}your-bot-ship}）才能响应。
  \item 话题回复：如果入站消息在话题中，OpenClaw 会在话题内回复。
  \item 媒体：\texttt{sendMedia} 回退为文本 + URL（无原生上传）。
\end{itemize}



\clearpage

% === 源文件: channels/troubleshooting.md ===
\section{渠道故障排除}


首先运行：

\begin{lstlisting}[language=bash]
openclaw doctor
openclaw channels status --probe
\end{lstlisting}


\texttt{channels status --probe} 会在检测到常见渠道配置错误时输出警告，并包含小型实时检查（凭据、部分权限/成员资格）。

\subsection{渠道}


\begin{itemize}
  \item Discord：\href{/channels/discord\#troubleshooting}{/channels/discord\#troubleshooting}
  \item Telegram：\href{/channels/telegram\#troubleshooting}{/channels/telegram\#troubleshooting}
  \item WhatsApp：\href{/channels/whatsapp\#troubleshooting-quick}{/channels/whatsapp\#troubleshooting-quick}
\end{itemize}


\subsection{Telegram 快速修复}


\begin{itemize}
  \item 日志显示 \texttt{HttpError: Network request for 'sendMessage' failed} 或 \texttt{sendChatAction} → 检查 IPv6 DNS。如果 \texttt{api.telegram.org} 优先解析为 IPv6 而主机缺少 IPv6 出站连接，请强制使用 IPv4 或启用 IPv6。参见 \href{/channels/telegram\#troubleshooting}{/channels/telegram\#troubleshooting}。
  \item 日志显示 \texttt{setMyCommands failed} → 检查到 \texttt{api.telegram.org} 的出站 HTTPS 和 DNS 可达性（常见于限制严格的 VPS 或代理环境）。
\end{itemize}



\clearpage

% === 源文件: channels/twitch.md ===
\section{Twitch（插件）}


通过 IRC 连接支持 Twitch 聊天。OpenClaw 以 Twitch 用户（机器人账户）身份连接，在频道中接收和发送消息。

\subsection{需要插件}


Twitch 作为插件发布，未与核心安装捆绑。

通过 CLI 安装（npm 注册表）：

\begin{lstlisting}[language=bash]
openclaw plugins install @openclaw/twitch
\end{lstlisting}


本地检出（从 git 仓库运行时）：

\begin{lstlisting}[language=bash]
openclaw plugins install ./extensions/twitch
\end{lstlisting}


详情：\href{/tools/plugin}{插件}

\subsection{快速设置（新手）}


\begin{enumerate}
  \item 为机器人创建一个专用的 Twitch 账户（或使用现有账户）。
  \item 生成凭证：\href{https://twitchtokengenerator.com/}{Twitch Token Generator}
\end{enumerate}

\begin{itemize}
  \item 选择 \textbf{Bot Token}
  \item 确认已选择 \texttt{chat:read} 和 \texttt{chat:write} 权限范围
  \item 复制 \textbf{Client ID} 和 \textbf{Access Token}
\end{itemize}

\begin{enumerate}
  \item 查找你的 Twitch 用户 ID：https://www.streamweasels.com/tools/convert-twitch-username-to-user-id/
  \item 配置令牌：
\end{enumerate}

\begin{itemize}
  \item 环境变量：\texttt{OPENCLAW\_TWITCH\_ACCESS\_TOKEN=...}（仅限默认账户）
  \item 或配置：\texttt{channels.twitch.accessToken}
  \item 如果两者都设置，配置优先（环境变量回退仅适用于默认账户）。
\end{itemize}

\begin{enumerate}
  \item 启动 Gateway 网关。
\end{enumerate}


\textbf{\emoji{⚠} 重要：} 添加访问控制（\texttt{allowFrom} 或 \texttt{allowedRoles}）以防止未授权用户触发机器人。\texttt{requireMention} 默认为 \texttt{true}。

最小配置：

\begin{lstlisting}[language=json]
{
  channels: {
    twitch: {
      enabled: true,
      username: "openclaw", // 机器人的 Twitch 账户
      accessToken: "oauth:abc123...", // OAuth Access Token（或使用 OPENCLAW_TWITCH_ACCESS_TOKEN 环境变量）
      clientId: "xyz789...", // Token Generator 中的 Client ID
      channel: "vevisk", // 要加入的 Twitch 频道聊天（必填）
      allowFrom: ["123456789"], // （推荐）仅限你的 Twitch 用户 ID - 从 https://www.streamweasels.com/tools/convert-twitch-username-to-user-id/ 获取
    },
  },
}
\end{lstlisting}


\subsection{它是什么}


\begin{itemize}
  \item 由 Gateway 网关拥有的 Twitch 渠道。
  \item 确定性路由：回复总是返回到 Twitch。
  \item 每个账户映射到一个隔离的会话键 \texttt{agent::twitch:}。
  \item \texttt{username} 是机器人账户（进行身份验证的账户），\texttt{channel} 是要加入的聊天室。
\end{itemize}


\subsection{设置（详细）}


\subsubsection{生成凭证}


使用 \href{https://twitchtokengenerator.com/}{Twitch Token Generator}：

\begin{itemize}
  \item 选择 \textbf{Bot Token}
  \item 确认已选择 \texttt{chat:read} 和 \texttt{chat:write} 权限范围
  \item 复制 \textbf{Client ID} 和 \textbf{Access Token}
\end{itemize}


无需手动注册应用。令牌在几小时后过期。

\subsubsection{配置机器人}


\textbf{环境变量（仅限默认账户）：}

\begin{lstlisting}[language=bash]
OPENCLAW_TWITCH_ACCESS_TOKEN=oauth:abc123...
\end{lstlisting}


\textbf{或配置：}

\begin{lstlisting}[language=json]
{
  channels: {
    twitch: {
      enabled: true,
      username: "openclaw",
      accessToken: "oauth:abc123...",
      clientId: "xyz789...",
      channel: "vevisk",
    },
  },
}
\end{lstlisting}


如果环境变量和配置都设置了，配置优先。

\subsubsection{访问控制（推荐）}


\begin{lstlisting}[language=json]
{
  channels: {
    twitch: {
      allowFrom: ["123456789"], // （推荐）仅限你的 Twitch 用户 ID
    },
  },
}
\end{lstlisting}


优先使用 \texttt{allowFrom} 作为硬性允许列表。如果你想要基于角色的访问控制，请改用 \texttt{allowedRoles}。

\textbf{可用角色：} \texttt{"moderator"}、\texttt{"owner"}、\texttt{"vip"}、\texttt{"subscriber"}、\texttt{"all"}。

\textbf{为什么用用户 ID？} 用户名可以更改，允许冒充。用户 ID 是永久的。

查找你的 Twitch 用户 ID：https://www.streamweasels.com/tools/convert-twitch-username-\%20to-user-id/（将你的 Twitch 用户名转换为 ID）

\subsection{令牌刷新（可选）}


来自 \href{https://twitchtokengenerator.com/}{Twitch Token Generator} 的令牌无法自动刷新 - 过期时需要重新生成。

要实现自动令牌刷新，请在 \href{https://dev.twitch.tv/console}{Twitch Developer Console} 创建你自己的 Twitch 应用并添加到配置中：

\begin{lstlisting}[language=json]
{
  channels: {
    twitch: {
      clientSecret: "your_client_secret",
      refreshToken: "your_refresh_token",
    },
  },
}
\end{lstlisting}


机器人会在令牌过期前自动刷新，并记录刷新事件。

\subsection{多账户支持}


使用 \texttt{channels.twitch.accounts} 配置每个账户的令牌。参阅 \href{/gateway/configuration}{\texttt{gateway/configuration}} 了解共享模式。

示例（一个机器人账户在两个频道中）：

\begin{lstlisting}[language=json]
{
  channels: {
    twitch: {
      accounts: {
        channel1: {
          username: "openclaw",
          accessToken: "oauth:abc123...",
          clientId: "xyz789...",
          channel: "vevisk",
        },
        channel2: {
          username: "openclaw",
          accessToken: "oauth:def456...",
          clientId: "uvw012...",
          channel: "secondchannel",
        },
      },
    },
  },
}
\end{lstlisting}


\textbf{注意：} 每个账户需要自己的令牌（每个频道一个令牌）。

\subsection{访问控制}


\subsubsection{基于角色的限制}


\begin{lstlisting}[language=json]
{
  channels: {
    twitch: {
      accounts: {
        default: {
          allowedRoles: ["moderator", "vip"],
        },
      },
    },
  },
}
\end{lstlisting}


\subsubsection{按用户 ID 允许列表（最安全）}


\begin{lstlisting}[language=json]
{
  channels: {
    twitch: {
      accounts: {
        default: {
          allowFrom: ["123456789", "987654321"],
        },
      },
    },
  },
}
\end{lstlisting}


\subsubsection{基于角色的访问（替代方案）}


\texttt{allowFrom} 是硬性允许列表。设置后，只允许这些用户 ID。
如果你想要基于角色的访问，请不设置 \texttt{allowFrom}，改为配置 \texttt{allowedRoles}：

\begin{lstlisting}[language=json]
{
  channels: {
    twitch: {
      accounts: {
        default: {
          allowedRoles: ["moderator"],
        },
      },
    },
  },
}
\end{lstlisting}


\subsubsection{禁用 @提及要求}


默认情况下，\texttt{requireMention} 为 \texttt{true}。要禁用并响应所有消息：

\begin{lstlisting}[language=json]
{
  channels: {
    twitch: {
      accounts: {
        default: {
          requireMention: false,
        },
      },
    },
  },
}
\end{lstlisting}


\subsection{故障排除}


首先，运行诊断命令：

\begin{lstlisting}[language=bash]
openclaw doctor
openclaw channels status --probe
\end{lstlisting}


\subsubsection{机器人不响应消息}


\textbf{检查访问控制：} 确保你的用户 ID 在 \texttt{allowFrom} 中，或临时移除 \texttt{allowFrom} 并设置 \texttt{allowedRoles: ["all"]} 来测试。

\textbf{检查机器人是否在频道中：} 机器人必须加入 \texttt{channel} 中指定的频道。

\subsubsection{令牌问题}


\textbf{"Failed to connect"或身份验证错误：}

\begin{itemize}
  \item 验证 \texttt{accessToken} 是 OAuth 访问令牌值（通常以 \texttt{oauth:} 前缀开头）
  \item 检查令牌具有 \texttt{chat:read} 和 \texttt{chat:write} 权限范围
  \item 如果使用令牌刷新，验证 \texttt{clientSecret} 和 \texttt{refreshToken} 已设置
\end{itemize}


\subsubsection{令牌刷新不工作}


\textbf{检查日志中的刷新事件：}

\begin{lstlisting}
Using env token source for mybot
Access token refreshed for user 123456 (expires in 14400s)
\end{lstlisting}


如果你看到"token refresh disabled (no refresh token)"：

\begin{itemize}
  \item 确保提供了 \texttt{clientSecret}
  \item 确保提供了 \texttt{refreshToken}
\end{itemize}


\subsection{配置}


\textbf{账户配置：}

\begin{itemize}
  \item \texttt{username} - 机器人用户名
  \item \texttt{accessToken} - 具有 \texttt{chat:read} 和 \texttt{chat:write} 权限的 OAuth 访问令牌
  \item \texttt{clientId} - Twitch Client ID（来自 Token Generator 或你的应用）
  \item \texttt{channel} - 要加入的频道（必填）
  \item \texttt{enabled} - 启用此账户（默认：\texttt{true}）
  \item \texttt{clientSecret} - 可选：用于自动令牌刷新
  \item \texttt{refreshToken} - 可选：用于自动令牌刷新
  \item \texttt{expiresIn} - 令牌过期时间（秒）
  \item \texttt{obtainmentTimestamp} - 令牌获取时间戳
  \item \texttt{allowFrom} - 用户 ID 允许列表
  \item \texttt{allowedRoles} - 基于角色的访问控制（\texttt{"moderator" | "owner" | "vip" | "subscriber" | "all"}）
  \item \texttt{requireMention} - 需要 @提及（默认：\texttt{true}）
\end{itemize}


\textbf{提供商选项：}

\begin{itemize}
  \item \texttt{channels.twitch.enabled} - 启用/禁用渠道启动
  \item \texttt{channels.twitch.username} - 机器人用户名（简化的单账户配置）
  \item \texttt{channels.twitch.accessToken} - OAuth 访问令牌（简化的单账户配置）
  \item \texttt{channels.twitch.clientId} - Twitch Client ID（简化的单账户配置）
  \item \texttt{channels.twitch.channel} - 要加入的频道（简化的单账户配置）
  \item \texttt{channels.twitch.accounts.} - 多账户配置（以上所有账户字段）
\end{itemize}


完整示例：

\begin{lstlisting}[language=json]
{
  channels: {
    twitch: {
      enabled: true,
      username: "openclaw",
      accessToken: "oauth:abc123...",
      clientId: "xyz789...",
      channel: "vevisk",
      clientSecret: "secret123...",
      refreshToken: "refresh456...",
      allowFrom: ["123456789"],
      allowedRoles: ["moderator", "vip"],
      accounts: {
        default: {
          username: "mybot",
          accessToken: "oauth:abc123...",
          clientId: "xyz789...",
          channel: "your_channel",
          enabled: true,
          clientSecret: "secret123...",
          refreshToken: "refresh456...",
          expiresIn: 14400,
          obtainmentTimestamp: 1706092800000,
          allowFrom: ["123456789", "987654321"],
          allowedRoles: ["moderator"],
        },
      },
    },
  },
}
\end{lstlisting}


\subsection{工具操作}


智能体可以调用 \texttt{twitch} 执行以下操作：

\begin{itemize}
  \item \texttt{send} - 向频道发送消息
\end{itemize}


示例：

\begin{lstlisting}[language=json]
{
  action: "twitch",
  params: {
    message: "Hello Twitch!",
    to: "#mychannel",
  },
}
\end{lstlisting}


\subsection{安全与运维}


\begin{itemize}
  \item \textbf{将令牌视为密码} - 永远不要将令牌提交到 git
  \item \textbf{使用自动令牌刷新} 用于长时间运行的机器人
  \item \textbf{使用用户 ID 允许列表} 而不是用户名进行访问控制
  \item \textbf{监控日志} 查看令牌刷新事件和连接状态
  \item \textbf{最小化令牌权限范围} - 只请求 \texttt{chat:read} 和 \texttt{chat:write}
  \item \textbf{如果卡住}：在确认没有其他进程拥有会话后重启 Gateway 网关
\end{itemize}


\subsection{限制}


\begin{itemize}
  \item 每条消息 \textbf{500 个字符}（在单词边界自动分块）
  \item 分块前会去除 Markdown
  \item 无速率限制（使用 Twitch 内置的速率限制）
\end{itemize}



\clearpage

% === 源文件: channels/whatsapp.md ===
\section{WhatsApp（网页渠道）}


状态：仅支持通过 Baileys 的 WhatsApp Web。Gateway 网关拥有会话。

\subsection{快速设置（新手）}


\begin{enumerate}
  \item 如果可能，使用\textbf{单独的手机号码}（推荐）。
  \item 在 \texttt{\textasciitilde{}/.openclaw/openclaw.json} 中配置 WhatsApp。
  \item 运行 \texttt{openclaw channels login} 扫描二维码（关联设备）。
  \item 启动 Gateway 网关。
\end{enumerate}


最小配置：

\begin{lstlisting}[language=json]
{
  channels: {
    whatsapp: {
      dmPolicy: "allowlist",
      allowFrom: ["+15551234567"],
    },
  },
}
\end{lstlisting}


\subsection{目标}


\begin{itemize}
  \item 在一个 Gateway 网关进程中支持多个 WhatsApp 账户（多账户）。
  \item 确定性路由：回复返回到 WhatsApp，无模型路由。
  \item 模型能看到足够的上下文来理解引用回复。
\end{itemize}


\subsection{配置写入}


默认情况下，WhatsApp 允许写入由 \texttt{/config set|unset} 触发的配置更新（需要 \texttt{commands.config: true}）。

禁用方式：

\begin{lstlisting}[language=json]
{
  channels: { whatsapp: { configWrites: false } },
}
\end{lstlisting}


\subsection{架构（谁拥有什么）}


\begin{itemize}
  \item \textbf{Gateway 网关}拥有 Baileys socket 和收件箱循环。
  \item \textbf{CLI / macOS 应用}与 Gateway 网关通信；不直接使用 Baileys。
  \item 发送出站消息需要\textbf{活跃的监听器}；否则发送会快速失败。
\end{itemize}


\subsection{获取手机号码（两种模式）}


WhatsApp 需要真实手机号码进行验证。VoIP 和虚拟号码通常会被封锁。在 WhatsApp 上运行 OpenClaw 有两种支持的方式：

\subsubsection{专用号码（推荐）}


为 OpenClaw 使用\textbf{单独的手机号码}。最佳用户体验，清晰的路由，无自聊天怪异问题。理想设置：\textbf{备用/旧 Android 手机 + eSIM}。保持 Wi-Fi 和电源连接，通过二维码关联。

\textbf{WhatsApp Business：} 你可以在同一设备上使用不同号码的 WhatsApp Business。非常适合将个人 WhatsApp 分开——安装 WhatsApp Business 并在那里注册 OpenClaw 号码。

\textbf{示例配置（专用号码，单用户允许列表）：}

\begin{lstlisting}[language=json]
{
  channels: {
    whatsapp: {
      dmPolicy: "allowlist",
      allowFrom: ["+15551234567"],
    },
  },
}
\end{lstlisting}


\textbf{配对模式（可选）：}
如果你想使用配对而不是允许列表，请将 \texttt{channels.whatsapp.dmPolicy} 设置为 \texttt{pairing}。未知发送者会收到配对码；使用以下命令批准：
\texttt{openclaw pairing approve whatsapp }

\subsubsection{个人号码（备选方案）}


快速备选方案：在\textbf{你自己的号码}上运行 OpenClaw。给自己发消息（WhatsApp"给自己发消息"）进行测试，这样就不会打扰联系人。在设置和实验期间需要在主手机上阅读验证码。\textbf{必须启用自聊天模式。}
当向导询问你的个人 WhatsApp 号码时，输入你将用于发送消息的手机（所有者/发送者），而不是助手号码。

\textbf{示例配置（个人号码，自聊天）：}

\begin{lstlisting}[language=json]
{
  "whatsapp": {
    "selfChatMode": true,
    "dmPolicy": "allowlist",
    "allowFrom": ["+15551234567"]
  }
}
\end{lstlisting}


当设置了 \texttt{identity.name} 时，自聊天回复默认为 \texttt{[\{identity.name\}]}（否则为 \texttt{[openclaw]}），
前提是 \texttt{messages.responsePrefix} 未设置。明确设置它可以自定义或禁用
前缀（使用 \texttt{""} 来移除）。

\subsubsection{号码获取提示}


\begin{itemize}
  \item \textbf{本地 eSIM} 来自你所在国家的移动运营商（最可靠）
  \item 奥地利：\href{https://www.hot.at}{hot.at}
  \item 英国：\href{https://www.giffgaff.com}{giffgaff} — 免费 SIM 卡，无合约
  \item \textbf{预付费 SIM 卡} — 便宜，只需接收一条验证短信
\end{itemize}


\textbf{避免：} TextNow、Google Voice、大多数"免费短信"服务——WhatsApp 会积极封锁这些。

\textbf{提示：} 该号码只需要接收一条验证短信。之后，WhatsApp Web 会话通过 \texttt{creds.json} 持久化。

\subsection{为什么不用 Twilio？}


\begin{itemize}
  \item 早期 OpenClaw 版本支持 Twilio 的 WhatsApp Business 集成。
  \item WhatsApp Business 号码不适合个人助手。
  \item Meta 强制执行 24 小时回复窗口；如果你在过去 24 小时内没有回复，商业号码无法发起新消息。
  \item 高频或"频繁"使用会触发激进的封锁，因为商业账户不适合发送大量个人助手消息。
  \item 结果：投递不可靠且频繁被封锁，因此该支持已被移除。
\end{itemize}


\subsection{登录 + 凭证}


\begin{itemize}
  \item 登录命令：\texttt{openclaw channels login}（通过关联设备扫描二维码）。
  \item 多账户登录：\texttt{openclaw channels login --account }（`\texttt{ = }accountId`）。
  \item 默认账户（省略 \texttt{--account} 时）：如果存在则为 \texttt{default}，否则为第一个配置的账户 id（排序后）。
  \item 凭证存储在 \texttt{\textasciitilde{}/.openclaw/credentials/whatsapp//creds.json}。
  \item 备份副本在 \texttt{creds.json.bak}（损坏时恢复）。
  \item 旧版兼容性：较旧的安装将 Baileys 文件直接存储在 \texttt{\textasciitilde{}/.openclaw/credentials/} 中。
  \item 登出：\texttt{openclaw channels logout}（或 \texttt{--account }）删除 WhatsApp 认证状态（但保留共享的 \texttt{oauth.json}）。
  \item 已登出的 socket => 错误提示重新关联。
\end{itemize}


\subsection{入站流程（私信 + 群组）}


\begin{itemize}
  \item WhatsApp 事件来自 \texttt{messages.upsert}（Baileys）。
  \item 收件箱监听器在关闭时分离，以避免在测试/重启时累积事件处理器。
  \item 状态/广播聊天被忽略。
  \item 直接聊天使用 E.164；群组使用群组 JID。
  \item \textbf{私信策略}：\texttt{channels.whatsapp.dmPolicy} 控制直接聊天访问（默认：\texttt{pairing}）。
  \item 配对：未知发送者会收到配对码（通过 \texttt{openclaw pairing approve whatsapp } 批准；码在 1 小时后过期）。
  \item 开放：需要 \texttt{channels.whatsapp.allowFrom} 包含 \texttt{"*"}。
  \item 你关联的 WhatsApp 号码是隐式信任的，因此自身消息会跳过 \texttt{channels.whatsapp.dmPolicy} 和 \texttt{channels.whatsapp.allowFrom} 检查。
\end{itemize}


\subsubsection{个人号码模式（备选方案）}


如果你在\textbf{个人 WhatsApp 号码}上运行 OpenClaw，请启用 \texttt{channels.whatsapp.selfChatMode}（见上面的示例）。

行为：

\begin{itemize}
  \item 出站私信永远不会触发配对回复（防止打扰联系人）。
  \item 入站未知发送者仍遵循 \texttt{channels.whatsapp.dmPolicy}。
  \item 自聊天模式（allowFrom 包含你的号码）避免自动已读回执并忽略提及 JID。
  \item 非自聊天私信会发送已读回执。
\end{itemize}


\subsection{已读回执}


默认情况下，Gateway 网关在接受入站 WhatsApp 消息后将其标记为已读（蓝色勾号）。

全局禁用：

\begin{lstlisting}[language=json]
{
  channels: { whatsapp: { sendReadReceipts: false } },
}
\end{lstlisting}


按账户禁用：

\begin{lstlisting}[language=json]
{
  channels: {
    whatsapp: {
      accounts: {
        personal: { sendReadReceipts: false },
      },
    },
  },
}
\end{lstlisting}


注意事项：

\begin{itemize}
  \item 自聊天模式始终跳过已读回执。
\end{itemize}


\subsection{WhatsApp 常见问题：发送消息 + 配对}


\textbf{当我关联 WhatsApp 时，OpenClaw 会给随机联系人发消息吗？}
不会。默认私信策略是\textbf{配对}，因此未知发送者只会收到配对码，他们的消息\textbf{不会被处理}。OpenClaw 只会回复它收到的聊天，或你明确触发的发送（智能体/CLI）。

\textbf{WhatsApp 上的配对是如何工作的？}
配对是未知发送者的私信门控：

\begin{itemize}
  \item 来自新发送者的第一条私信返回一个短码（消息不会被处理）。
  \item 使用以下命令批准：\texttt{openclaw pairing approve whatsapp }（使用 \texttt{openclaw pairing list whatsapp} 列出）。
  \item 码在 1 小时后过期；每个渠道的待处理请求上限为 3 个。
\end{itemize}


\textbf{多个人可以在一个 WhatsApp 号码上使用不同的 OpenClaw 实例吗？}
可以，通过 \texttt{bindings} 将每个发送者路由到不同的智能体（peer \texttt{kind: "dm"}，发送者 E.164 如 \texttt{+15551234567}）。回复仍然来自\textbf{同一个 WhatsApp 账户}，直接聊天会折叠到每个智能体的主会话，因此\textbf{每人使用一个智能体}。私信访问控制（\texttt{dmPolicy}/\texttt{allowFrom}）是每个 WhatsApp 账户全局的。参见\href{/concepts/multi-agent}{多智能体路由}。

\textbf{为什么向导会询问我的手机号码？}
向导使用它来设置你的\textbf{允许列表/所有者}，以便允许你自己的私信。它不会用于自动发送。如果你在个人 WhatsApp 号码上运行，请使用相同的号码并启用 \texttt{channels.whatsapp.selfChatMode}。

\subsection{消息规范化（模型看到的内容）}


\begin{itemize}
  \item \texttt{Body} 是带有信封的当前消息正文。
  \item 引用回复上下文\textbf{始终附加}：
\begin{lstlisting}
  [Replying to +1555 id:ABC123]
  <quoted text or <media:...>>
  [/Replying]
\end{lstlisting}

  \item 回复元数据也会设置：
  \item \texttt{ReplyToId} = stanzaId
  \item \texttt{ReplyToBody} = 引用正文或媒体占位符
  \item \texttt{ReplyToSender} = 已知时为 E.164
  \item 纯媒体入站消息使用占位符：
  \item ``
\end{itemize}


\subsection{群组}


\begin{itemize}
  \item 群组映射到 \texttt{agent::whatsapp:group:} 会话。
  \item 群组策略：\texttt{channels.whatsapp.groupPolicy = open|disabled|allowlist}（默认 \texttt{allowlist}）。
  \item 激活模式：
  \item \texttt{mention}（默认）：需要 @提及或正则匹配。
  \item \texttt{always}：始终触发。
  \item \texttt{/activation mention|always} 仅限所有者，必须作为独立消息发送。
  \item 所有者 = \texttt{channels.whatsapp.allowFrom}（如果未设置则为自身 E.164）。
  \item \textbf{历史注入}（仅待处理）：
  \item 最近\textit{未处理}的消息（默认 50 条）插入在：
\end{itemize}

\texttt{[Chat messages since your last reply - for context]}（已在会话中的消息不会重新注入）
\begin{itemize}
  \item 当前消息在：
\end{itemize}

\texttt{[Current message - respond to this]}
\begin{itemize}
  \item 附加发送者后缀：\texttt{[from: Name (+E164)]}
  \item 群组元数据缓存 5 分钟（主题 + 参与者）。
\end{itemize}


\subsection{回复投递（线程）}


\begin{itemize}
  \item WhatsApp Web 发送标准消息（当前 Gateway 网关无引用回复线程）。
  \item 此渠道忽略回复标签。
\end{itemize}


\subsection{确认表情（收到时自动回应）}


WhatsApp 可以在收到传入消息时立即自动发送表情回应，在机器人生成回复之前。这为用户提供即时反馈，表明他们的消息已收到。

\textbf{配置：}

\begin{lstlisting}[language=json]
{
  "whatsapp": {
    "ackReaction": {
      "emoji": "👀",
      "direct": true,
      "group": "mentions"
    }
  }
}
\end{lstlisting}


\textbf{选项：}

\begin{itemize}
  \item \texttt{emoji}（字符串）：用于确认的表情（例如"\emoji{👀}"、"\emoji{✅}"、"\emoji{📨}"）。为空或省略 = 功能禁用。
  \item \texttt{direct}（布尔值，默认：\texttt{true}）：在直接/私信聊天中发送表情回应。
  \item \texttt{group}（字符串，默认：\texttt{"mentions"}）：群聊行为：
  \item \texttt{"always"}：对所有群消息做出回应（即使没有 @提及）
  \item \texttt{"mentions"}：仅在机器人被 @提及时做出回应
  \item \texttt{"never"}：从不在群组中做出回应
\end{itemize}


\textbf{按账户覆盖：}

\begin{lstlisting}[language=json]
{
  "whatsapp": {
    "accounts": {
      "work": {
        "ackReaction": {
          "emoji": "✅",
          "direct": false,
          "group": "always"
        }
      }
    }
  }
}
\end{lstlisting}


\textbf{行为说明：}

\begin{itemize}
  \item 表情回应在消息收到时\textbf{立即}发送，在输入指示器或机器人回复之前。
  \item 在 \texttt{requireMention: false}（激活：always）的群组中，\texttt{group: "mentions"} 会对所有消息做出回应（不仅仅是 @提及）。
  \item 即发即忘：表情回应失败会被记录但不会阻止机器人回复。
  \item 群组表情回应会自动包含参与者 JID。
  \item WhatsApp 忽略 \texttt{messages.ackReaction}；请改用 \texttt{channels.whatsapp.ackReaction}。
\end{itemize}


\subsection{智能体工具（表情回应）}


\begin{itemize}
  \item 工具：\texttt{whatsapp}，带有 \texttt{react} 动作（\texttt{chatJid}、\texttt{messageId}、\texttt{emoji}，可选 \texttt{remove}）。
  \item 可选：\texttt{participant}（群组发送者）、\texttt{fromMe}（对自己的消息做出回应）、\texttt{accountId}（多账户）。
  \item 表情移除语义：参见 \href{/tools/reactions}{/tools/reactions}。
  \item 工具门控：\texttt{channels.whatsapp.actions.reactions}（默认：启用）。
\end{itemize}


\subsection{限制}


\begin{itemize}
  \item 出站文本按 \texttt{channels.whatsapp.textChunkLimit} 分块（默认 4000）。
  \item 可选换行分块：设置 \texttt{channels.whatsapp.chunkMode="newline"} 在长度分块之前按空行（段落边界）分割。
  \item 入站媒体保存受 \texttt{channels.whatsapp.mediaMaxMb} 限制（默认 50 MB）。
  \item 出站媒体项受 \texttt{agents.defaults.mediaMaxMb} 限制（默认 5 MB）。
\end{itemize}


\subsection{出站发送（文本 + 媒体）}


\begin{itemize}
  \item 使用活跃的网页监听器；如果 Gateway 网关未运行则报错。
  \item 文本分块：每条消息最大 4k（可通过 \texttt{channels.whatsapp.textChunkLimit} 配置，可选 \texttt{channels.whatsapp.chunkMode}）。
  \item 媒体：
  \item 支持图片/视频/音频/文档。
  \item 音频作为 PTT 发送；\texttt{audio/ogg} => \texttt{audio/ogg; codecs=opus}。
  \item 仅在第一个媒体项上添加标题。
  \item 媒体获取支持 HTTP(S) 和本地路径。
  \item 动画 GIF：WhatsApp 期望带有 \texttt{gifPlayback: true} 的 MP4 以实现内联循环。
  \item CLI：\texttt{openclaw message send --media  --gif-playback}
  \item Gateway 网关：\texttt{send} 参数包含 \texttt{gifPlayback: true}
\end{itemize}


\subsection{语音消息（PTT 音频）}


WhatsApp 将音频作为\textbf{语音消息}（PTT 气泡）发送。

\begin{itemize}
  \item 最佳效果：OGG/Opus。OpenClaw 将 \texttt{audio/ogg} 重写为 \texttt{audio/ogg; codecs=opus}。
  \item WhatsApp 忽略 \texttt{[[audio\_as\_voice]]}（音频已作为语音消息发送）。
\end{itemize}


\subsection{媒体限制 + 优化}


\begin{itemize}
  \item 默认出站上限：5 MB（每个媒体项）。
  \item 覆盖：\texttt{agents.defaults.mediaMaxMb}。
  \item 图片自动优化为上限以下的 JPEG（调整大小 + 质量扫描）。
  \item 超大媒体 => 错误；媒体回复降级为文本警告。
\end{itemize}


\subsection{心跳}


\begin{itemize}
  \item \textbf{Gateway 网关心跳}记录连接健康状态（\texttt{web.heartbeatSeconds}，默认 60 秒）。
  \item \textbf{智能体心跳}可以按智能体配置（\texttt{agents.list[].heartbeat}）或通过
\end{itemize}

\texttt{agents.defaults.heartbeat} 全局配置（当没有设置按智能体条目时的降级）。
\begin{itemize}
  \item 使用配置的心跳提示词（默认：\texttt{Read HEARTBEAT.md if it exists (workspace context). Follow it strictly. Do not infer or repeat old tasks from prior chats. If nothing needs attention, reply HEARTBEAT\_OK.}）+ \texttt{HEARTBEAT\_OK} 跳过行为。
  \item 投递默认为最后使用的渠道（或配置的目标）。
\end{itemize}


\subsection{重连行为}


\begin{itemize}
  \item 退避策略：\texttt{web.reconnect}：
  \item \texttt{initialMs}、\texttt{maxMs}、\texttt{factor}、\texttt{jitter}、\texttt{maxAttempts}。
  \item 如果达到 maxAttempts，网页监控停止（降级）。
  \item 已登出 => 停止并要求重新关联。
\end{itemize}


\subsection{配置快速映射}


\begin{itemize}
  \item \texttt{channels.whatsapp.dmPolicy}（私信策略：pairing/allowlist/open/disabled）。
  \item \texttt{channels.whatsapp.selfChatMode}（同手机设置；机器人使用你的个人 WhatsApp 号码）。
  \item \texttt{channels.whatsapp.allowFrom}（私信允许列表）。WhatsApp 使用 E.164 手机号码（无用户名）。
  \item \texttt{channels.whatsapp.mediaMaxMb}（入站媒体保存上限）。
  \item \texttt{channels.whatsapp.ackReaction}（消息收到时的自动回应：\texttt{\{emoji, direct, group\}}）。
  \item \texttt{channels.whatsapp.accounts..*}（按账户设置 + 可选 \texttt{authDir}）。
  \item \texttt{channels.whatsapp.accounts..mediaMaxMb}（按账户入站媒体上限）。
  \item \texttt{channels.whatsapp.accounts..ackReaction}（按账户确认回应覆盖）。
  \item \texttt{channels.whatsapp.groupAllowFrom}（群组发送者允许列表）。
  \item \texttt{channels.whatsapp.groupPolicy}（群组策略）。
  \item \texttt{channels.whatsapp.historyLimit} / \texttt{channels.whatsapp.accounts..historyLimit}（群组历史上下文；\texttt{0} 禁用）。
  \item \texttt{channels.whatsapp.dmHistoryLimit}（私信历史限制，按用户轮次）。按用户覆盖：\texttt{channels.whatsapp.dms[""].historyLimit}。
  \item \texttt{channels.whatsapp.groups}（群组允许列表 + 提及门控默认值；使用 \texttt{"*"} 允许全部）
  \item \texttt{channels.whatsapp.actions.reactions}（门控 WhatsApp 工具表情回应）。
  \item \texttt{agents.list[].groupChat.mentionPatterns}（或 \texttt{messages.groupChat.mentionPatterns}）
  \item \texttt{messages.groupChat.historyLimit}
  \item \texttt{channels.whatsapp.messagePrefix}（入站前缀；按账户：\texttt{channels.whatsapp.accounts..messagePrefix}；已弃用：\texttt{messages.messagePrefix}）
  \item \texttt{messages.responsePrefix}（出站前缀）
  \item \texttt{agents.defaults.mediaMaxMb}
  \item \texttt{agents.defaults.heartbeat.every}
  \item \texttt{agents.defaults.heartbeat.model}（可选覆盖）
  \item \texttt{agents.defaults.heartbeat.target}
  \item \texttt{agents.defaults.heartbeat.to}
  \item \texttt{agents.defaults.heartbeat.session}
  \item \texttt{agents.list[].heartbeat.*}（按智能体覆盖）
  \item \texttt{session.*}（scope、idle、store、mainKey）
  \item \texttt{web.enabled}（为 false 时禁用渠道启动）
  \item \texttt{web.heartbeatSeconds}
  \item \texttt{web.reconnect.*}
\end{itemize}


\subsection{日志 + 故障排除}


\begin{itemize}
  \item 子系统：\texttt{whatsapp/inbound}、\texttt{whatsapp/outbound}、\texttt{web-heartbeat}、\texttt{web-reconnect}。
  \item 日志文件：\texttt{/tmp/openclaw/openclaw-YYYY-MM-DD.log}（可配置）。
  \item 故障排除指南：\href{/gateway/troubleshooting}{Gateway 网关故障排除}。
\end{itemize}


\subsection{故障排除（快速）}


\textbf{未关联 / 需要二维码登录}

\begin{itemize}
  \item 症状：\texttt{channels status} 显示 \texttt{linked: false} 或警告"Not linked"。
  \item 修复：在 Gateway 网关主机上运行 \texttt{openclaw channels login} 并扫描二维码（WhatsApp → 设置 → 关联设备）。
\end{itemize}


\textbf{已关联但断开连接 / 重连循环}

\begin{itemize}
  \item 症状：\texttt{channels status} 显示 \texttt{running, disconnected} 或警告"Linked but disconnected"。
  \item 修复：\texttt{openclaw doctor}（或重启 Gateway 网关）。如果问题持续，通过 \texttt{channels login} 重新关联并检查 \texttt{openclaw logs --follow}。
\end{itemize}


\textbf{Bun 运行时}

\begin{itemize}
  \item \textbf{不推荐} Bun。WhatsApp（Baileys）和 Telegram 在 Bun 上不可靠。
\end{itemize}

请使用 \textbf{Node} 运行 Gateway 网关。（参见入门指南运行时说明。）


\clearpage

% === 源文件: channels/zalo.md ===
\section{Zalo (Bot API)}


状态：实验性。仅支持私信；根据 Zalo 文档，群组即将推出。

\subsection{需要插件}


Zalo 以插件形式提供，不包含在核心安装中。

\begin{itemize}
  \item 通过 CLI 安装：\texttt{openclaw plugins install @openclaw/zalo}
  \item 或在新手引导期间选择 \textbf{Zalo} 并确认安装提示
  \item 详情：\href{/tools/plugin}{插件}
\end{itemize}


\subsection{快速设置（初学者）}


\begin{enumerate}
  \item 安装 Zalo 插件：
\end{enumerate}

\begin{itemize}
  \item 从源代码检出：\texttt{openclaw plugins install ./extensions/zalo}
  \item 从 npm（如果已发布）：\texttt{openclaw plugins install @openclaw/zalo}
  \item 或在新手引导中选择 \textbf{Zalo} 并确认安装提示
\end{itemize}

\begin{enumerate}
  \item 设置 token：
\end{enumerate}

\begin{itemize}
  \item 环境变量：\texttt{ZALO\_BOT\_TOKEN=...}
  \item 或配置：\texttt{channels.zalo.botToken: "..."}。
\end{itemize}

\begin{enumerate}
  \item 重启 Gateway 网关（或完成新手引导）。
  \item 私信访问默认为配对模式；首次联系时批准配对码。
\end{enumerate}


最小配置：

\begin{lstlisting}[language=json]
{
  channels: {
    zalo: {
      enabled: true,
      botToken: "12345689:abc-xyz",
      dmPolicy: "pairing",
    },
  },
}
\end{lstlisting}


\subsection{它是什么}


Zalo 是一款专注于越南市场的即时通讯应用；其 Bot API 让 Gateway 网关可以运行一个用于一对一对话的 bot。
它非常适合需要确定性路由回 Zalo 的支持或通知场景。

\begin{itemize}
  \item 由 Gateway 网关拥有的 Zalo Bot API 渠道。
  \item 确定性路由：回复返回到 Zalo；模型不会选择渠道。
  \item 私信共享智能体的主会话。
  \item 群组尚不支持（Zalo 文档标注"即将推出"）。
\end{itemize}


\subsection{设置（快速路径）}


\subsubsection{1）创建 bot token（Zalo Bot 平台）}


\begin{enumerate}
  \item 前往 \textbf{https://bot.zaloplatforms.com} 并登录。
  \item 创建新 bot 并配置其设置。
  \item 复制 bot token（格式：\texttt{12345689:abc-xyz}）。
\end{enumerate}


\subsubsection{2）配置 token（环境变量或配置）}


示例：

\begin{lstlisting}[language=json]
{
  channels: {
    zalo: {
      enabled: true,
      botToken: "12345689:abc-xyz",
      dmPolicy: "pairing",
    },
  },
}
\end{lstlisting}


环境变量选项：\texttt{ZALO\_BOT\_TOKEN=...}（仅适用于默认账户）。

多账户支持：使用 \texttt{channels.zalo.accounts} 配置每账户 token 和可选的 \texttt{name}。

\begin{enumerate}
  \item 重启 Gateway 网关。当 token 被解析（环境变量或配置）时，Zalo 启动。
  \item 私信访问默认为配对模式。当 bot 首次被联系时批准配对码。
\end{enumerate}


\subsection{工作原理（行为）}


\begin{itemize}
  \item 入站消息被规范化为带有媒体占位符的共享渠道信封。
  \item 回复始终路由回同一 Zalo 聊天。
  \item 默认使用长轮询；可通过 \texttt{channels.zalo.webhookUrl} 启用 webhook 模式。
\end{itemize}


\subsection{限制}


\begin{itemize}
  \item 出站文本按 2000 字符分块（Zalo API 限制）。
  \item 媒体下载/上传受 \texttt{channels.zalo.mediaMaxMb} 限制（默认 5）。
  \item 由于 2000 字符限制使流式传输效果不佳，默认阻止流式传输。
\end{itemize}


\subsection{访问控制（私信）}


\subsubsection{私信访问}


\begin{itemize}
  \item 默认：\texttt{channels.zalo.dmPolicy = "pairing"}。未知发送者会收到配对码；消息在批准前会被忽略（配对码 1 小时后过期）。
  \item 通过以下方式批准：
  \item \texttt{openclaw pairing list zalo}
  \item \texttt{openclaw pairing approve zalo }
  \item 配对是默认的令牌交换方式。详情：\href{/channels/pairing}{配对}
  \item \texttt{channels.zalo.allowFrom} 接受数字用户 ID（无用户名查找功能）。
\end{itemize}


\subsection{长轮询与 webhook}


\begin{itemize}
  \item 默认：长轮询（不需要公共 URL）。
  \item Webhook 模式：设置 \texttt{channels.zalo.webhookUrl} 和 \texttt{channels.zalo.webhookSecret}。
  \item Webhook secret 必须为 8-256 个字符。
  \item Webhook URL 必须使用 HTTPS。
  \item Zalo 发送事件时带有 \texttt{X-Bot-Api-Secret-Token} 头用于验证。
  \item Gateway 网关 HTTP 在 \texttt{channels.zalo.webhookPath} 处理 webhook 请求（默认为 webhook URL 路径）。
\end{itemize}


\textbf{注意：} 根据 Zalo API 文档，getUpdates（轮询）和 webhook 是互斥的。

\subsection{支持的消息类型}


\begin{itemize}
  \item \textbf{文本消息}：完全支持，2000 字符分块。
  \item \textbf{图片消息}：下载和处理入站图片；通过 \texttt{sendPhoto} 发送图片。
  \item \textbf{贴纸}：已记录但未完全处理（无智能体响应）。
  \item \textbf{不支持的类型}：已记录（例如来自受保护用户的消息）。
\end{itemize}


\subsection{功能}



\begin{table}[h]
\centering
\small
\begin{tabular}{|l|l|}
\hline
功能 & 状态 \\
\hline
私信 & \emoji{✅} 支持 \\
\hline
群组 & \emoji{❌} 即将推出（根据 Zalo 文档） \\
\hline
媒体（图片） & \emoji{✅} 支持 \\
\hline
表情回应 & \emoji{❌} 不支持 \\
\hline
主题 & \emoji{❌} 不支持 \\
\hline
投票 & \emoji{❌} 不支持 \\
\hline
原生命令 & \emoji{❌} 不支持 \\
\hline
流式传输 & \emoji{⚠} 已阻止（2000 字符限制） \\
\hline
\end{tabular}
\end{table}



\subsection{投递目标（CLI/cron）}


\begin{itemize}
  \item 使用聊天 id 作为目标。
  \item 示例：\texttt{openclaw message send --channel zalo --target 123456789 --message "hi"}。
\end{itemize}


\subsection{故障排除}


\textbf{Bot 不响应：}

\begin{itemize}
  \item 检查 token 是否有效：\texttt{openclaw channels status --probe}
  \item 验证发送者已被批准（配对或 allowFrom）
  \item 检查 Gateway 网关日志：\texttt{openclaw logs --follow}
\end{itemize}


\textbf{Webhook 未收到事件：}

\begin{itemize}
  \item 确保 webhook URL 使用 HTTPS
  \item 验证 secret token 为 8-256 个字符
  \item 确认 Gateway 网关 HTTP 端点在配置的路径上可访问
  \item 检查 getUpdates 轮询未在运行（它们是互斥的）
\end{itemize}


\subsection{配置参考（Zalo）}


完整配置：\href{/gateway/configuration}{配置}

提供商选项：

\begin{itemize}
  \item \texttt{channels.zalo.enabled}：启用/禁用渠道启动。
  \item \texttt{channels.zalo.botToken}：来自 Zalo Bot 平台的 bot token。
  \item \texttt{channels.zalo.tokenFile}：从文件路径读取 token。
  \item \texttt{channels.zalo.dmPolicy}：\texttt{pairing | allowlist | open | disabled}（默认：pairing）。
  \item \texttt{channels.zalo.allowFrom}：私信允许列表（用户 ID）。\texttt{open} 需要 \texttt{"*"}。向导会询问数字 ID。
  \item \texttt{channels.zalo.mediaMaxMb}：入站/出站媒体上限（MB，默认 5）。
  \item \texttt{channels.zalo.webhookUrl}：启用 webhook 模式（需要 HTTPS）。
  \item \texttt{channels.zalo.webhookSecret}：webhook secret（8-256 字符）。
  \item \texttt{channels.zalo.webhookPath}：Gateway 网关 HTTP 服务器上的 webhook 路径。
  \item \texttt{channels.zalo.proxy}：API 请求的代理 URL。
\end{itemize}


多账户选项：

\begin{itemize}
  \item \texttt{channels.zalo.accounts..botToken}：每账户 token。
  \item \texttt{channels.zalo.accounts..tokenFile}：每账户 token 文件。
  \item \texttt{channels.zalo.accounts..name}：显示名称。
  \item \texttt{channels.zalo.accounts..enabled}：启用/禁用账户。
  \item \texttt{channels.zalo.accounts..dmPolicy}：每账户私信策略。
  \item \texttt{channels.zalo.accounts..allowFrom}：每账户允许列表。
  \item \texttt{channels.zalo.accounts..webhookUrl}：每账户 webhook URL。
  \item \texttt{channels.zalo.accounts..webhookSecret}：每账户 webhook secret。
  \item \texttt{channels.zalo.accounts..webhookPath}：每账户 webhook 路径。
  \item \texttt{channels.zalo.accounts..proxy}：每账户代理 URL。
\end{itemize}



\clearpage

% === 源文件: channels/zalouser.md ===
\section{Zalo Personal（非官方）}


状态：实验性。此集成通过 \texttt{zca-cli} 自动化\textbf{个人 Zalo 账户}。

\begin{quote}
\textbf{警告：}这是一个非官方集成，可能导致账户被暂停/封禁。使用风险自负。
\end{quote}


\subsection{需要插件}


Zalo Personal 作为插件提供，不包含在核心安装中。

\begin{itemize}
  \item 通过 CLI 安装：\texttt{openclaw plugins install @openclaw/zalouser}
  \item 或从源码检出安装：\texttt{openclaw plugins install ./extensions/zalouser}
  \item 详情：\href{/tools/plugin}{插件}
\end{itemize}


\subsection{前置条件：zca-cli}


Gateway 网关机器必须在 \texttt{PATH} 中有可用的 \texttt{zca} 二进制文件。

\begin{itemize}
  \item 验证：\texttt{zca --version}
  \item 如果缺失，请安装 zca-cli（参见 \texttt{extensions/zalouser/README.md} 或上游 zca-cli 文档）。
\end{itemize}


\subsection{快速设置（新手）}


\begin{enumerate}
  \item 安装插件（见上文）。
  \item 登录（QR，在 Gateway 网关机器上）：
\end{enumerate}

\begin{itemize}
  \item \texttt{openclaw channels login --channel zalouser}
  \item 用 Zalo 手机应用扫描终端中的二维码。
\end{itemize}

\begin{enumerate}
  \item 启用渠道：
\end{enumerate}


\begin{lstlisting}[language=json]
{
  channels: {
    zalouser: {
      enabled: true,
      dmPolicy: "pairing",
    },
  },
}
\end{lstlisting}


\begin{enumerate}
  \item 重启 Gateway 网关（或完成新手引导）。
  \item 私信访问默认为配对模式；首次联系时批准配对码。
\end{enumerate}


\subsection{这是什么}


\begin{itemize}
  \item 使用 \texttt{zca listen} 接收入站消息。
  \item 使用 \texttt{zca msg ...} 发送回复（文本/媒体/链接）。
  \item 专为"个人账户"使用场景设计，适用于 Zalo Bot API 不可用的情况。
\end{itemize}


\subsection{命名}


渠道 ID 为 \texttt{zalouser}，以明确表示这是自动化\textbf{个人 Zalo 用户账户}（非官方）。我们保留 \texttt{zalo} 用于未来可能的官方 Zalo API 集成。

\subsection{查找 ID（目录）}


使用目录 CLI 发现联系人/群组及其 ID：

\begin{lstlisting}[language=bash]
openclaw directory self --channel zalouser
openclaw directory peers list --channel zalouser --query "name"
openclaw directory groups list --channel zalouser --query "work"
\end{lstlisting}


\subsection{限制}


\begin{itemize}
  \item 出站文本分块为约 2000 字符（Zalo 客户端限制）。
  \item 默认阻止流式传输。
\end{itemize}


\subsection{访问控制（私信）}


\texttt{channels.zalouser.dmPolicy} 支持：\texttt{pairing | allowlist | open | disabled}（默认：\texttt{pairing}）。
\texttt{channels.zalouser.allowFrom} 接受用户 ID 或名称。向导会在可用时通过 \texttt{zca friend find} 将名称解析为 ID。

通过以下方式批准：

\begin{itemize}
  \item \texttt{openclaw pairing list zalouser}
  \item \texttt{openclaw pairing approve zalouser }
\end{itemize}


\subsection{群组访问（可选）}


\begin{itemize}
  \item 默认：\texttt{channels.zalouser.groupPolicy = "open"}（允许群组）。使用 \texttt{channels.defaults.groupPolicy} 在未设置时覆盖默认值。
  \item 通过以下方式限制为允许列表：
  \item \texttt{channels.zalouser.groupPolicy = "allowlist"}
  \item \texttt{channels.zalouser.groups}（键为群组 ID 或名称）
  \item 阻止所有群组：\texttt{channels.zalouser.groupPolicy = "disabled"}。
  \item 配置向导可以提示输入群组允许列表。
  \item 启动时，OpenClaw 将允许列表中的群组/用户名称解析为 ID 并记录映射；未解析的条目保持原样。
\end{itemize}


示例：

\begin{lstlisting}[language=json]
{
  channels: {
    zalouser: {
      groupPolicy: "allowlist",
      groups: {
        "123456789": { allow: true },
        "Work Chat": { allow: true },
      },
    },
  },
}
\end{lstlisting}


\subsection{多账户}


账户映射到 zca 配置文件。示例：

\begin{lstlisting}[language=json]
{
  channels: {
    zalouser: {
      enabled: true,
      defaultAccount: "default",
      accounts: {
        work: { enabled: true, profile: "work" },
      },
    },
  },
}
\end{lstlisting}


\subsection{故障排除}


\textbf{找不到 \texttt{zca}：}

\begin{itemize}
  \item 安装 zca-cli 并确保它在 Gateway 网关进程的 \texttt{PATH} 中。
\end{itemize}


\textbf{登录不保持：}

\begin{itemize}
  \item \texttt{openclaw channels status --probe}
  \item 重新登录：\texttt{openclaw channels logout --channel zalouser \&\& openclaw channels login --channel zalouser}
\end{itemize}



\clearpage
\section{IRC}
% === 源文件: channels/irc.md

% === 源文件: channels/irc.md ===
Use IRC when you want OpenClaw in classic channels (\texttt{\#room}) and direct messages.
IRC ships as an extension plugin, but it is configured in the main config under \texttt{channels.irc}.

\subsection{Quick start}


\begin{enumerate}
  \item Enable IRC config in \texttt{\textasciitilde{}/.openclaw/openclaw.json}.
  \item Set at least:
\end{enumerate}


\begin{lstlisting}[language=json]
{
  "channels": {
    "irc": {
      "enabled": true,
      "host": "irc.libera.chat",
      "port": 6697,
      "tls": true,
      "nick": "openclaw-bot",
      "channels": ["#openclaw"]
    }
  }
}
\end{lstlisting}


\begin{enumerate}
  \item Start/restart gateway:
\end{enumerate}


\begin{lstlisting}[language=bash]
openclaw gateway run
\end{lstlisting}


\subsection{Security defaults}


\begin{itemize}
  \item \texttt{channels.irc.dmPolicy} defaults to \texttt{"pairing"}.
  \item \texttt{channels.irc.groupPolicy} defaults to \texttt{"allowlist"}.
  \item With \texttt{groupPolicy="allowlist"}, set \texttt{channels.irc.groups} to define allowed channels.
  \item Use TLS (\texttt{channels.irc.tls=true}) unless you intentionally accept plaintext transport.
\end{itemize}


\subsection{Access control}


There are two separate “gates” for IRC channels:

\begin{enumerate}
  \item \textbf{Channel access} (\texttt{groupPolicy} + \texttt{groups}): whether the bot accepts messages from a channel at all.
  \item \textbf{Sender access} (\texttt{groupAllowFrom} / per-channel \texttt{groups["\#channel"].allowFrom}): who is allowed to trigger the bot inside that channel.
\end{enumerate}


Config keys:

\begin{itemize}
  \item DM allowlist (DM sender access): \texttt{channels.irc.allowFrom}
  \item Group sender allowlist (channel sender access): \texttt{channels.irc.groupAllowFrom}
  \item Per-channel controls (channel + sender + mention rules): \texttt{channels.irc.groups["\#channel"]}
  \item \texttt{channels.irc.groupPolicy="open"} allows unconfigured channels (\textbf{still mention-gated by default})
\end{itemize}


Allowlist entries should use stable sender identities (\texttt{nick!user@host}).
Bare nick matching is mutable and only enabled when \texttt{channels.irc.dangerouslyAllowNameMatching: true}.

\subsubsection{Common gotcha: allowFrom is for DMs, not channels}


If you see logs like:

\begin{itemize}
  \item \texttt{irc: drop group sender alice!ident@host (policy=allowlist)}
\end{itemize}


…it means the sender wasn’t allowed for \textbf{group/channel} messages. Fix it by either:

\begin{itemize}
  \item setting \texttt{channels.irc.groupAllowFrom} (global for all channels), or
  \item setting per-channel sender allowlists: \texttt{channels.irc.groups["\#channel"].allowFrom}
\end{itemize}


Example (allow anyone in \texttt{\#tuirc-dev} to talk to the bot):

\begin{lstlisting}[language=json]
{
  channels: {
    irc: {
      groupPolicy: "allowlist",
      groups: {
        "#tuirc-dev": { allowFrom: ["*"] },
      },
    },
  },
}
\end{lstlisting}


\subsection{Reply triggering (mentions)}


Even if a channel is allowed (via \texttt{groupPolicy} + \texttt{groups}) and the sender is allowed, OpenClaw defaults to \textbf{mention-gating} in group contexts.

That means you may see logs like \texttt{drop channel … (missing-mention)} unless the message includes a mention pattern that matches the bot.

To make the bot reply in an IRC channel \textbf{without needing a mention}, disable mention gating for that channel:

\begin{lstlisting}[language=json]
{
  channels: {
    irc: {
      groupPolicy: "allowlist",
      groups: {
        "#tuirc-dev": {
          requireMention: false,
          allowFrom: ["*"],
        },
      },
    },
  },
}
\end{lstlisting}


Or to allow \textbf{all} IRC channels (no per-channel allowlist) and still reply without mentions:

\begin{lstlisting}[language=json]
{
  channels: {
    irc: {
      groupPolicy: "open",
      groups: {
        "*": { requireMention: false, allowFrom: ["*"] },
      },
    },
  },
}
\end{lstlisting}


\subsection{Security note (recommended for public channels)}


If you allow \texttt{allowFrom: ["*"]} in a public channel, anyone can prompt the bot.
To reduce risk, restrict tools for that channel.

\subsubsection{Same tools for everyone in the channel}


\begin{lstlisting}[language=json]
{
  channels: {
    irc: {
      groups: {
        "#tuirc-dev": {
          allowFrom: ["*"],
          tools: {
            deny: ["group:runtime", "group:fs", "gateway", "nodes", "cron", "browser"],
          },
        },
      },
    },
  },
}
\end{lstlisting}


\subsubsection{Different tools per sender (owner gets more power)}


Use \texttt{toolsBySender} to apply a stricter policy to \texttt{"*"} and a looser one to your nick:

\begin{lstlisting}[language=json]
{
  channels: {
    irc: {
      groups: {
        "#tuirc-dev": {
          allowFrom: ["*"],
          toolsBySender: {
            "*": {
              deny: ["group:runtime", "group:fs", "gateway", "nodes", "cron", "browser"],
            },
            "id:eigen": {
              deny: ["gateway", "nodes", "cron"],
            },
          },
        },
      },
    },
  },
}
\end{lstlisting}


Notes:

\begin{itemize}
  \item \texttt{toolsBySender} keys should use \texttt{id:} for IRC sender identity values:
\end{itemize}

\texttt{id:eigen} or \texttt{id:eigen!\textasciitilde{}eigen@174.127.248.171} for stronger matching.
\begin{itemize}
  \item Legacy unprefixed keys are still accepted and matched as \texttt{id:} only.
  \item The first matching sender policy wins; \texttt{"*"} is the wildcard fallback.
\end{itemize}


For more on group access vs mention-gating (and how they interact), see: \href{/channels/groups}{/channels/groups}.

\subsection{NickServ}


To identify with NickServ after connect:

\begin{lstlisting}[language=json]
{
  "channels": {
    "irc": {
      "nickserv": {
        "enabled": true,
        "service": "NickServ",
        "password": "your-nickserv-password"
      }
    }
  }
}
\end{lstlisting}


Optional one-time registration on connect:

\begin{lstlisting}[language=json]
{
  "channels": {
    "irc": {
      "nickserv": {
        "register": true,
        "registerEmail": "bot@example.com"
      }
    }
  }
}
\end{lstlisting}


Disable \texttt{register} after the nick is registered to avoid repeated REGISTER attempts.

\subsection{Environment variables}


Default account supports:

\begin{itemize}
  \item \texttt{IRC\_HOST}
  \item \texttt{IRC\_PORT}
  \item \texttt{IRC\_TLS}
  \item \texttt{IRC\_NICK}
  \item \texttt{IRC\_USERNAME}
  \item \texttt{IRC\_REALNAME}
  \item \texttt{IRC\_PASSWORD}
  \item \texttt{IRC\_CHANNELS} (comma-separated)
  \item \texttt{IRC\_NICKSERV\_PASSWORD}
  \item \texttt{IRC\_NICKSERV\_REGISTER\_EMAIL}
\end{itemize}


\subsection{Troubleshooting}


\begin{itemize}
  \item If the bot connects but never replies in channels, verify \texttt{channels.irc.groups} \textbf{and} whether mention-gating is dropping messages (\texttt{missing-mention}). If you want it to reply without pings, set \texttt{requireMention:false} for the channel.
  \item If login fails, verify nick availability and server password.
  \item If TLS fails on a custom network, verify host/port and certificate setup.
\end{itemize}



\clearpage

% === 源文件: channels/synology-chat.md ===
\section{Synology Chat (plugin)}


Status: supported via plugin as a direct-message channel using Synology Chat webhooks.
The plugin accepts inbound messages from Synology Chat outgoing webhooks and sends replies
through a Synology Chat incoming webhook.

\subsection{Plugin required}


Synology Chat is plugin-based and not part of the default core channel install.

Install from a local checkout:

\begin{lstlisting}[language=bash]
openclaw plugins install ./extensions/synology-chat
\end{lstlisting}


Details: \href{/tools/plugin}{Plugins}

\subsection{Quick setup}


\begin{enumerate}
  \item Install and enable the Synology Chat plugin.
  \item In Synology Chat integrations:
\end{enumerate}

\begin{itemize}
  \item Create an incoming webhook and copy its URL.
  \item Create an outgoing webhook with your secret token.
\end{itemize}

\begin{enumerate}
  \item Point the outgoing webhook URL to your OpenClaw gateway:
\end{enumerate}

\begin{itemize}
  \item \texttt{https://gateway-host/webhook/synology} by default.
  \item Or your custom \texttt{channels.synology-chat.webhookPath}.
\end{itemize}

\begin{enumerate}
  \item Configure \texttt{channels.synology-chat} in OpenClaw.
  \item Restart gateway and send a DM to the Synology Chat bot.
\end{enumerate}


Minimal config:

\begin{lstlisting}[language=json]
{
  channels: {
    "synology-chat": {
      enabled: true,
      token: "synology-outgoing-token",
      incomingUrl: "https://nas.example.com/webapi/entry.cgi?api=SYNO.Chat.External&method=incoming&version=2&token=...",
      webhookPath: "/webhook/synology",
      dmPolicy: "allowlist",
      allowedUserIds: ["123456"],
      rateLimitPerMinute: 30,
      allowInsecureSsl: false,
    },
  },
}
\end{lstlisting}


\subsection{Environment variables}


For the default account, you can use env vars:

\begin{itemize}
  \item \texttt{SYNOLOGY\_CHAT\_TOKEN}
  \item \texttt{SYNOLOGY\_CHAT\_INCOMING\_URL}
  \item \texttt{SYNOLOGY\_NAS\_HOST}
  \item \texttt{SYNOLOGY\_ALLOWED\_USER\_IDS} (comma-separated)
  \item \texttt{SYNOLOGY\_RATE\_LIMIT}
  \item \texttt{OPENCLAW\_BOT\_NAME}
\end{itemize}


Config values override env vars.

\subsection{DM policy and access control}


\begin{itemize}
  \item \texttt{dmPolicy: "allowlist"} is the recommended default.
  \item \texttt{allowedUserIds} accepts a list (or comma-separated string) of Synology user IDs.
  \item In \texttt{allowlist} mode, an empty \texttt{allowedUserIds} list is treated as misconfiguration and the webhook route will not start (use \texttt{dmPolicy: "open"} for allow-all).
  \item \texttt{dmPolicy: "open"} allows any sender.
  \item \texttt{dmPolicy: "disabled"} blocks DMs.
  \item Pairing approvals work with:
  \item \texttt{openclaw pairing list synology-chat}
  \item \texttt{openclaw pairing approve synology-chat }
\end{itemize}


\subsection{Outbound delivery}


Use numeric Synology Chat user IDs as targets.

Examples:

\begin{lstlisting}[language=bash]
openclaw message send --channel synology-chat --target 123456 --text "Hello from OpenClaw"
openclaw message send --channel synology-chat --target synology-chat:123456 --text "Hello again"
\end{lstlisting}


Media sends are supported by URL-based file delivery.

\subsection{Multi-account}


Multiple Synology Chat accounts are supported under \texttt{channels.synology-chat.accounts}.
Each account can override token, incoming URL, webhook path, DM policy, and limits.

\begin{lstlisting}[language=json]
{
  channels: {
    "synology-chat": {
      enabled: true,
      accounts: {
        default: {
          token: "token-a",
          incomingUrl: "https://nas-a.example.com/...token=...",
        },
        alerts: {
          token: "token-b",
          incomingUrl: "https://nas-b.example.com/...token=...",
          webhookPath: "/webhook/synology-alerts",
          dmPolicy: "allowlist",
          allowedUserIds: ["987654"],
        },
      },
    },
  },
}
\end{lstlisting}


\subsection{Security notes}


\begin{itemize}
  \item Keep \texttt{token} secret and rotate it if leaked.
  \item Keep \texttt{allowInsecureSsl: false} unless you explicitly trust a self-signed local NAS cert.
  \item Inbound webhook requests are token-verified and rate-limited per sender.
  \item Prefer \texttt{dmPolicy: "allowlist"} for production.
\end{itemize}



\clearpage

